% Compile with pdflatex paper.tex until the cross-references stabilize.
\documentclass[11pt,a4paper]{article}
\usepackage[T1]{fontenc}
\usepackage{lmodern}
\usepackage{amsmath,amssymb,amsthm}
\usepackage{geometry,xcolor,hyperref,fancyhdr}
\geometry{left=28mm,right=28mm,top=26mm,bottom=25mm,headheight=15pt}
\definecolor{accent}{HTML}{214A63}
\hypersetup{unicode=true,colorlinks=true,linkcolor=accent,citecolor=accent,urlcolor=accent,
  pdftitle={The Small-Volume Spectral Gap of Four-Dimensional SU(2) Wilson Gauge Theory},
  pdfsubject={The fixed-volume ultraviolet limit and the weak-coupling spectrum},
  pdfkeywords={Yang-Mills, Wilson transfer operator, finite volume, spectral gap}}
\urlstyle{same}
\newtheorem{theorem}{Theorem}
\newtheorem{proposition}[theorem]{Proposition}
\newtheorem{lemma}[theorem]{Lemma}
\newtheorem{corollary}[theorem]{Corollary}
\numberwithin{equation}{section}
\newcommand{\SU}{\mathrm{SU}}
\newcommand{\DTV}{d_{\mathrm{TV}}}
\newcommand{\HS}{\mathrm{HS}}
\newcommand{\Tr}{\operatorname{Tr}}
\newcommand{\op}{\mathrm{op}}
\newcommand{\hper}{\mathfrak h}
\setlength{\parskip}{.35\baselineskip}
\setlength{\parindent}{1.35em}
\setlength{\emergencystretch}{2em}
\linespread{1.08}
\allowdisplaybreaks[2]
\raggedbottom
\pagestyle{fancy}
\fancyhf{}
\fancyhead[L]{\footnotesize\color{accent}Small-volume Wilson spectrum}
\fancyhead[R]{\footnotesize\color{accent}Finite-volume spectral theory}
\fancyfoot[C]{\thepage}
\renewcommand{\headrulewidth}{0.25pt}
\setcounter{tocdepth}{1}
\makeatletter
\newcommand{\articlecontents}{%
  \begingroup\small
  \setlength{\parskip}{0pt}%
  \renewcommand{\l@section}{\@dottedtocline{1}{0pt}{1.8em}}%
  \tableofcontents
  \endgroup
}
\makeatother
\title{The Small-Volume Spectral Gap of Four-Dimensional \texorpdfstring{\(\SU(2)\)}{SU(2)} Wilson Gauge Theory}
\author{}
\date{November 23, 2026}
\begin{document}
\maketitle
\begin{abstract}
Consider the four-dimensional \(\SU(2)\) Wilson transfer operator on a fixed spatial three-torus of side \(L_0\), with no magnetic twist, in the sector of zero electric flux and exact zero total momentum. Along a selected sequence of Wilson bare couplings, the first energy gap has a full ultraviolet limit at each fixed terminal inverse coupling \(v\). Set \(g=(4/v)^{1/2}\), and let \(b\) be the reference inverse coupling, asymptotic to \(v\). In the subsequent limit \(b\to\infty\), the gap is
\[
 \frac{g^{2/3}}{2L_0}(\nu_1-\nu_0)\bigl[1+O(b^{-1/8})\bigr],
\]
where \(\nu_0<\nu_1\) are the first two distinct eigenvalues of the nine-dimensional quartic Schr\"odinger operator in the color-singlet subspace. A short-period heat-operator comparison gives the full cutoff limit. A longer period resolves the low-energy cluster and preserves the center and spatial symmetries. The continuous-time vacuum representation generated by bounded Wilson observables at fixed physical resolution has the same first gap and multiplicity. The cyclic subspace generated by the specified separated renormalized local fields inherits the lower bound. Equality there requires positive spectral weight in the first excited cluster.
\end{abstract}

\noindent\textbf{Keywords.} Wilson transfer operator, finite volume, spectral gap, renormalization group, quartic matrix quantum mechanics.

\articlecontents

\section{The finite-volume problem and its physical origin}

The Wilson lattice regularization uses group elements on links and preserves exact gauge invariance at every lattice spacing. The transfer operator connects Euclidean correlations to the energy spectrum on the Hilbert space satisfying Gauss's law. L\"uscher proved strict positivity of the Wilson transfer operator. Osterwalder and Seiler developed the corresponding framework using lattice reflection positivity. A spectral statement after removal of the lattice cutoff also requires comparison between transfer operators at different cutoffs \cite{wilson,luscher-transfer,osterwalder-seiler}.

The predicted weak-coupling behavior dates to work from 1983 to 1987. In 1983, Martin L\"uscher developed the \(g^{2/3}\) expansion on a periodic small box and carried it through one loop. In 1986 and 1987, Jeffrey Koller and Pierre van Baal studied the effective Hamiltonian of the spatially constant modes, electric-flux energies and their tunneling corrections, and the low-lying spectrum using analytic and variational methods. Van Baal's 2000 review, \emph{QCD in a Finite Volume}, gives a systematic account of the Hamiltonian approach, gauge fixing, and finite-volume spectra \cite{luscher-torus,koller-vanbaal,vanbaal-koller-torus,vanbaal-review}.

The exponent is determined by the spatial zero modes. On a three-torus of side \(L_0\), each of the three spatial directions has a constant mode in the three-dimensional Lie algebra \(\mathfrak{su}(2)\). Write \(c_i\), \(1\le i\le3\), for its dimensionless color coordinates. For constant modes the linear magnetic curvature vanishes and the leading term is the Lie bracket. The magnetic potential is quartic and vanishes when the three color vectors are parallel. The zero-point energy of the transverse oscillators confines motion along these flat directions. If \(r\) is the typical amplitude of the constant modes, the kinetic and magnetic energies scale as \(g^2r^{-2}/L_0\) and \(g^{-2}r^4/L_0\), respectively. Balancing the two gives \(r\asymp g^{2/3}\) and the energy scale \(g^{2/3}/L_0\).

The Millennium Prize problem formulated by Jaffe and Witten asks for a nontrivial quantum Yang--Mills theory on \(\mathbb R^4\) for any compact simple gauge group, together with a strictly positive mass gap. The present result concerns the first excitation on a fixed periodic spatial three-torus \cite{jaffe-witten}.

OpenAI's Lean construction supplies selected sequences of Wilson bare couplings and a fixed-volume ultraviolet limit for local fields on equal-period Euclidean tori. We use its exact blocking construction and scale estimates as inputs. To our knowledge, this is the first derivation of the gap asymptotic with a quantitative error from the original Wilson transfer spectrum along the periodic trajectory selected by that construction. We identify the same first gap and multiplicity in the continuous-time vacuum representation of bounded Wilson observables. The physical side length and terminal inverse coupling are fixed during removal of the cutoff, and the weak-coupling limit is taken afterward. The cyclic subspace of specified renormalized local fields satisfies the same spectral lower bound \cite{formal,armstrong}.

A fixed short period allows comparison of ordered heat eigenvalues at different ultraviolet cutoffs in the same physical box. A long period, comparable to the inverse of the predicted gap, separates the low-lying eigenvalues of the heat operator. The periods use the same bare coupling, and the comparison retains the joint distribution of the actual coarse-grained endpoint observables. This finite-cutoff comparison relates the classical effective-Hamiltonian calculation to the full-sequence limit of the Wilson transfer spectrum.

\section{Theorem, symmetry sectors, and order of limits}
\label{sec:main}

Let \(\mathbb T^3_{L_0}=(\mathbb R/L_0\mathbb Z)^3\) and \(G=\SU(2)\). At the coarsest level there are \(P\) sites in each spatial direction, with \(P/8\) an admissible period divisible by four. The level \(f\) increases toward finer scales. The odd integer \(\ell_f\) is the ratio for blocking level \(f+1\) to level \(f\). Set
\[
 D_0=1,\qquad D_f=\prod_{j<f}\ell_j,\qquad n_f=PD_f,\qquad a_f=L_0/n_f.
\]
Thus \(n_f\) is the number of sites per spatial direction and \(a_f\) the lattice spacing. The ultraviolet cutoff is at level \(N\). For an oriented plaquette \(p\), let \(U_p\) be the ordered product of links around its boundary. The Wilson weight is \(\exp[-\beta\sum_p(1-\tfrac12\operatorname{Re}\Tr U_p)]\), with the matrix trace in the defining representation.

On the Euclidean four-torus of side \(L_0\), tune the bare inverse coupling so that its blocking trajectory has inverse coupling \(v\) at level zero. Write \(\beta_{N,P,f}(v)\) for the inverse coupling at level \(f\), and \(\beta_{N,P,N}(v)\) for the bare value. Use that same bare value for every time period at the same spatial size \(L_0\). In Wilson's convention the coupling associated with inverse coupling \(\beta\) is \((4/\beta)^{1/2}\). Throughout, \(g=(4/v)^{1/2}\) is the coupling defined by the terminal inverse coupling, and \(b\) is the reference inverse coupling about which the terminal band is chosen.

Let \(K_N\) be the positive transfer operator for one Euclidean time step \(a_N\). On the product Haar space of a spatial slice, \(M\) multiplies by the square root of the spatial plaquette weight. The tensor product \(C_\beta\) of central convolutions has single-link convolution density proportional to \(\exp(\beta\chi_{1/2}/2)\), where \(\chi_{1/2}\) is the defining character and \(\beta=\beta_{N,P,N}(v)\). Integration over time links implements the Haar average \(\mathcal G\) over slice gauge transformations. The physical Hilbert space is \(\mathcal H_N=\operatorname{ran}\mathcal G\).

Normalize the convolution density to be a probability density. The operator \(K_N=\mathcal G M C_\beta M\mathcal G\) on \(\mathcal H_N\) then differs from the Wilson one-step transfer operator by a field-independent factor. Both the spatial weight and the central convolution commute with \(\mathcal G\). The tensor powers \(\chi_{1/2}^r\) have nonnegative character coefficients. For spin \(j\), the coefficient from \(r=2j\) is positive. Peter--Weyl theory therefore gives strict positivity of \(C_\beta\), and \(M>0\) transfers it to \(K_N\). Its kernel is smooth and strictly positive on finitely many compact group factors. Consequently \(K_N\) is compact and has a simple top eigenvalue with strictly positive normalized ground state \(\Omega_N\). The Wilson weight preserves center transformations, spatial translations, and inversion. Define the vacuum-shifted Hamiltonian, heat semigroup, and heat trace by
\[
 H_N=-a_N^{-1}\log(K_N/\|K_N\|),\qquad
 S_N(t)=e^{-tH_N},\qquad Z_N^0(t)=\Tr S_N(t).
\]
Unless otherwise indicated, heat traces are on the full Gauss-law space. For integer \(t/a_N\), let \(Z_N(t)=\Tr K_N^{t/a_N}\) be the trace before vacuum shifting. It differs from the Wilson partition function by a power of the one-step normalization. That factor cancels in the heat-trace ratios at fixed bare coupling and fixed spatial section used here.

Use periodic boundary conditions with no magnetic twist. In each spatial direction the center transformation multiplies the links crossing a dual two-torus by \(-I\). These transformations preserve plaquette weights and generate an action of \((\mathbb Z/2)^3\) on \(\mathcal H_N\). Its trivial-character subspace \(\mathcal H_N^0\) is the zero-electric-flux sector. Its intersection with the invariants of the three one-site translations is the exact zero-total-momentum sector \(\mathcal H_N^{00}\). Denote the first positive energy of \(H_N\) on \(\mathcal H_N^{00}\) by \(m_N^{00}(L_0,v)\).

Positive constants \(c,C\) can change between estimates. Weak-coupling bounds hold uniformly in the specified terminal band on a common cutoff tail. The notation \(C_g=g^{-o(1)}\) denotes at most subpower growth as \(g\downarrow0\), namely an eventual \(O_\epsilon(g^{-\epsilon})\) bound for every \(\epsilon>0\). Operator and Hilbert--Schmidt norms are denoted by \(\|\cdot\|_{\op}\) and \(\|\cdot\|_{\HS}\).

\begin{theorem}[Fixed-volume Wilson transfer spectrum]
\label{thm:main}
Fix \(L_0>0\) and \(0<c_0<1/64\). Fix the stage parameters and blocking ratio for the weak-coupling stopping scale first. For every sufficiently large reference inverse coupling \(b\), one can choose an admissible blocking schedule and a base period \(P=P(b)\) that work throughout the real terminal band
\[
                   |v-b|\le c_0 b^{31/32}.
\]
For each \(v\) in this band, \(m_N^{00}(L_0,v)\) has a full-sequence limit \(m^{00}(L_0,v)\) as \(N\to\infty\). Let
\[
 H_4=-\Delta_X+\sum_{i<j}|X_i\times X_j|^2,
 \qquad X=(X_1,X_2,X_3)\in(\mathbb R^3)^3
\]
be the Friedrichs realization, with subscript \(4\) denoting the degree of the potential. The color-singlet subspace consists of functions invariant under simultaneous rotations \(X_i\mapsto RX_i\), \(R\in\mathrm{SO}(3)\). Let \(\nu_0<\nu_1\) be the first two distinct eigenvalues of \(H_4\) restricted to that subspace. Then
\begin{equation}
 \label{eq:main}
 L_0 m^{00}(L_0,v)
  =\frac{\nu_1-\nu_0}{2}(4/v)^{1/3}
           \bigl[1+O(b^{-1/8})\bigr].
\end{equation}
The error constant is uniform in the stated real terminal band. The weak-coupling estimate holds on a common tail \(N\ge N_*(b)\) of the original cutoff sequence. First take \(N\to\infty\) at fixed \(b,v\), then take \(b\to\infty\).
\end{theorem}

\begin{proposition}[Vacuum representation on the fixed torus]
\label{prop:representation}
Along the same bare trajectory, the continuous-time vacuum representation generated by bounded, center-invariant Wilson observables at fixed nonzero physical resolution has the first gap and multiplicity of the Wilson limit. For the specified renormalized local insertions satisfying the separated reflection comparison on the short square torus, the cyclic space generated by products with positive time separations reduces the same time semigroup. Its Hamiltonian on the orthogonal complement of the vacuum is bounded below by \(m^{00}(L_0,v)\).
\end{proposition}

The square four-torus of short period \(L_0\) gives full-sequence convergence at a fixed terminal inverse coupling. To identify the first excited cluster, take a long period \(T\) for which the dimensionless time \(Tg^{2/3}/L_0\) stays in a compact subinterval of \((0,\infty)\). At this time scale, separated levels of order \(g^{2/3}/L_0\) correspond to separated heat eigenvalues.

\section{Quantitative comparison of Wilson heat operators}
\label{sec:kernel-comparison}

Let \(\kappa_G>0\) be the one-loop drift coefficient for the inverse coupling in the public blocking construction. Use the reference inverse coupling \(\widehat\beta_f=b+\kappa_G\log D_f\) at level \(f\). The small-field threshold and the covering radius in the support norm are
\[
 \lambda_f=\widehat\beta_f,\qquad M_f=\exp\bigl((\log\log\lambda_f)^2\bigr),\qquad
 t_f=M_f/\sqrt{\lambda_f},\qquad B_f^{\mathrm{cov}}\asymp(1+\log\lambda_f)^2.
\]
These reference quantities are fixed throughout the real terminal band. Each stage uses a fixed blocking ratio and fixed constants for the local estimates. After choosing these data, increase \(b\). Within the chosen stage take a stopping level \(m\) and a finer auxiliary level \(m_b>m\) with
\begin{equation}
 n=PD_m\asymp b^{1/24},\qquad PD_{m_b}\asymp b^{1/12}.
 \label{eq:paper-stops}
\end{equation}
The finite prefix before the chosen stage can have total logarithmic scale \(o(\log b)\). The reference periods \(P/8,P,2P\) satisfy the locality requirements of the blocking maps. A fixed power of \(\log b\) suffices for the base period of the finite stage. Later stages begin once their dependence ranges are sufficiently small relative to the actual period. On the selected cutoff tail, the accumulated drift satisfies
\begin{equation}
 \big|\beta_{N,P,f}(v)-v-\kappa_G\log D_f\big|
 \le \tfrac14\kappa_G\log D_f,\qquad
 \beta_{N,P,m}(v)=v+O(\log b).
 \label{eq:paper-coupling}
\end{equation}
These bounds are uniform in the prescribed real terminal band. Both cylinders use the bare sequence fixed on the square four-torus. Write \(\beta_m=\beta_{N,P,m}(v)\) and \(g_m=(4/\beta_m)^{1/2}\) for the actual stopping-scale inverse coupling and coupling.

Fix a compact interval \(\mathcal I\subset(0,\infty)\). For target time \(u_{\mathrm{tar}}\in\mathcal I\), take a physical period \(T=kL_0\) and dimensionless time \(u=Tg^{2/3}/L_0\) satisfying
\begin{equation}
 k\in\mathbb N,\qquad u=kg^{2/3}=u_{\mathrm{tar}}+O(g^{2/3}),\qquad
 q=kn\asymp g^{-3/4}.
 \label{eq:paper-longtime}
\end{equation}
Here \(q\) is the number of time steps at level \(m\), and \(T/2\) is an integer number of fine time steps. A \emph{coarse-grained observable} is a collection of retained link variables at a fixed level, or a measurable function of those links. The blocking kernels include random branch choices and conditional Haar noise. A \emph{large-field flag} records a local curvature or matching test that exceeds its threshold. Local integrals retain the corresponding event and support. All comparisons are based on the joint law of one Wilson cylinder. For densities \(p_1,p_2\) with respect to a common reference measure, use
\[
 d_{\mathrm H}(p_1,p_2)=\|\sqrt{p_1}-\sqrt{p_2}\|_{L^2},\qquad
 \DTV(p_1,p_2)=\tfrac12\|p_1-p_2\|_{L^1}
\]
for Hellinger distance and total variation distance.

\subsection{Short heat traces and the long-cylinder density}

Pressure comparison between the short reference tori gives
\begin{equation}
 \Tr S_N(s_*)\le \frac{Z_N(s_*)^2}{Z_N(2s_*)}\le C_{b,v},
 \qquad s_*=\frac{L_0}{4}.
 \label{eq:paper-short-trace}
\end{equation}
The first inequality follows from
\(\Tr K_N^{2s_*/a_N}\le \|K_N\|^{s_*/a_N}\Tr K_N^{s_*/a_N}\).
The second compares partition functions of different shapes at the same bare coupling. The pressure difference per lattice cell is \(O(a_N^4)\), leaving a cutoff-uniform bound after multiplication by the number of cells at fixed physical volume. For the weak-coupling comparison we also need
\begin{equation}
  \log Z_N^0(L_0)\le C_\varepsilon b^\varepsilon\qquad(\varepsilon>0).
 \label{eq:paper-short-entropy}
\end{equation}
For fixed \(\varepsilon\), choose a coarse-graining level \(f_\alpha\) on the actual short square torus with \(n_\alpha=n_{f_\alpha}\asymp b^\alpha\), where \(0<\delta_\alpha<\alpha\le1/12\). Retain a band of width \(L_0b^{-\alpha+\delta_\alpha}\) at each end, containing \(O(b^{3\alpha+\delta_\alpha})\) links. Given the parent field and blocking labels, the link noises at the last step are independent. Their Haar densities are bounded by \(C(1+\kappa)^{3/2}\), where \(\kappa\) is the penalty coefficient for the conditional noise. The normalization integral on a ball of radius \(O(\kappa^{-1/2})\) gives this bound. The endpoint mutual information is thus \(O(b^{3\alpha+\delta_\alpha}\log(2+b))\).

Write \(e_f\) for the mean total variation error in the true posterior comparison between adjacent levels, and \(B_f^{\mathrm{lik}}\) for the corresponding log-likelihood bound. On the short torus, \(\sup_N\sum_{f>f_\alpha}B_f^{\mathrm{lik}}e_f=O_A(b^{-A})\). The Hellinger affinity of the opposite-cut joint law and the product of its marginals, the inner product of their square-root densities, is at most \(Z_N^0(L_0)^{-1/2}\). Relative entropy is at least minus twice the logarithm of that affinity. Taking \(\alpha,\delta_\alpha\) sufficiently small gives \eqref{eq:paper-short-entropy}. This short-torus bound is obtained before the long-cylinder estimates in Appendix~\ref{cam-bound-section}.

The comparison of blocking on a long cylinder and on a square torus retains the integer paths, dependence domains, activity tests, and conditional noise at each anchor. The activities are local renormalization-group coefficients indexed by finite supports. The polymer expansion sums compatible supports and retains the original trigger events for large-field activities. For a support \(Y\), let \(\operatorname{cov}(Y)\) be the minimum number of balls of radius \(B_f^{\mathrm{cov}}\) that cover \(Y\) with connected union. Let \(o\) denote a lattice anchor and \(\sigma\) the fixed exponential-weight parameter. The difference \(\delta A_Y\) of regular activities is measured in the norm
\[
 \sup_o\sum_{Y\ni o}\exp((\sigma-2)\operatorname{cov}(Y))
 \sup_{\text{common configuration domain}}|\delta A_Y|.
\]
In a real local direction, use a centered difference of step \(\eta\). If \(U\) bounds the value difference and \(V\) bounds its fourth derivative, the difference in second derivatives is bounded by \(4U/\eta^2+V\eta^2/12\). Summing with the common weight before choosing the step gives one square-root loss. We truncate both the dependence radius and the weighted path length, so that long paths with many returns in a small region also have an exponential tail. Differences of local inverses use \(R_1^{-1}-R_2^{-1}=R_1^{-1}(R_2-R_1)R_2^{-1}\).

Fix \(\vartheta=1/2\), then choose odd blocking ratios with \(\vartheta\ell_f^{1/4}>4\). Let \(j\) denote the stage of level \(f\), and \(d_f\) the anchor-uniform weighted discrepancy of regular activities between the two geometries. Then
\begin{equation}
 d_f\le C_j(1+n_f+\widehat\beta_f)^{C_j}
   \left(d_{f+1}
      +\exp\!\left[-\frac{c_j n_f^{1/3}}
       {1+B_f^{\mathrm{cov}}+B_{f+1}^{\mathrm{cov}}}\right]
   \right)^\vartheta,\qquad d_N=0 .
 \label{eq:paper-shape-recurrence}
\end{equation}
The initial discrepancy at level \(N\) vanishes because the two local bare Wilson actions agree. Since \(\log b=O(\log n_f)\), the choice of stage entrances and blocking ratios gives \(d_f\le e^{-c n_f^{1/4}}\) by induction toward coarser levels. We extract the six affine curvature directions separately and retain them in the relative density of the original positive law. Summing over \(O(kn_f^4)\) anchors gives cumulative log-density oscillation \(O(ke^{-c'n^{1/4}})\). This uses the affine calibration at all frequencies and the likelihood ratio of the actual blocking noise. See Appendix~\ref{app:geometry}.

On the short square torus, the centered total regular action is bounded by \(n^4\epsilon\), where \(\epsilon=O(\widehat\beta_m^{-11/10}+e^{-c n^{1/4}})\). To transfer this bound, fold a segment of the long cylinder into the square torus of period \(n\) in the truncated local formula. The four stages are a pause, forward traversal, a second pause, and backward traversal. During a pause the spatial links are held constant and the time links are the identity. During backward traversal the time links are inverted. Take two traversal lengths \(r_1=\lfloor n/8\rfloor\), \(r_2=\lfloor n/4\rfloor\), both exceeding twice the local dependence radius. Averaging over the starting time makes the seam contribution independent of the traversal length. If \(A_r\) is the averaged folded action and \(\bar f^*\) the static contribution per time step, the long-cylinder average per time step is
\[
 \bar f=\bar f^*+\frac{A_{r_2}-A_{r_1}}{2(r_2-r_1)}.
\]
The short-torus bound gives \(17kn^4\epsilon\). Including the errors from truncating the local actions, the centered log-density oscillation of the long-cylinder law on the small-curvature region relative to the common analytic reference is
\begin{equation}
 O(kn^4 b^{-11/10}+kn^4e^{-c n^{1/4}})
      =O(b^{-3/5+o(1)}).
 \label{eq:paper-density}
\end{equation}
If \(F\) is the log-density difference, the separately normalized ratio on that region is \(e^F/\mathbb E e^F\) and gives the same order for total variation. The large-field region has mass \(O_A(b^{-A})\), so the full Hellinger error is \(O(b^{-3/10+o(1)})\). See Appendices~\ref{app:geometry} and~\ref{cam-bound-section}.

\subsection{A common operator space for fine slices and coarse observations}

Fix \(0<\delta\le10^{-3}\). Retain the fine-cut configurations \(x,y\) and the full coarse observables \(u_1,u_2\) in two time bands of half-width \(h\asymp L_0b^{-1/24+\delta}\) about opposite cuts. Write \(\mu\) for their joint Wilson-cylinder law, \(\mu_{u_1,u_2}\) for the marginal of the coarse observations, and use analogous notation for other marginals and conditionals.

The comparison from coarse to fine levels uses the short-torus positive curvature penalty and reflection-positive chessboard estimates for large-field events. Let \(P_f^{\mathrm{thr}}\) be the fixed logarithmic power dividing \(t_f\) to produce the tightened threshold. Let \(S_f^{\mathrm{cell}}\) be the product of adjacent blocking ratios needed to contain the full dependence domain in a reflection cell. The failure probability for a specified test is controlled by an exponential with negative exponent of order
\[
 \frac{cM_f^{1/5}}{(S_f^{\mathrm{cell}})^4(P_f^{\mathrm{thr}})^2}
      - C\log\lambda_f-\frac{C\log Z_N^0(L_0)}{n_f^4}.
\]
The last term is controlled by \eqref{eq:paper-short-entropy} and \(n_f^4\ge n^4\). For conditional noise the single-link tail is \(C(1+\kappa)^{3/2}\exp[-c\kappa(t_f/P_f^{\mathrm{thr}})^2]\). Given the parent field and original labels, the noises are independent. At level \(f\) take buffer width \(h_f^{\mathrm{buf}}=ch\sqrt{n/n_f}\), with sum below \(h/4\). Fix a sufficiently large exponent \(K\). For \(n_f\le\lambda_f^K\), the cylinder has at most a fixed power of \(b\) sites. A union bound then controls failures and a common small-field domain gives the weighted Hessian lower bound. At finer levels a finite-range positive probability approximation is used. A failure then requires a connected chain of failing cells of length at least \(n_f^{1/4}\). Both the failure and conditional mixing errors can be bounded by any prescribed negative power of \(b\). Summing over levels gives
\begin{equation}
 d_{\mathrm{TV}}\big(\mu,\,
    \mu_{u_1,u_2}\,\mu_{x\mid u_1}\mu_{y\mid u_2}\big)
         =O_A(b^{-A})\qquad(A>0).
 \label{eq:paper-backward}
\end{equation}
If local mixing first produces candidate one-end kernels, conditioning the actual one-end marginal on its coarse observable replaces each candidate by the true conditional kernel at an additional cost of at most the original error. See Appendix~\ref{cam-bound-section}.

To generate the coarse observations from fine cuts, take ground-state-transformed Markov paths on each side of each cut. Making the ends conditionally independent given the fine cuts has squared Hellinger error at most \(2[Z_N^0(T-4h)/Z_N^0(T)-1]\). The logarithm of the heat trace is decreasing and convex, so
\begin{equation}
 0\le \log\frac{Z_N^0(T-4h)}{Z_N^0(T)}
   \le \frac{4h}{T-4h-L_0}\log Z_N^0(L_0).
 \label{eq:paper-forward}
\end{equation}
By \eqref{eq:paper-short-entropy}, taking the square root gives \(O(b^{-3/16+\delta/2+o(1)})\). Let \(r_x\) be the conditional density of the one-end observation generated from the ground-state path. The map \(f(x)\mapsto\sqrt{r_x(u_1)}f(x)\) is an isometric embedding. The reverse embedding uses the one-end conditional density of the actual finite-period law. The square roots of the two comparison densities give the zero extensions of two operators on the common space. The first is the normalized half-period heat operator
\[
 A_N(T)=\frac{S_N(T/2)}{\sqrt{Z_N^0(T)}}.
\]
The second, \(C_{N,m}\), has integral kernel given by the square root of the joint density of the two coarse observables. Spatial inversion, center transformations, and translations by one stopping-scale lattice step preserve the conditional laws. Both embeddings are therefore equivariant. The Hilbert--Schmidt error is controlled by the square root of \eqref{eq:paper-forward}. See Appendix~\ref{cam-bound-section}.

\subsection{Integration outside the center neighborhoods}

At the stopping scale, let \(P_i\) be the spatial holonomy based at a common site and winding once along direction \(i\). Giving each of the three holonomies a value in \(\{\pm I\}\) produces eight center configurations. A center transformation takes the selected center to the identity. Use \(C_i=\log P_i\) in the resulting holonomy chamber. Use the Pauli basis \(t_a=-i\sigma_a/2\) of \(\mathfrak{su}(2)\), with \([t_a,t_b]=\sum_c\epsilon_{abc}t_c\). Identify the Lie algebra with \(\mathbb R^3\) in this basis. On each spatial slice, ordinary Coulomb gauge gives
\[
 U_i(x)=\exp(c_i/n+w_i(x)),\qquad \sum_xw_i(x)=0,\qquad d_0^*w=0 .
\]
Here \(c_i\) is \(n\) times the mean link logarithm, \(w_i(x)\) is the zero-mean remainder, and \(d_0^*\) is the adjoint of the flat lattice differential. Let \(\tau=g^{1-o(1)}\) be the curvature bound used for this chart. Based paths and plaquette fillings give \(\|[P_i,P_j]-I\|\le Cn^2\tau\), with \([P_i,P_j]\) the group commutator. Distributing fractional powers of the holonomies along the three axes gives a periodic representative in the chamber. Let \(D_C\) be the derivative of the gauge action at the constant configuration \(\exp(C_i/n)\), and set \(M_C=d_0^*D_C\). With \(\omega_k\asymp\min(1,|k|/n)\) for nonzero spatial Fourier index \(k\), one has \(\|M_C(k)^{-1}\|\le C\omega_k^{-2}\) and
\[
 \|M_C^{-1}d_0^*\|_{\infty\to\infty}\le Cn,qquad
 \|D_C M_C^{-1}d_0^*\|_{\infty\to\infty}\le C(1+\log n).
\]
The contraction condition is \(Cn^4\tau(1+\log n)=o(1)\). Propagating based gauge transformations link by link in the same chamber gives uniqueness of the small-field representative and
\[
 \|w\|_\infty\le Cn^2\tau(1+\log n),\qquad
 |c-C|\le Cn^3\tau(1+\log n).
\]

On each nonzero real Fourier block, the volume factor for the gauge-orbit and transverse coordinates is \(|\det B_D|\sqrt{\det\Delta_D}\), where \(D\) is the homogeneous reference connection. The transverse map is \(B_D=P_D\mathcal J_D|_{\ker d_0^*}\), with \(\mathcal J_D\) the differential from logarithmic to link-tangent coordinates and \(P_D\) the orthogonal projection onto the covariantly divergence-free space. If \(d_D\) is the covariant lattice differential at \(D\), then \(\Delta_D=d_D^*d_D\). The transverse and Gauss changes of variables cancel the first factor, and the Wilson reference Gauss integral cancels the second. The remaining determinant is transverse and relative to the Wilson reference action. The spatial zero-mode Haar factor changes in logarithm by \(O(q|D|^2/n^2)\). The same positive flat-reference normalization \(C_{n,q,g}\) applies inside and outside the center neighborhood.

The comparison retains the actual temporal holonomy and four charged spatial zero-mode oscillators, from the two transverse spatial directions and two transverse color directions at a rank-one background. Set \(\rho=g^{5/9}\) as the center-ball radius. Write \(\mathcal R_1\) for the cone of matrices of rank at most one, with the three color vectors as columns. Freezing the nonzero-mode determinant at a nearest \(D\in\mathcal R_1\) has second-order error at most \(C_gn^3\operatorname{dist}(c,\mathcal R_1)^2\), absorbed by the quartic magnetic potential. The four charged zero-point energies and the frozen nonzero-mode potential in the same normalization together contribute at least \(c|D|-o(\rho)\). Near the rank-one cone use a single anchor for the full path. After fixing five tangential paths, moving the center in time cannot increase the four transverse Gaussian integrals. The resulting positive mass comparison is
\begin{equation}
 \frac{\text{unnormalized mass outside the center neighborhoods}}
      {\text{unnormalized mass of a center ball in the same normalization}}
       \le C\exp[-c g^{-1/9+o(1)}],
 \label{eq:paper-outside}
\end{equation}
where \(\rho q/n\asymp g^{-1/9}\).

A path crossing chambers changes the average of the parallel winding traces by a fixed amount in at least one direction. A sum of squares of the actual time-space plaquette deviations controls the difference of this average between adjacent slices. Cauchy--Schwarz along the time arc gives curvature cost \(cn/q\). Multiplying by \(\beta_m\) gives cost \(g^{-4/3+o(1)}\), while the logarithmic cost of a fixed rescaled center ball in the same normalization is \(O(g^{-1+o(1)})\). Chamber crossings are thus suppressed beyond all powers. The spatial center cocycles permute the eight centers, giving equal center weights. The comparison \eqref{eq:paper-outside} holds first under the common analytic action and transfers to the original field by \eqref{eq:paper-density} and the long-cylinder failure bound. See Appendix~\ref{app:cam-soft}.

\subsection{The low-energy model in a center neighborhood}

Use \(c=g_m^{2/3}X\) and the cutoff radius \(M=C_A(1+|\log g|)\) in \(X\). For any prescribed \(A\), both the actual constant-mode path and the model bridge have mass \(O(g^A)\) outside this ball. Let \(A_M=g_m^{2/3}M\) and \(u_c=A_M/n\). A temporal convolution makes the electric quadratic form a sum of squares of nearest-neighbor differences. We call this change of variables whitening. Write \(Q\) for the whitened electric quadratic form, \(Q_B\) for the quartic magnetic action, and \(R\) for the remainder. Then
\begin{equation}
 |R|\le C_g\big[(u_c^2+g^{4/3})(Q+Q_B)+u_c\sqrt{QQ_B}\big],
      \qquad C_g=g^{-o(1)} .
 \label{eq:paper-action}
\end{equation}
The static nonzero-mode factor contributes a further error \(O(g^{5/12-o(1)}M^2)\), and the log Haar density changes by \(O(q u_c^2)=O(g^{3/4}M^2)\). The square-root error from the nonlinear fast-mode integration is at most \(O(g^{5/12-o(1)})\). With Young's parameter \(g^{3/8}\), the electric remainder after multiplication by \(q\) is \(O(g^{3/8-o(1)}(1+M^2))\), and the magnetic remainder is \(O(g^{3/8}M^4)\). With the common small-mass reference, \(\mathbb E[(1+Q)^2e^{C\delta_e Q}]\le Cq^2\), where \(\delta_e\) bounds the relative error in the electric form. Applying \(|e^{-R/2}-1|^2\le R^2e^{|R|}/4\) gives action error \(O(g^{3/8-o(1)})\) for the square roots.

The complete coarse observation retains the nonzero spatial modes. The relative covariance error for integer Fourier block \(\mathbf k\ne0\) is \(O(g^{-o(1)}n\rho/|\mathbf k|^2)\). Summing \(|\mathbf k|^{-4}\) over \(\mathbb Z^3\setminus\{0\}\) and retaining both complete time bands gives a squared Hellinger error for the Gaussian marginal of
\begin{equation}
 O\bigl(g^{-2\delta-o(1)}(n\rho)^2\bigr)
       =O(g^{17/18-2\delta-o(1)}).
 \label{eq:paper-normal}
\end{equation}
Conditional generation of the observation retains both nonzero-mode Gaussian noise and the Haar variables on the gauge orbit. Decay of the same nonzero-mode inverse controls the the effects of the actual temporal holonomy and residual correlations between ends.

After factoring out \(\beta_m n/8\), the temporal electric symbol of the spatial zero mode is \(\omega(\theta)^2 a_g(\theta)\), where \(\theta\) is time frequency and \(\omega(\theta)=2|\sin(\theta/2)|\). The blocked curvature symbol and the zero-frequency moments of the inherited quadratic loop term give \(a_g(0)=1\). Truncate the whitening convolution and its inverse at \(g^{-\delta'}\) lattice steps, with \(0<\delta'<\delta\). Analyticity in a strip gives a tail \(\exp[-g^{-\delta'+o(1)}]\), which controls the growth of the model kernel's Lipschitz constant on a logarithmically growing ball. The dimensionless time step is \(d_s=g_m^{2/3}/n\asymp g^{3/4}\). Odd terms have zero expectation under the short Brownian bridge, giving an error of \(O(d_s(1+|\log d_s|)^3)\) in the square roots of path densities. The whitening displacement in the square root of the nonnegative magnetic weight costs \(O(d_s^{1/2}g^{-o(1)})\). The four short model propagations at the two ends give \(O((d_s b^\delta)^{1/2})=O(g^{3/8-\delta})\).

The finite-range inverse convolution recovers the center configuration and the color orbit of the constant mode from each complete endpoint observation. Let \(F\ge0\) be the normalized model half-period kernel and take conditional expectations under the reference law in which the two endpoint observations are generated independently. Then
\[
 \left\|F-\sqrt{\mathbb E[F^2\mid u_1,u_2]}\right\|_2^2
       \le 2\,\mathbb E\operatorname{Var}(F\mid u_1,u_2).
\]
This controls the complete coarse observation kernel. See Appendix~\ref{app:cam-soft}.

Both the fine-slice and the nine-dimensional model comparisons use \(C_{N,m}\). Identify its common subspace and take the orthogonal direct sum of the two complements. The resulting equivariant Hilbert--Schmidt error is
\begin{equation}
 \begin{split}
  &O(b^{-3/10+o(1)})
   +O(b^{-3/16+\delta/2+o(1)})
   +O(g^{3/8-\delta-o(1)})+O_A(b^{-A})\\
  &\hspace{28mm}=O(b^{-3/16+\delta/2+o(1)})=O(b^{-1/8}).
 \end{split}
 \label{eq:paper-total}
\end{equation}
All estimates hold together on the same cutoff tail \(N\ge N_*(b)\), uniformly over the chosen real terminal band and compact dimensionless time interval. Rounding the physical time to lattice steps changes the dimensionless time by \(O(g^{2/3})\), and replacing the stopping-scale inverse coupling by \(v\) changes it by a relative \(O(\log b/b)\). Both errors are within \eqref{eq:paper-total}.

\section{The nine-dimensional model and the Wilson low-energy spectrum}
\label{sec:soft-spectrum}

\subsection{Physical time and model normalization}

In the Lie algebra metric \(\langle A,B\rangle=-\operatorname{Re}\Tr(AB)/2\), the basis satisfies \(\|t_a\|^2=1/4\). Write \(c_i\cdot t=\sum_a c_i^a t_a\) for the dimensionless spatial constant mode and \(V_4(c)=\sum_{i<j}|c_i\times c_j|^2\). The physical time step at the stopping level is \(a_m=L_0/n\). The leading spatial curvature is \([c_i\cdot t,c_j\cdot t]/n^2\), and the leading electric curvature is \((\delta_t c_i)\cdot t/n\), where \(\delta_t\) is the adjacent-slice difference. The zero-frequency normalization of the curvature term and the vanishing second moments of the closed paths give electric and magnetic coefficients \(\beta_m n/8\) and \(\beta_m/(8n)\), respectively (Appendix~\ref{app:cam-soft}). One electric step has diffusion variance \(g_m^2/n\). Converting diffusion and the magnetic potential to physical time gives
\begin{equation}
 L_0H_{\mathrm{eff}}
   = \frac{g_m^2}{2}(-\Delta_c)+\frac{V_4(c)}{2g_m^2}
   = \frac{g_m^{2/3}}{2} H_4,\qquad
 H_4=-\Delta_X+V_4(X),\qquad c=g_m^{2/3}X .
 \label{eq:model-scale}
\end{equation}
Gauss's law restricts the model to color singlets. At the actual stopping-scale coupling the dimensionless time is \(u_m=Tg_m^{2/3}/L_0\). Since \(\beta_m=v+O(\log b)\), replacing it by \(u=Tg^{2/3}/L_0\) incurs a relative time error \(O(\log b/b)\). The model half-period operator on the eight center copies, normalized by the full-period trace, is
\begin{equation}
 A_*(u)=\frac{1}{\sqrt 8}\,I_{\mathrm{cent}}\otimes
         \frac{e^{-uH_4/4}}
              {(\operatorname{Tr}_{\mathrm{sing}}e^{-uH_4/2})^{1/2}} .
 \label{eq:model-root}
\end{equation}
Here \(\operatorname{Tr}_{\mathrm{sing}}\) is the trace on color singlets and \(I_{\mathrm{cent}}\) is the identity on the eight-dimensional center regular representation. The Hilbert--Schmidt norm of \(A_*(u)\) is one.

\subsection{Quantum confinement and parity}

The quartic potential can vanish when the three columns are parallel. Transverse zero-point energy makes the spectrum discrete. For fixed \(X_i\), split \(X_j\) into components parallel and perpendicular to \(X_i\). The two-dimensional transverse oscillator yields
\[
 -\Delta_{X_j}+|X_i\times X_j|^2\ge 2|X_i|\qquad(i\ne j)
\]
as a form inequality. Half the sum over the six ordered pairs gives \(H_4\ge2\sum_i|X_i|\). Averaging this with \(H_4\ge-\Delta_X\) gives \(H_4\ge-\tfrac12\Delta_X+\sum_i|X_i|\). This bound controls local Sobolev norms and the \(L^2\) tail at infinity, so Rellich compactness gives a compact resolvent. The comparison operator has finite heat trace at all positive times. The Golden--Thompson bound integrates the product of the constant free heat-kernel diagonal and \(\exp(-t\sum_i|X_i|)\). The Feynman--Kac kernel is strictly positive, so the ground state is simple and strictly positive. Color rotations and inversion \(X\mapsto-X\) preserve it. The first two distinct singlet eigenvalues thus satisfy \(\nu_1-\nu_0>0\).

Exact zero momentum requires the parity of the whole first excited eigenspace. View \(X\) as a \(3\times3\) matrix with Gram matrix \(\mathsf G=X^{\mathsf T}X\). A smooth odd color singlet vanishes on \(\det X=0\), since a simultaneous color half-turn sends all three columns to their negatives when they lie in a plane. On the full-rank set it has the form \(\psi(X)=\det X\,F(\mathsf G)\). The Gram pushforward of Lebesgue measure for real \(d\times3\) matrices has density proportional to \((\det\mathsf G)^{(d-4)/2}\). Multiplication by \(|\det X|^2=\det\mathsf G\) takes the \(d=3\) density to the \(d=5\) density. Let \(\Delta_{(d)}\) be the Laplacian on \(\mathbb R^{d\times3}\). On Gram functions its row-dependent drift is \(2d\operatorname{tr}\partial_{\mathsf G}\), the sum of the diagonal derivatives multiplied by \(2d\). The determinant is harmonic, and the product rule contributes \(4\det X\operatorname{tr}\partial_{\mathsf G}F\). Thus
\[
 \Delta_{(3)}\big((\det X)F(\mathsf G)\big)=(\det X)\Delta_{(5)}F(\mathsf G).
\]
The potential is unchanged since \(V_4(X)=\sum_{i<j}(\mathsf G_{ii}\mathsf G_{jj}-\mathsf G_{ij}^2)\). The form identity holds first on a compactly supported core away from rank loss. A compactly supported smooth odd singlet satisfies \(|\psi(X)|\le C\sigma_{\min}(X)\), with \(\sigma_{\min}\) the least singular value. On cutting off near \(\sigma_{\min}=0\), the gradient of the cutoff times \(\psi\) tends to zero in \(L^2\) because the rank-deficient set is Lebesgue-null. Hence the comparison form bound holds on the entire odd singlet form domain.

Let \(H^{(5)}\) be \(-\Delta_{(5)}\) plus the same Gram potential on \(\mathbb R^{5\times3}\). Retain half the kinetic energy for each column and pair the remainder cyclically with the potential. With each axial column fixed, the four-dimensional transverse oscillator yields
\[
 H^{(5)}\ge \sum_{i=1}^3\left(-\tfrac12\Delta_{X_i}+2\sqrt 2\,|X_i|\right).
\]
Radial reduction and rescaling in the lowest angular momentum give \(4^{1/3}\) times the Dirichlet half-line operator \(-\partial_r^2+2/r^2+r\). For \(-\partial_r^2+r\) use \(\phi(r)=r\exp[-2r^{3/2}/(3\sqrt2)]\). The positive-function transform gives a lower bound of \(2\), since
\[
 -\phi''(r)/\phi(r)+r=\tfrac12 r+\tfrac5{2\sqrt2}r^{-1/2},\qquad
 \left(\inf_{r>0}[-\phi''(r)/\phi(r)+r]\right)^3=\tfrac{675}{64}>8.
\]
Split the linear potential as \((1-a)r+ar\). Rescaling gives \(2(1-a)^{2/3}\) for the first part with the kinetic term. Minimizing \(2/r^2+ar\) gives \(3a^{2/3}/2^{1/3}\) for the second. Optimizing over \(a\in[0,1]\) yields \((43/2)^{1/3}\) for the half-line operator, and hence a lower bound \(3\sqrt[3]{86}\) on the odd singlet spectrum.

Use the even trial functions \(e^{-\alpha|X|^2/2}\) and \((|X_1|^2-|X_2|^2)e^{-\alpha|X|^2/2}\). Exchanging the first two columns makes both the inner product and cross form vanish. The Rayleigh quotients are \(9\alpha/2+9/(2\alpha^2)\) and \(13\alpha/2+15/(2\alpha^2)\). For \(\alpha=(30/13)^{1/3}\), the cube of the larger quotient is \(68445/32<27\cdot86\). By the min--max principle both \(\nu_0\) and \(\nu_1\) lie below the odd spectrum. Simplicity of \(\nu_0\) places \(\nu_1\) second when eigenvalues are counted with multiplicity, regardless of the multiplicity of \(\nu_1\) itself.

\subsection{Exact zero momentum and the ordered limits}

The common isometric embeddings preserve center transformations, inversion, and stopping-scale translations (Appendix~\ref{app:cam-soft}). In the zero-electric-flux sector, \eqref{eq:model-root} retains the symmetric combination of the eight copies. The model eigenvalue at \(\nu_1\) gives an isolated heat eigenvalue. Choose thresholds between it and the neighboring heat eigenvalues. Their distances from the model spectrum are bounded below uniformly on the compact dimensionless time interval. The Hilbert--Schmidt comparison controls the operator norm, and min--max controls the corresponding heat eigenvalues. The isolated spectral projections converge in norm. The \emph{first excited cluster} of the original Wilson operator is its full spectral subspace between the two thresholds. At sufficiently weak coupling the equivariant projections differ by less than one, so the dimension and the inversion and stopping-scale translation actions agree with the model.

Let \(\mathsf T_i\) be the fine one-site translation in direction \(i\), \(J\) the fine inversion, and \(r_N=D_N/D_m\) the ratio from the fine grid to the stopping scale. The model cluster is even under inversion and trivial under stopping-scale translations. Thus \(J=\mathsf T_i^{r_N}=I\) on the Wilson cluster. The relation \(J\mathsf T_iJ=\mathsf T_i^{-1}\) gives \(\mathsf T_i^2=I\). The oddness of \(r_N\) then gives \(\mathsf T_i=I\). The full cluster has exact zero total momentum.

Write \(\lambda_{0,N},\lambda_{1,N}\) for the vacuum and largest excited heat eigenvalues of \(A_N(T)\) in the zero-electric-flux sector. They are uniformly bounded away from zero on the chosen time interval and differ from the corresponding values in \eqref{eq:model-root} by \(O(b^{-1/8})\). Consequently
\begin{equation}
 m_N^{00}
 =-\frac2T\log\frac{\lambda_{1,N}}{\lambda_{0,N}}
 =\frac{g^{2/3}}{2L_0}(\nu_1-\nu_0)
     +O\!\left(\frac{g^{2/3}}{L_0}\,b^{-1/8}\right).
 \label{eq:original-gap-readout}
\end{equation}
The factor \(2/T\) is the reciprocal of the evolution time \(T/2\).

On the fixed short square torus, fix \(b,v\) and an auxiliary observation level \(f\). Its joint law converges along the full cutoff sequence (Appendix~\ref{app:cam-fixed-law}). The conditional comparison at both ends bounds corresponding ordered eigenvalue differences at two fine cutoffs by \(C_{b,v}n_f^{-1/8}\) plus an error vanishing as both cutoffs are removed. Sending \(f\to\infty\) gives the Cauchy property. With \(t_0=L_0/2>s_*\), each state of energy at most \(E\) contributes at least \(e^{-t_0E}\) to \(Z_N^0(t_0)\). The short-time bound thus controls the eigenvalue count in fixed energy windows uniformly in the cutoff. It also bounds the short normalized vacuum eigenvalue away from zero, and \eqref{eq:original-gap-readout} provides an upper energy bound that does the same for the first excited heat eigenvalue. First take the full cutoff limit of the short-period eigenvalue ratio at fixed \(b,v\). The long-period error \eqref{eq:original-gap-readout} on the same Wilson trajectory then passes to this energy difference. This proves \eqref{eq:main} with the common cutoff tail fixed for the entire real terminal band.

\section{Continuous-time vacuum representation on the spatial torus}
\label{sec:vacuum-reports}

Fix \(b,v,L_0\) on the same Wilson bare sequence. Operator comparisons in this section are restricted to \(\mathcal H_N^0\), using the same notation for \(H_N,S_N(t)\), and \(\Omega_N\). The heat traces \(Z_N^0(t)\) remain on the full Gauss-law space. Set \(t_0=L_0/2\). The entire first excited cluster has exact zero momentum, so the first gap in \(\mathcal H_N^0\) equals \(m_N^{00}\) on a sufficiently fine tail. Write \(m_{\mathrm{gap}}=m^{00}(L_0,v)>0\), and choose \(0<m_{\mathrm{lb}}<m_{\mathrm{gap}}\) with \(m_N^{00}\ge m_{\mathrm{lb}}\) on that tail.

\subsection{Vacuum correlations and the strongly continuous semigroup}

Choose a bounded, gauge- and center-invariant Wilson observable \(F\) at fixed nonzero physical resolution. It is a function of links from one or more fixed blocking levels and can also be specified by Wilson loop traces at those levels. Place the observation in the central time band of the short half-cylinder so that its fine-lattice dependence domain is at distance at least \(d_F\) from both ends, with \(2d_F>s_*=L_0/4\). The unmarked intervals adjacent to the observation then have uniform short-time heat-trace bounds. Let \(\mathcal T_N(F)\) be the vacuum-shifted half-cylinder operator with the insertion \(F\), and let \(B_N(F)\) be the insertion operator with the two free intervals removed. Time reflection and complex conjugation give the adjoint. The positive path integral gives
\begin{equation}\label{eq:rep-buffer}
 \mathcal T_N(F)=S_N(d_F)B_N(F)S_N(d_F),
 \qquad \|B_N(F)\|_{\op}\leq\|F\|_\infty .
\end{equation}

The short half-period operators with and without an insertion use the same endpoint embeddings. Keep the observation's blocking random variables and their complete fine-lattice dependence domain. Let \(p(x,y,\gamma)\) be the joint density of the endpoints and observation \(\gamma\), and \(p_0=\int p\,d\gamma\) its endpoint marginal. Use \(\widetilde p,\widetilde p_0\) for a reference law. The insertion kernels satisfy
\begin{equation}\label{eq:rep-mark-inequality}
 \left\|\int \left|\frac{p}{\sqrt{p_0}}-\frac{\widetilde p}{\sqrt{\widetilde p_0}}\right|d\gamma\right\|_{L^2}
    \leq 3\|\sqrt p-\sqrt{\widetilde p}\|_{L^2}.
\end{equation}
Set terms with zero denominator to zero. At each endpoint pair write the densities as products of marginal and conditional densities. The marginal term gives \(|\sqrt{p_0}-\sqrt{\widetilde p_0}|\). The conditional term is bounded by twice the conditional Hellinger distance times the smaller square root of the marginal. Both are controlled by the square-root distance of the joint laws. At fixed auxiliary level \(f\), the coarse observation kernels with and without an insertion converge along the full cutoff sequence. In the common embedding, the original and fixed-level kernels differ in Hilbert--Schmidt norm by at most \(\varepsilon_f\|F\|_\infty\) in the limit superior, with \(\varepsilon_f\to0\). The comparison is made for the positive law on the full Gauss-law space and then restricted by equivariance to zero electric flux. The top eigenvalue of the unmarked short operator is uniformly bounded away from zero, giving a common vacuum normalization. See Appendices~\ref{cam-bound-section} and~\ref{app:cam-vacuum}.

Let \(W_N\) be a finite product of operators \(\mathcal T_N(F)\), their adjoints, and \(S_N(t)\) at fixed positive times. We call such an ordered product an observation word. The identity \(S_N(t)=S_N(t_0)^{t/t_0}\) extends the comparison by continuous functional calculus. In the zero-flux sector the vacuum projection of \(S_N(t_0)\) is an isolated rank-one projection, with other eigenvalues at most \(e^{-t_0m_{\mathrm{lb}}}\). The common operator norm approximation controls the difference of the vacuum projections and, by rank one, their trace norm difference. Hence \(\langle\Omega_N,W_N\Omega_N\rangle\) converges along the full cutoff sequence. Compare two fine cutoffs at fixed auxiliary level \(f\), then send \(f\to\infty\). The limiting positive vacuum functional defines a semidefinite inner product on words. Quotienting by its null space and completing gives \(\mathcal H_{\mathrm{rep}}\) and vacuum \(\Omega\).

The positive contraction semigroup laws pass to the limit. By \eqref{eq:rep-buffer},
\[
 \|(S_N(s)-I)\mathcal T_N(F)\|_{\op}
   \leq \|F\|_\infty
       \sup_{E\geq0}(1-e^{-sE})e^{-d_FE}\longrightarrow 0
          \qquad(s\downarrow0),
\]
uniformly in the fine cutoff. Adjoints, words starting with a positive free interval, and the vacuum obey the corresponding continuity estimates. The limiting semigroup is strongly continuous, with nonnegative self-adjoint generator \(H_{\mathrm{rep}}\). Reflection inner products pass from the original path integral. The vacuum moments also arise from an infinitely long time cylinder at fixed spatial volume. For \(l>s_*\) the short heat-trace bound gives
\[
 \|S_N(l)|_{\mathcal H_N^0}-|\Omega_N\rangle\langle\Omega_N|\|_1
       \leq e^{-m_{\mathrm{lb}}(l-s_*)}Z_N^0(s_*).
\]
Append a free interval of length \(l\) to the fixed observation word. After removal of the cutoff and then letting \(l\to\infty\), the zero-flux periodic trace gives the same vacuum moment. Passing the finite-lattice spectral inequality to centered word vectors gives
\begin{equation}\label{eq:rep-lower}
 H_{\mathrm{rep}}\big|_{\Omega^\perp}\geq m_{\mathrm{gap}} .
\end{equation}

\subsection{Visibility and multiplicity of the first cluster}

Let \(\psi_0>0\) be the normalized model ground state and \(\{\psi_j\}_{j=1}^d\) a real orthonormal basis of the full first excited eigenspace. Take a fixed smooth compactly supported cutoff \(\chi\), invariant under color rotations and inversion and equal to one on a sufficiently large ball. The functions \(f_i=\chi\psi_i/\psi_0\) are bounded and Lipschitz, and the matrix
\[
 \bigl(\langle\psi_j,f_i\psi_0\rangle\bigr)_{j,i=1}^d
\]
is invertible because it tends to the identity as the cutoff ball grows. On a central stopping-scale slice take \(F_i=f_i(g^{-2/3}c)\) with the same function in each of the eight center charts, extended by zero outside and averaged over center transformations. The observation and its level are fixed once \(b,v\) are fixed.

The long-period comparison embeds the fine and model slice spaces in a common space containing the observation variable. In each observation fiber identify the unit lines generated by the respective square roots of the conditional densities and take the direct sum of the complements. Multiplication by \(F_i\) is scalar in each fiber and is preserved. Let \(Y_s\) be the stationary model ground-state process and write the local time-convolution output as \(X_0=\sum_j a_jY_{jd_s}\), where \(d_s=g_m^{2/3}/n\), \(\sum_j a_j=1+o(1)\), and \(\sum_j |a_j|\sqrt{|j|}=g^{-o(1)}\). The bound \(\mathbb E|Y_s-Y_0|^2\le C|s|\) follows from the ground-state Dirichlet form identity \(\langle y_{ia}\psi_0,(H_4-\nu_0)y_{ia}\psi_0\rangle=\|\psi_0\|^2\), with \(y_{ia}(X)=X_i^a\), using bounded truncations and integration by parts. Consequently
\[
 \|X_0-Y_0\|_{L^2}
 \le C\sqrt{d_s}\sum_j|a_j|\sqrt{|j|}
       +\left|\sum_j a_j-1\right|\|Y_0\|_{L^2}=o(1).
\]
At a fixed positive model time the heat-kernel diagonal contains a positive ground-state contribution, so the endpoint tail bound also controls the vacuum path tail. In the model the fixed-terminal observation is \(f_i((g_m/g)^{2/3}X_0)\), which tends to \(f_i(Y_0)\) by the convolution bound and \(g_m/g=1+O(\log b/b)\). Norm convergence of the vacuum and cluster projections in the common space transfers the strictly positive least singular value of the model observation matrix to the Wilson cluster. See Appendices~\ref{app:cam-soft} and~\ref{app:cam-vacuum}.

To move the observation into the short half-cylinder, let \(\pi_N\) be the ground-state probability of density \(|\Omega_N|^2\) relative to the slice reference measure, put \(p_\varphi=\varphi/\Omega_N\), and let \(\mathcal U_s\) be the stationary slice process at physical time \(s\). For \(\varphi\) in the first excited cluster with largest cluster energy \(E_{\max,N}\),
\[
 \mathbb E_{\pi_N}|p_\varphi(\mathcal U_s)-p_\varphi(\mathcal U_0)|^2
 =2\langle\varphi,(I-S_N(s))\varphi\rangle
 \le 2(1-e^{-sE_{\max,N}})\|\varphi\|^2 .
\]
At \(s=t_0/2\), the bound \(E_{\max,N}=O(g^{2/3}/L_0)\) makes the matrix element error at most \(O(g^{1/3})\|\varphi\|\). Let the \(i\)th column of \(V_N\colon\mathbb C^d\to\mathcal H_N^0\) be \(\mathcal T_N(F_i)\Omega_N\) and let \(P_{\mathrm{cl},N}\) project onto the full cluster. There is a \(\gamma>0\) such that at sufficiently weak coupling, on the same fine cutoff tail,
\begin{equation}\label{eq:rep-frame}
 P_{\mathrm{cl},N}V_NV_N^*P_{\mathrm{cl},N}
       \geq \gamma P_{\mathrm{cl},N}.
\end{equation}
In particular the bound holds on any lowest subcluster produced by splitting.

For fixed \(b,v\), choose a continuous cutoff \(\eta\) equal to one near an isolated energy of interest and zero near the other levels and at high energies. Continuous functional calculus for the short heat operators gives
\[
 V_N^*\eta(H_N)V_N\longrightarrow V^*\eta(H_{\mathrm{rep}})V,
 \qquad V=(\mathcal T(F_i)\Omega)_i ,
\]
where \(\mathcal T(F_i)\) is the limiting insertion. Take windows shrinking to \(m_{\mathrm{gap}}\) whose boundaries avoid limiting eigenvalues of the ordered short heat operators. The positive eigenvalue bound in \eqref{eq:rep-frame} passes to the spectral-weight matrix at \(m_{\mathrm{gap}}\). Conversely, if the limiting isolated energy window had more independent spectral vectors, one could approximate them by vacuum observation words and transfer their Gram matrix back to a sufficiently fine lattice using the same continuous cutoff. This would exceed the lattice eigenvalue count in the window. The first energy and multiplicity are thus preserved. On a fixed blocking level, contractible Wilson loop traces separate gauge and center orbits on a set of full Haar measure. The last blocking step gives a uniform joint Haar density bound at that level. Bounded continuous functions of finitely many such traces therefore approximate the \(F_i\), and conditional Jensen's inequality preserves a strict margin in \eqref{eq:rep-frame}. See Appendix~\ref{app:cam-vacuum}.

\subsection{The renormalized local-field subspace}

Take a renormalized local insertion \(\mathcal O_N\) satisfying the separated chessboard estimate in the same short box. Place the insertion and its reflection in separate cells at a fixed positive distance from the cell boundaries. The blocking estimate gives a bounded approximation \(F_{\varepsilon,N}\) at a fixed level such that
\begin{equation}\label{eq:rep-local-approx}
 \limsup_{N\to\infty}
       \|\mathcal T_N(\mathcal O_N)-\mathcal T_N(F_{\varepsilon,N})\|_{\HS}
     \leq \varepsilon\sup_N\sqrt{\Tr S_N(L_0)} .
\end{equation}
The norm is evaluated by the reflected, separated two-point function. The argument uses the single-insertion replication and coarse approximation of \cite[Lemmas 14.14 and 14.29]{armstrong}, with disjoint blocking dependence domains verified on the chosen short torus. The approximation coefficients are bounded in a fixed finite-dimensional space. If their subsequential limits give \(F_\varepsilon,F'_\varepsilon\), take the limit of the true mixed moment
\(\Tr(\mathcal T_N(F_\varepsilon-F'_\varepsilon)^*\mathcal T_N(\mathcal O_N))\)
on both subsequences. Equation~\eqref{eq:rep-local-approx} and Cauchy--Schwarz control the difference of the candidate operators. Hence one can fix a candidate and send \(\varepsilon\downarrow0\) to define the buffered local-field operator on \(\mathcal H_{\mathrm{rep}}\). A telescoping estimate applies to finite products with positive time separations, and the compressed matrices between fixed isolated spectral windows converge. See Appendices~\ref{cam-bound-section} and~\ref{app:cam-vacuum}.

Let \(\mathcal H_{\mathrm{loc}}\) be the closed space generated from \(\Omega\) by these separated local words and positive heat times. Invariance under the self-adjoint heat semigroup makes it reducing. Its Hamiltonian satisfies \(H_{\mathrm{loc}}|_{\Omega^\perp}\ge m^{00}(L_0,v)\) by \eqref{eq:rep-lower}. For the local-field gap to have the same leading asymptotic, one needs positive projection of a local-word vacuum vector onto the actual first excited cluster. Choose a continuous cutoff \(\eta_{\mathrm{soft}}\) on dimensionless energy, equal to one near \((\nu_1-\nu_0)/2\) and zero near the vacuum and all other model levels. At weak coupling it is one on the entire actual first cluster. The desired spectral weight is
\begin{equation}\label{eq:rep-local-weight}
     \lim_{N\to\infty}
       \big\|\eta_{\mathrm{soft}}(L_0g^{-2/3}H_N)
                    W_{N,g}\Omega_N\big\|^2>0 .
\end{equation}
The word \(W_{N,g}\) and its positive time separations must be specified at fixed \(b,v\) before removing the cutoff. In the model set \(B=(X_2\times X_3,X_3\times X_1,X_1\times X_2)\). The six Gram entries \(q_{ij}=(B^{\mathsf T}B)_{ij}\) separate full-rank even singlet orbits because \(B^{\mathsf T}B=\operatorname{adj}(X^{\mathsf T}X)\). Let \(K=(H_4-\nu_0)/2\). Since \(|q_{ij}|\le |B|^2\), the operators \(e^{-sK}q_{ij}e^{-sK}\) are bounded and self-adjoint. The cyclic subspace generated by them and the heat semigroup has projection commuting with both. Sending \(s\downarrow0\) on the energy form domain gives commutation with the multiplication forms. Smooth compactly supported functions are form cores for the nonnegative weights \(|B|^2\) and \(|B|^2+q_{ij}\), giving commutation with the bounded functional calculus of the six entries. Strict ground-state positivity makes their bounded functions times the ground state dense in the even singlet space. Hence a finite family of separated model words has a full-rank observation matrix on the model first cluster. Establishing \eqref{eq:rep-local-weight} for the continuum local fields requires the corresponding sourced Wilson estimates. Exact identification of the lowest actual excited energy and its multiplicity requires full rank on its lowest subcluster. See Appendix~\ref{app:cam-vacuum}.

\section{The nonabelian mechanism, finite volume, and the trace anomaly}
\label{sec:discussion}

\paragraph{Maxwell comparison.}
In compact Maxwell theory on the same periodic spatial box, the flat spatial holonomy angles are uniform in the unprojected heat trace. Photon frequencies are independent of these angles, so there is no analogue of the suppression from the four charged transverse oscillators. In the sector of zero magnetic flux, zero electric flux, and exact zero total momentum, the flat-angle rotor wave function is constant. The least nonzero photon momentum costs \(2\pi/L_0\). The lowest nonvacuum state at zero total momentum contains two transverse photons with opposite momenta and energy \(4\pi/L_0\). This gives a separate finite-volume mechanism for a positive gap. The \(g^{2/3}\) law here arises from the commutator magnetic potential and quantum motion of the spatial constant modes.

\paragraph{Variation of a thermal partition function.}
The infinite-volume gap problem involves changing the physical scale. For a Wilson family at fixed lattice spacing, write \(Z(L,T)\) for the unnormalized partition function with spatial side \(L\) and time period \(T\). Let \(A_\alpha\) be the plaquette-action sum for orientation class \(\alpha\) and \(\beta_\alpha(r)\) its differentiable bare coefficient. Differentiating the finite-dimensional integral gives
\[
 \frac{d}{dr}\log\frac{Z(L,L/2)^2}{Z(L,L)}
 =\sum_\alpha\dot\beta_\alpha(r)
 \left(\langle A_\alpha\rangle_{L,L}
               -2\langle A_\alpha\rangle_{L,L/2}\right).
\]
This is related to the integral method of lattice thermodynamics, where differences of action expectations determine partition functions \cite{thermodynamics}. When \(r\) denotes a continuum physical scale, the insertion on the right needs identification in a common continuum renormalization. The trace anomaly is the natural candidate for that scale insertion, and its sign in this difference depends on the two thermal expectations. Positivity for the four-dimensional \(a\)-theorem comes from a distinct spectral and scattering argument \cite{a-theorem}. Theorem~\ref{thm:main} uses only comparisons at fixed physical spatial size.

\paragraph{Scope.}
Equation~\eqref{eq:main} quantifies the commutator zero-mode gap in the fixed box, and Proposition~\ref{prop:representation} identifies it in the continuous-time representation of bounded Wilson observables. Equation~\eqref{eq:rep-lower} gives the lower bound inherited by the specified renormalized local fields, and \eqref{eq:rep-local-weight} specifies the extra positive weight required for their gap asymptotic. An infinite-volume result requires estimates uniform in growing physical volumes. The chosen weak-coupling schedule gives a sufficiently large \(b\) in asymptotic terms. Its admissibility and error bounds have no numerical onset in this argument.

\section*{Acknowledgments}
\addcontentsline{toc}{section}{Acknowledgments}
We acknowledge assistance from Claude Fable 5.2, GPT-6.2 Sol, and GPT-6.2 Astra in developing, checking, and presenting the argument.

\paragraph{Supplementary checks.}
The companion scripts check the weak-coupling exponents and parity bounds by exact rational arithmetic, the time-folding identity on finite examples, and the Cartan coarea, bridge, and endpoint-kernel calculations on finite models. Run each script in \path{checks/} with Python 3.10 or later. The coarea check uses NumPy. The finite examples illustrate the analytic identities and estimates stated in the paper.

\appendix
\section{Inputs from the public blocking construction}
\label{app:inputs}

Positive transfer and gluing use the Wilson plaquette weight and its Gauss projection \cite{luscher-transfer,osterwalder-seiler}. Definitions 5.1 and 5.2 in \cite{armstrong} specify the exact positive blocking kernel and its blocked density. Theorems 5.5--5.6 and Lemmas 13.23 and 14.4 give the real terminal shooting and layerwise drift used here. Lemmas 12.3 and 12.4 supply the absolute path weights and the zero-frequency electric and magnetic normalization, respectively. Proposition 5.7 gives convergence at each fixed layer on the original equal-period tori at the center of the terminal band. Lemma 13.24 supplies the fixed-layer comparison for the selected real terminal band. Theorem 5.9 and Lemmas 14.14 and 14.29 concern separated local insertions and their coarse approximations. The local-field comparison here checks fixed dependence domains and separation margins in the actual short torus. The spectral argument uses the whole cutoff sequence at fixed volume.

The public shooting and weighted activity propagation in \cite{armstrong} concern the fixed equal-period box. The rectangular cylinder here retains the same bare coupling and actual positive blocking kernels. The comparison keeps the full dependence domains in both shapes, and the field-independent normalizations are canceled in the required pressure difference. Appendix~\ref{app:geometry} gives these steps. Propagation along the four short arms and the relative mass estimate outside the center neighborhoods are finite-cutoff estimates on the same bare trajectory. They appear in Appendices~\ref{cam-bound-section} and~\ref{app:cam-soft}. The fixed-volume vacuum representation uses these bounds and the fixed-level probability limit. The local-field comparison in Appendix~\ref{app:cam-local-mark} uses the public insertions.
\section{Blocking on rectangular tori and short-time heat traces}
\label{app:geometry}

\subsection{One-step conditional integration and local dependence}
\label{lem:geometry-local}

We use the blocking sequence of Theorem \ref{thm:main}. The integer \(N\) is the ultraviolet cutoff level, \(f\le N\) is the blocking level, and the odd integer \(\ell_f\) is the blocking ratio from level \(f+1\) to level \(f\). The cumulative blocking ratio is \(D_f=\prod_{i<f}\ell_i\), with \(D_0=1\). The reference square torus has \(P\) lattice steps along each side at level zero. At level \(f\), each spatial side has \(n_f=P D_f\) lattice steps and the physical lattice spacing is \(L_0/n_f\). Set \(s_f=\log D_f\). The terminal coupling band is centered at \(b\), and \(v\) is the actual terminal coupling in this band. Let \(\kappa_G>0\) be the coupling drift per unit logarithmic scale in \cite{armstrong}. The reference coupling is \(\widehat\beta_f=b+\kappa_Gs_f\), and the actual coupling on the square torus is \(\beta_{N,P,f}(v)\). Set
\[
 B_f^{\mathrm{cov}}=\lceil(\log\widehat\beta_f)^2\rceil,\qquad
 t_f=\widehat\beta_f^{-1/2}M_f,\qquad
 M_f=\exp\big((\log\log\widehat\beta_f)^2\big).
\]
Here \(B_f^{\mathrm{cov}}\) is the covering radius in lattice steps, \(t_f\) is the local small-field amplitude, and \(M_f\) is its subpower correction. For a set \(X\), let \(\operatorname{cov}_{B}(X)\) be the minimum number of lattice balls of radius \(B\) needed to cover \(X\) with connected union. The dependence set of a local factor includes every link needed to evaluate its thresholds, tags, paths, and conditional integrals. The root is the lattice site to which the factor is assigned. Distances are measured after lifting from that root to the integer lattice. The comparison uses local variables admitting the same lift on both tori. At level \(f\), a site incident to a plaquette of chord length greater than \(t_f\) carries a large-field flag. A support satisfies the small-field condition when no site in it carries such a flag. The block integral also selects threshold tests on fine plaquettes and coarse-link mismatches. These selections, together with the large-field flags at the output level, determine the local integrals.

Stage \(j\) consists of a range of levels that share blocking ratios and constants in the local estimates. We write \(C_j,c_j>0\) for constants depending only on these data, and \(\Pi_j\) for a polynomial majorant with the same dependence that may be increased from one estimate to the next.

We use the stochastic blocking and local integral estimates of \cite[Definition 5.1 and Lemmas 11.12--11.14 and 12.3--12.4]{armstrong}. Let \(U\) and \(V\) be the fine and coarse fields, \(\lambda\) the random blocking tag, and \(e\) a coarse link. The fine field determines a predictor \(Q_e(U,\lambda)\) for the coarse link. On the small-field branch it is the group element minimizing the weighted sum of squared chord distances from the local parallel transports. If the test for this branch fails, the tag selects one transport with the original path weights. Set \(\operatorname{pr}W=1-\tfrac12\operatorname{Re}\Tr W\). The conditional density of the coarse field given the fine field, relative to product probability Haar measure on the coarse links, is
\begin{equation}
 k_\beta(U,V)=\sum_\lambda \mathfrak p_\lambda(U)
       \prod_e q_{\mathsf b\beta}\big(V_e Q_e(U,\lambda)^{-1}\big),
 \quad q_\gamma(W)=\frac{e^{-\gamma\,\operatorname{pr} W}}
 {\int e^{-\gamma\,\operatorname{pr} W}\,dW}.
 \label{eq:geom-noise}
\end{equation}
Here \(\beta\) is the coupling used in this blocking step, \(\mathsf b>0\) is the fixed blocking penalty, and \(\mathfrak p_\lambda(U)\) is the original tag probability conditional on \(U\). Both geometries retain all admissible tags within the shared dependence range, with their respective conditional probabilities. For a fixed selection of local tests in the original integral, freeze the coordinates prescribed by the tests and write the quadratic energy in the remaining coordinates \(u\) as \(Q(u)=|D_{\rm loc}u+d|^2/2\). The matrix \(D_{\rm loc}\) comprises the linear constraints retained in the quadratic energy, and \(d\) includes the linear displacement of the coarse field. Write \(A=D_{\rm loc}^*D_{\rm loc}\) for the precision matrix, \(\xi=D_{\rm loc}^*d\) for the affine source, and \(d_{\rm act}\) for the number of active coordinates. Unused coordinates are padded with an identity block for matrix operations. Let \(\beta'\) be the coefficient of the quadratic energy in the local integral, and set \(g'=(\beta')^{-1/2}\). With \(P_{\rm loc}\) the local approximate inverse, set \(\mu=-P_{\rm loc}^*\xi\) and \(\Gamma=(P_{\rm loc}+P_{\rm loc}^*)/2\). Under \(u=\mu+g'w\), the normalization factor and the exponential correction retained inside the Gaussian integral are, respectively,
\begin{equation}
 (2\pi(g')^2)^{d_{\rm act}/2}(\det\Gamma)^{1/2} e^{-\beta' Q(\mu)},
 \qquad
 - (g')^{-1}(A\mu+\xi)\cdot w+\tfrac12 w^*(\Gamma^{-1}-A)w .
 \label{eq:geom-affine}
\end{equation}
Each coarse cell contains \(3\ell_f^4+1\) free fine links. The integral retains the Haar density factor at the identity for each free link and four single-link denominators from \eqref{eq:geom-noise} per coarse cell. When the output is the identity field \(V=\mathbf 1\) and no test is selected, \(\mu=0\), and the precision matrix is the periodization of a translation-invariant positive matrix. The field-independent logarithmic normalization on this branch consists of the flat Gaussian normalization and the remaining regular logarithmic terms from the same integral. The original integral has already removed the dimensions of the frozen coordinates, and this count is used again in forming the relative density.

Set \(B'=1+B_f^{\mathrm{cov}}+B_{f+1}^{\mathrm{cov}}\). The local matrix estimates in \cite[Lemmas 11.13--11.14]{armstrong} use a uniform spectral lower bound, weighted row and column sums, and a boundary residual for the local inverse. Retain or remove all three color indices of each link together. Let \(J\) be the retained index set and \(A_J\) the corresponding principal submatrix. For column \(k\), let \(\Omega_k\) be the original inversion window. Define column \(k\) of \(P_J\) by extending column \(k\) of \(A_{J\cap\Omega_k}^{-1}\) by zero to \(J\). The residual satisfies
\[
 (I-A_JP_J)_{zk}=
 \begin{cases}
 0,&z\in J\cap\Omega_k,\\
 -\displaystyle\sum_{x\in J\cap\Omega_k}
                      A_{zx}(A_{J\cap\Omega_k}^{-1})_{xk},
     &z\in J\setminus\Omega_k .
 \end{cases}
\]
The conjugation error for a diagonal weight restricts to the corresponding principal submatrix, and the spectral lower bound on the precision is preserved. Thus the original weighted inverse estimate applies to the restricted matrix, with a smaller range of summation in the boundary residual. Symmetrization gives
\[
 \Gamma_J=(P_J+P_J^*)/2,\qquad
 I-A_J\Gamma_J=
 \tfrac12\big[(I-A_JP_J)+A_J(I-A_JP_J)^*A_J^{-1}\big].
\]
Estimating the two terms in the energy norm of \(A_J\) gives
\(\|I-A_J^{1/2}\Gamma_J A_J^{1/2}\|\le\epsilon\),
where \(\epsilon\le\widehat\beta_{f+1}^{-c_j}\). More precisely, if \(B\) is the length scale of the local inversion window, the window contains a ball of radius \(cB\). The boundary residual also contains a finite-range matrix factor. Write \(\theta>0\) for the decay rate in the exponentially weighted inverse-column estimate. The original estimate bounds the residual by \(\widehat\beta_{f+1}^{o(1)}e^{-c_j'\theta B}\), where \(\theta B\) is at least a fixed multiple of \(\log\widehat\beta_{f+1}\). The larger window used in the original integral gives a stronger exponential estimate for the linear and quadratic terms in \(w\) in \eqref{eq:geom-affine}. For localizing the determinant, an inverse-power bound on the principal-submatrix residual suffices. The fourth- and sixth-order directional derivatives use the original polynomial bounds for matrix derivatives of fixed order.

For the normalization of \(\Gamma\), set \(F=I-A\Gamma\). Contraction in the energy norm gives the absolutely convergent identities
\begin{equation}
 \log\det\Gamma=-\log\det A-\sum_{h\ge1}\frac{\operatorname{Tr}F^h}{h},
 \qquad \Gamma^{-1}-A=\sum_{h\ge1}F^h A .
 \label{eq:geom-determinant}
\end{equation}
The representation \(A\mu+\xi=(I-AP_{\rm loc}^*)\xi\) specifies the local dependence of the linear correction. For its norm estimate one can use \(A(I-AP_{\rm loc})^*A^{-1}\xi\). The original integral retains separately the actual paths omitted from the finite-range truncation. The two families of linear constraints at the flat free minimizer satisfy an isometry identity, giving cancellation of the untruncated linear term. At local small-field amplitude \(t_f\), truncation of long paths incurs an error at most
\(\Pi_j(B)(t_f e^{-\tau_0R_*}+t_f^2)\), where \(R_*=\lfloor B/10^5\rfloor\) and \(\tau_0>0\) is a common exponential decay rate for the fixed-stage path kernels. The admission conditions for the original local estimate supply the required decay. In a local gauge, any test vector \(W\) gives the identity
\[
 D_{\rm loc}\mu+d=(I-D_{\rm loc}A^{-1}D_{\rm loc}^*)(D_{\rm loc}W+d)
                +D_{\rm loc}A^{-1}(A\mu+\xi).
\]
The two terms on the right are orthogonal in the Euclidean space of constraints. The quadratic energy of coordinates frozen by a local test stays with the original connected component.

To decompose \(-\log\det A\) into link contributions, assign independent continuous priorities to links, with the same priority for all three colors of a link. Fix the ordering and let \(J_i\) consist of the indices of link \(i\) and the links preceding it. Schur elimination gives
\[
       -\log\det A=\sum_i\log\det G_i,\qquad
       G_i=(A_{J_i}^{-1})_{ii}.
\]
The identity holds for every ordering, so it also holds after averaging over the priorities. For link \(i\), take the local diagonal block \(Q_i=(\Gamma_{J_i})_{ii}\) as reference. Fix one such index set \(J=J_i\) with precision \(A_J\) and residual \(F_J=I-A_J\Gamma_J\). Let \(t\) be the parameter along a normalized real field direction and use subscripts \(t\) for evaluation along that direction. Define
\[
 G_t^0=\Gamma_0^{-1/2}\Gamma_t\Gamma_0^{-1/2},\quad
 L_t^0=\Gamma_0^{1/2}(\Gamma_t^{-1}-A_{J,t})\Gamma_0^{1/2},
 \quad A_t^0=L_t^0G_t^0 .
\]
Here \(\Gamma_0=\Gamma_J|_{t=0}\),
\(A_t^0=\Gamma_0^{1/2}F_{J,t}\Gamma_0^{-1/2}\),
\(G_0^0=I\), and \(\|A_0^0\|\le\epsilon\).
Let \(\iota_i\) embed the three-dimensional coordinate block of link \(i\). The normalized embedding \(J_i^0=\Gamma_0^{1/2}\iota_i Q_i^{-1/2}\) is isometric. A complex parameter \(z\) records the degree in local matrix products. The normalized matrix for link \(i\) is
\[
 M_i(t,z)=(J_i^0)^*G_t^0(I-zA_t^0)^{-1}J_i^0.
\]
For \(|z|\le (4\epsilon)^{-1}\), \(\|M_i(0,z)-I\|\le1/3\). Taking at most six directional derivatives replaces at most six factors by directional derivatives of \(zA_t^0\) or \(G_t^0\), and the remaining inverses are uniformly bounded. With \([z^h]\) denoting the degree-\(h\) coefficient of the power series, the Cauchy estimate gives
\begin{equation}
 \left|\left.\partial_t^r[z^h]\log\det M_i(t,z)\right|_{t=0}\right|
       \le C_r\widehat\beta_{f+1}^{C_r}\epsilon^{-r}(4\epsilon)^h,
       \qquad 0\le r\le6 .
 \label{eq:geom-cauchy}
\end{equation}
The bound \(\widehat\beta_{f+1}^{-c_j}\) can be used to choose one Cauchy circle. Differentiation in \(t\) takes place on the same real small-field interval, while the complex circle concerns only \(z\). At \(z=1\) the logarithm is \(\log\det G_i-\log\det Q_i\). With the support count below, discarding terms above a sufficiently large fixed degree yields a weighted remainder bounded by \(\widehat\beta_{f+1}^{-2}\).

Fix a dimensionless constant \(L\). With the root as origin, partition the integer lattice into cells of side lengths between \(LB'\) and \(2LB'\), assigning any remaining sites to the far cells. Keep the root cell, which contains the full dependence set of \(Q_i\). For a set \(S\) of exterior cells, restrict \(J_i\) to indices covered by the root cell and \(S\), and compute the degree-\(h\) coefficient from \eqref{eq:geom-cauchy} on that restriction. Denote it by \(a_h(S)\). This auxiliary principal submatrix is used to expand the coefficient. The test regions selected in the original integral stay fixed. For a set \(T\) of exterior cells, take the finite difference
\[
         d_h(T)=\sum_{S\subset T}(-1)^{|T|-|S|}a_h(S).
\]
Inclusion-exclusion gives \(a_h(\text{all exterior cells})=\sum_Td_h(T)\). Each summand is bounded in absolute value by \(2^{|T|}\) times the right-hand side of \eqref{eq:geom-cauchy}. A product of local matrices in the coefficient gives a walk through link indices starting and ending at the root, of total length \(O(hB')\). Taking a point along the walk every \(LB'\) steps covers its dependence sets by a connected set of at most \(C(1+h/L)\) cells containing the root. Cells on which the product does not depend cancel in the finite difference. Thus \(T\) contributes only if it lies in such a connected set. On a graph of bounded degree there are at most \(C^s\) rooted connected sets of \(s\) cells, by encoding each with a \(2s\)-step traversal of a spanning tree. Counting subsets adds a factor \(2^s\). The \(B'\)-covering number of the connected cells and their local inverse windows is at most \(C L^4(1+h/L)\), so any fixed exponential covering weight contributes at most \(\exp(C_L(1+h))\). If the coupling at the start of the stage is sufficiently large, multiplication by \((4\epsilon)^h\) still gives fixed exponential decay, part of which can be reserved for weighting the degree of the matrix product. Spatial growth of fixed-order directional derivatives contributes at most another factor \((1+hB')^C\). Summing over terms whose supports contain a given site introduces only one polynomial factor for possible roots and intracell coordinates. The traces of matrix powers in \eqref{eq:geom-determinant} are handled in the same way, by extracting coefficients at each root and then taking finite differences in the support. The estimate uses contraction of the matrix products in the energy norm.

When the neighborhood of the root satisfies the small-field condition, a local spanning-tree gauge bounds the link coordinates by \(t_f\) times a fixed polylogarithmic factor. The bounded-degree part of the determinant expansion differs from its flat value by \(\widehat\beta_f^{-1/2+o(1)}\). When its dependence set meets a selected test, use the original integral on the test region. Each test acquires at most a fixed polylogarithmic number of additional coordinates, and the absolute logarithm of its extra factor has a fixed polylogarithmic bound. The energy on the selected region has coefficient \(\beta'\), of which the original test suppression bound uses only \(\widehat\beta_{f+1}/2\). If \(\beta'\ge(1/2+\delta)\widehat\beta_{f+1}\) for fixed \(\delta>0\), the difference supplies additional exponential suppression per test. In a fixed stage the added exponent can be taken to be a fixed fraction of \(c_\delta M_{f+1}^{2}\), which exceeds the logarithmic bound on the bounded-degree term. The weaker margin \(2\delta\log(1024/\epsilon_f^{\rm for})\) required by the original connected grouping also suffices, where \(\epsilon_f^{\rm for}=\exp(-M_f^{1/5})\) is the weight bound for terms with triggered tests. Frozen dimensions, local Haar masses, and positive normalization factors are all computed in the same original integral.

The bounded-degree terms that depend on the test selection can be incorporated into the original exact grouping by triggered tests. Fix a compatible test selection \(\mathfrak a\), and let \(g_\gamma(U,\mathfrak a)\) be such a correction, with full dependence set \(Z_\gamma\). It depends only on tests meeting \(Z_\gamma\) and vanishes if there is no such test. With \(\mathfrak a\) fixed, expand \(\prod_\gamma[1+(e^{g_\gamma}-1)]\). Group all retained terms and original tests into connected components according to intersections of their dependence sets. Each retained term is connected to at least one original test and depends only on tests in its component. A triggered test on a component's support must be assigned to an original factor in that component, since otherwise the dependence sets of two components would intersect. The record of triggered tests in the new component is thus the union of the original records. Conversely, each admissible family of new components determines uniquely the original selection and the retained terms. This gives an identity at finite truncation, with the original Borel indicators kept fixed as integration conditions. The first grouping occurs inside the same Gaussian and Haar integrals. The second takes the logarithm of the remaining hard-core polymer expansion on compatible supports. The two absolute tree sums use the respective rooted summation margins in \cite[Lemmas 13.4--13.5 and 13.11]{armstrong}. The support \(X\) offers at most \((2B'+1)^4\operatorname{cov}_{B'}(X)\) possible intersection sites. Thus passing from a weighted sum at a fixed site to the sum over possible rooted branches costs this local volume factor once.

This representation uses directional derivatives through order six and spatial weights of order two. The regular activities have weighted norm at most \(\omega_f=\widehat\beta_f^{-1/512}\). Terms in the original one-step integral that contain no interaction from the preceding finer level satisfy the bound \(\widehat\beta_f^{-7/16}\). The bounded-degree part, the long-range remainder, and the window boundary correction of the new small-field Schur terms each satisfy a stronger inverse-power estimate. Multiplying the weights of small regular activities by \(\widehat\beta_f^{1/4096}\) preserves the tree-summation conditions. Removing this multiplier from graphs containing two or more small terms gains an additional inverse power. When there is a single interaction from the finer level, its leading average is handled by the radial expansion of \cite[Lemma 13.6]{armstrong}. After subtraction of the constant and the full affine Hessian, the first remaining term carries five inverse powers of the blocking ratio. Multiplication by the four-dimensional cell volume bounds its coefficient by \(V_j/\ell_f\) plus the original comparison remainder, where \(V_j\) depends only on the stage. The ratio of \(\omega_f\) on adjacent levels tends to one. The original smallness conditions at the start of the stage therefore propagate this norm bound. We call this the extended-support representation. It is used to compare the two geometries. The total density on the actual square torus can also be calculated from the original representation, whose dependence ranges are smaller and whose remainder bounds are stronger.

After calibration of the coupling, terms with triggered tests are merged into connected components. Let \(e^{-2K_f}\) bound the original weighted sum at any site for these terms after integration and before merging. Set \(\mathcal Q_f=(2B_f^{\mathrm{cov}}+1)^4\), and let \(\gamma_{\log}\) denote the weighted sum at a site of the new regular logarithmic terms, so that \(\mathcal Q_f\gamma_{\log}=o(1)\). Extract \(2\mathcal Q_f e^{-2K_f}\) from the weights of terms with tests. The remaining sum at each site is at most \(1/(2\mathcal Q_f)\). The tree estimate of Lemma 13.11 gives
\[
 2\mathcal Q_f e^{-2K_f}
    \left(\exp\{1/(2\mathcal Q_f)+\gamma_{\log}\}-1\right)
 = (1+o(1))e^{-2K_f}.
\]
The original admission conditions \(e^{-K_f}(e^{\gamma_f^{\rm old}}-1)\le\epsilon_f^{\rm for}/1024\) and \(\gamma_f^{\rm old}\ge e^{-K_f}\) give \(e^{-2K_f}\le\epsilon_f^{\rm for}/1024\), where \(\gamma_f^{\rm old}\) is the tree-summation parameter of the original polymer expansion. The original bound for terms with tests is therefore preserved after extending their support. The estimates concern absolute values of the signed expansion terms.

\subsection{Two independent truncations and the one-step comparison}

At level \(f\), set \(r=n_f^{1/3}\). For each term in the one-step expansion, record its maximum distance from the root to the full dependence set, and separately a weighted degree \(\mathcal N\). Each occurrence of a regular factor and each triggered test counts once. A determinant term is counted by its degree in local matrix products and a path by its length divided by \(B'\). An interaction from the previous finer level counts as one factor at this step, with its entire dependence set still included in the support. Repeated use of a local variable counts with multiplicity. Reserve fixed fractions of the path decay, \eqref{eq:geom-cauchy}, the extra weights on small regular activities, and the energy suppression for tests. In the weaker covering norm the two tree sums then allow the factor \(e^{\lambda\mathcal N}\), with \(\lambda>0\) fixed before the stage starts. The resulting absolute sum at each site still satisfies a stage-dependent polynomial bound. Terms with dependence radius greater than \(r\) are controlled by the covering weight. For terms with small dependence radius but \(\mathcal N>r/B'\), use the separate bound \(e^{-\lambda r/B'}\). This controls, in particular, path products that repeatedly backtrack inside a small window.

Every remaining term has the same lift to the integer lattice in both geometries and consists of at most \(O(r)\) elementary steps. Compare the periodization of the underlying matrix before applying local matrix operations. For common indices at distance at most \(Cr\), the unmatched periodic images are shifted by at least \(cn_f\). The resulting entrywise difference is at most
\begin{equation}
 \Pi_j(1+n_f+\widehat\beta_f)\exp(-c_j n_f/B').
 \label{eq:geom-periodization}
\end{equation}
This estimate also applies when distant entries are truncated after periodization. In direction \(a\), a period of \(M_a\) lattice steps gives Fourier phases \(2\pi u_a/M_a\) for integers \(u_a\). The coarse phase \(p_a\) corresponds to fine phases \((p_a+2\pi z_a)/\ell_f\), where \(z_a\) ranges over the residue classes modulo \(\ell_f\). Extending the time direction preserves the shared local terms, with the same pointwise bound on the sum over periodic images. Let \(A_1,A_2\) be the positive principal submatrices being compared. For \(A_s=(1-s)A_2+sA_1\), use
\[
 A_1^{-1}-A_2^{-1}=A_1^{-1}(A_2-A_1)A_2^{-1},\qquad
 \log\det A_1-\log\det A_2=\int_0^1\Tr\big(A_s^{-1}(A_1-A_2)\big)\,ds
\]
to compare inverses and determinants. Differences of affine means, Gaussian integrals in the original coordinates, and finite products are taken one factor at a time. For a given test, compare the positive densities with different means and covariances on the same fixed Borel domain in physical coordinates. The product of the other factors in a fixed stage is bounded crudely by \(\exp(C_jr(1+\log(n_f+\widehat\beta_f))^{C_j})\), which is absorbed by \eqref{eq:geom-periodization}. Combining the two truncations, if the interactions on the adjacent finer level agree, the one-step shape discrepancy satisfies
\begin{equation}
 \varepsilon_f^{\mathrm{geom}}\le \Pi_j(1+n_f+\widehat\beta_f)e^{-c_jn_f^{1/3}/B'} .
 \label{eq:geom-seed}
\end{equation}
On the rectangle we use the local paths generated by lattice coordinate rotations on the square torus, with their original weights. The local lift preserves these paths. The full rectangle only needs the reflection symmetries used in the comparison.

Let \(d_{f+1}\) denote the supremum over roots of the covering-weighted sum of the absolute differences of local terms evaluated on the shared variables, as specified before \eqref{eq:paper-shape-recurrence}. For regular exponential factors use
\[
 e^{-E}-e^{-E'}=(E'-E)\int_0^1 e^{-sE-(1-s)E'}\,ds.
\]
Products containing terms with triggered tests are telescoped one factor at a time. In both tree sums, each term in the difference contains just one factor to which the difference is applied. Fix the fine links on which this factor depends. Derivatives in a coarse-field direction then act on the shared conditional density and on the other smooth terms. The weighted value difference at a site, before local merging, is bounded by \(U\le \Pi_j(1+n_f+\widehat\beta_f)(d_{f+1}+\varepsilon_f^{\mathrm{geom}})\). In the same real affine direction, extract the second-order coefficient on the small-field interval \([-h_0,h_0]\) by a centered difference of step size \(0<h\le h_0\). If \(u_2\) bounds the difference in values and \(v_4\) bounds the difference in fourth derivatives, the difference of second derivatives is bounded by \(4u_2/h^2+v_4h^2/12\). One can use the same relative step size \(h/h_0\) in the weighted sum, or optimize at each root and sum using Cauchy--Schwarz. This gives
\begin{equation}
 C h_0^{-2}(U+\sqrt{UV}),
 \label{eq:geom-hessian}
\end{equation}
where \(V\) is the corresponding weighted sum of fourth-derivative bounds in the normalized direction. Affine link amplitudes on the shared dependence set are \(O(1+r)\), so \(h_0\asymp t_f/(1+r)\) is admissible. Polarization gives mixed coefficients from finitely many real directions. The extracted quadratic terms are subtracted from the local functions. We then adjust the real path-action coefficients and merge the remaining value differences. The procedure uses square-root interpolation once and gives
\begin{equation}
 d_f+\max_{a<c}|d_{ac,f}|+|\delta c_f^{\mathrm{norm}}|
 \le \Pi_j(1+n_f+\widehat\beta_f)(d_{f+1}+\varepsilon_f^{\mathrm{geom}})^{1/2}.
 \label{eq:geom-one-step}
\end{equation}
The indices \(a,c\) run over distinct coordinate directions, and \(d_{ac,f}\) is the difference of the corresponding affine-curvature coefficients on the rectangle and the square torus. The difference of the per-cell field-independent logarithmic normalization increments on two square tori at the same bare coupling is \(\delta c_f^{\mathrm{norm}}\). It is evaluated at the identity field and includes the Gaussian normalization with no test selected and the remaining regular terms. The periodic images of the Gaussian term for the empty test selection also satisfy \eqref{eq:geom-periodization}. Field-independent normalization factors cancel in the relative density on the rectangle. The output supports shrink on the coarse lattice. Differences use the weaker covering weight, and new factors use the stronger weight already established.

\subsection{Fourier calibration of the coupling on a fixed square torus}
\label{lem:geometry-shoot}

Fix a square torus and a blocking level. Let \(\mathsf K_f\) be the array of path coefficients of the quadratic curvature action and let \(\mathcal L_f\) be the closed-loop action inherited from the finer level. Let \(\rho\) be an admissible representation of the positive density at that level, with \(\beta_f\) the coefficient multiplying \(\mathsf K_f\). Fix two axes \(a\ne c\) and a unit color vector \(S\in\mathfrak{su}(2)\). Use \(e^{ipy_a}S\) as the complex test direction on links parallel to axis \(c\), where \(y_a\) is lattice coordinate \(a\) of the base point of the link. Set the test direction to zero on the other links. Use polarization of the real and imaginary parts to interpret the complex Hessian. Write \(n=n_f\) in this paragraph. For every discrete phase \(p=2\pi j/n\), let \(h_\rho(p)\) be the per-cell Hessian of \(-\log\rho-\mathcal L_f\) evaluated on this direction and its complex conjugate. The same test applied to \(\mathsf K_f\) defines \(h_K(p)\). Set
\[
 \mathfrak D_n h=-\sum_{r\in I_n}r^2\widehat h_n(r),\qquad
 \widehat h_n(r)=\frac1n\sum_{j=0}^{n-1}
                    e^{-2\pi i jr/n}h(2\pi j/n),
\]
where \(I_n\) is a centered system of residue classes modulo \(n\). If a regular activity lifts into a lattice ball of radius \(n/10\), the inverse Fourier transform has no aliasing and \(\mathfrak D_n\) gives the second derivative of the resulting trigonometric polynomial at zero. With \(x\) the lattice coordinate in direction \(a\) relative to the root and \(H\) the activity's Hessian, that derivative is \(2H(xS,xS)-H(S,x^2S)-H(x^2S,S)\). The first term vanishes because the affine Hessian has already been subtracted. The constant direction \(S\) is a gradient on the local lift, so the last two terms vanish by local gauge invariance and the vanishing of the first derivative at the identity field. For the remaining supports, the test at each phase is bounded by \(Ct_f^{-2}\) times the regular norm of the activity. A support beyond the local lifting range requires at least \(cn_f/(2B_f^{\mathrm{cov}}+1)\) covering balls. Also \(\|\mathfrak D_n\|_{\ell^\infty\to\mathbb R}\le Cn^3\). Thus, provided \(|\mathfrak D_n h_K|\ge1\),
\begin{equation}
 \left|\frac{\mathfrak D_n h_\rho}{\mathfrak D_n h_K}-\beta_f\right|
 \le Cn_f^3 t_f^{-2}\,\omega_f\exp[-c n_f/(2B_f^{\mathrm{cov}}+1)] .
 \label{eq:geom-beta-readout}
\end{equation}
The symbol of the quadratic curvature action in this direction is \(4\sin^2(p/2)\widehat{\mathsf H}_{ac,ac}(pe_a)\), where \(e_a\) is the coordinate unit vector and \(\widehat{\mathsf H}\) is the Fourier symbol of the flat quadratic curvature kernel. We have \(\widehat{\mathsf H}(0)=I\) and exponential decay of its Fourier coefficients \cite[Lemmas 12.3--12.4]{armstrong}. The second-moment aliasing from coefficients outside the centered representatives modulo \(n\) is bounded by
\(C(1+(A_j^{\rm path})^2)(n_f^2+(B_f^{\mathrm{cov}})^2)e^{-c n_f/B_f^{\mathrm{cov}}}\),
where \(A_j^{\rm path}\) is the weighted absolute sum of the flat-kernel path coefficients at stage \(j\). The denominator therefore tends to \(2\) on the chosen cutoff tail. Taking the second moment over all discrete phases removes the contribution of each short-support regular activity exactly.

This calibration compares the extended-support representation on the same square torus with the actual probability law obtained by tuning to the prescribed terminal coupling. Both joint laws start with the same bare Wilson field and use the same random tags and blocking rules. Their coupling coefficients may differ at each step. Set
\(\delta_f=|\beta_f^{(1)}-\beta_f^{(2)}|\) and
\(w_f=\sum_{i=f}^{N-1}n_i^4\delta_{i+1}\).
By \(\partial_\gamma\log q_\gamma(W)=-\operatorname{pr}W+\mathbb E_{q_\gamma}\operatorname{pr}\) and the count \(4n_i^4\) of coarse links at level \(i\), the log likelihood ratio of the two joint densities along any complete field history has absolute value at most \(Cw_f\). Given the same coarse field \(V\), the total variation distance between the two conditional laws of the finer levels is at most \(C\min(1,w_f)\). Let \(S_1,S_2\) be the first derivatives of the logarithm of the last conditional blocking density in the two test directions, and \(S_{12}\) its mixed second derivative. Write \(p_f\) for the actual marginal density of the coarse field after integrating out the fine field. With \(D_{12}\) the mixed derivative along these directions, at the identity coarse field we have
\[
 D_{12}\log p_f=\mathbb E(S_{12}\mid V)
                  +\operatorname{Cov}(S_1,S_2\mid V).
\]
For directions of amplitude at most one, the absolute values of these scores are bounded by \(C n_f^4\beta_{f+1}\). For the same fine field and tag the score differences are at most \(C n_f^4\delta_{f+1}\). Dividing by the number of cells and applying the conditional-law estimate gives
\[
 \sup_p|h_{\rho^{(1)}}(p)-h_{\rho^{(2)}}(p)|
 \le C n_f^4(1+\widehat\beta_{f+1})^2(w_f+\delta_{f+1}).
\]
The quadratic curvature path arrays and the inherited loop actions agree for the two representations. Let \(\mathfrak e_f\) be the right-hand side of \eqref{eq:geom-beta-readout}. We obtain
\begin{equation}
 \delta_f\le2\mathfrak e_f+C n_f^7(1+\widehat\beta_{f+1})^2 w_f,\qquad
 w_{f-1}=w_f+n_{f-1}^4\delta_f,\qquad w_N=0 .
 \label{eq:geom-likelihood-recurrence}
\end{equation}
At the first step at which a coupling could leave the admissible real interval, the input couplings still have a fixed relative margin from the interval boundary. The Hessian calibration can therefore be applied before merging the local terms. From \(n_{f+1}\ge3n_f\) and the original stage conditions,
\(\prod_{i\le f}(1+C n_i^{11}(1+\widehat\beta_{i+1})^2)\)
is at most \(\exp(C(1+\log n_f)^2)\). At levels \(n_f\ge b^\alpha\) for fixed \(\alpha>0\), we have \(\log\widehat\beta_{f+1}=O(1+\log n_f)\). Thus \(\exp[-c n_f/(2B_f^{\mathrm{cov}}+1)]\) absorbs this product. Iterating \eqref{eq:geom-likelihood-recurrence} from the same bare coupling gives \(\delta_f+w_f\le C e^{-c\sqrt{n_f}}\). For sufficiently large \(b\), this difference is smaller than the margin from the original square-torus coupling to the boundary of the admissible interval, ruling out a first exit. To compare absolute pressures at fixed \(b\), one may instead stop at a shared, sufficiently fine fixed level. This conditional likelihood comparison takes place on the joint configuration space of a single square torus.

\subsection{Six curvature coefficients and their effect on the law}

The rectangle has reflection symmetry separately in each coordinate direction. The product of affine curvatures in two distinct coordinate planes changes sign under at least one of these reflections, so the cross terms of the affine Hessian vanish. The rectangle has six diagonal coefficients \(\Delta_{ac}^{\mathrm{rect}}\), and the square torus has one common coefficient \(\Delta^{\mathrm{box}}\) by coordinate-permutation symmetry. First subtract the full affine Hessian of each locally liftable term. Then write
\begin{equation}
 \Delta_{ac}^{\mathrm{rect}}=\Delta^{\mathrm{box}}+d_{ac,f},\qquad
 \mathcal D_f(U)=\tfrac12\sum_{a<c}d_{ac,f}
       \sum_o h_{\{o\}}(U)\,|F^{\rm cl}_{ac}(U,o)|^2 .
 \label{eq:geom-reset}
\end{equation}
Here \(o\) runs over lattice sites, \(h_{\{o\}}\) is the original small-field indicator at \(o\), and \(F^{\rm cl}_{ac}\) averages the Lie-algebra projections of the four positively oriented plaquette loops in the specified coordinate plane based at \(o\). The projection for a loop is \(\operatorname{proj}_{\mathfrak g}(U_\square-U_\square^{-1})/2\), where \(\mathfrak g=\mathfrak{su}(2)\) and \(U_\square\) is the ordered product of the plaquette links. Thus the first-order term of the clover field strength in an affine test field is the curvature in the corresponding plane. Use \(\Delta^{\mathrm{box}}\) for the original path replacement, and merge terms with triggered tests according to the original support rules \cite[Lemmas 13.10--13.13]{armstrong}. If the support indicator for several sites differs from the single-site indicator, keep the original admissibility conditions and flag records and include the discrepancy in the same merge. The remaining differences in all six directions appear in \(\mathcal D_f(U)\), with the same formula on configurations with triggered tests. Let \(\mathcal H_f^{\mathrm{rect}}\) be the actual positive pushforward density on the rectangle at that level and \(\rho_f^{\mathrm{rect}}\) the comparison density after removing the field-independent factor. On the full configuration space,
\[
 \mathcal H_f^{\mathrm{rect}}(U)/\mathcal H_f^{\mathrm{rect}}(\mathbf 1)
   =e^{-\mathcal D_f(U)}\rho_f^{\mathrm{rect}}(U).
\]
The output regular activities have had their full quadratic terms in the six directions subtracted, so the next blocking step again uses interactions with zero affine Hessian. The square-torus coefficient used in the comparison has already been calibrated from the actual field density before grouping into connected components.

The positive reweighting at each level can be written as one likelihood ratio on the full history. Let \(\mathbb P\) be the original joint law of the fields \(X_f\) at every level, the random blocking tags, and the designated end-face variables. Let \(\widetilde{\mathbb P}\) be the joint law obtained by multiplying at each level by \(e^{\mathcal D_f(X_f)}\) and renormalizing. Substitution of the normalization factors at each step gives
\begin{equation}
 \frac{d\widetilde{\mathbb P}}{d\mathbb P}
 =\frac{e^{\sum_{f\ge m}\mathcal D_f(X_f)}}
        {\mathbb E_{\mathbb P}e^{\sum_{f\ge m}\mathcal D_f(X_f)}} .
 \label{eq:geom-history-resets}
\end{equation}
The sum extends to the cutoff level \(N\), and \(m\) is the coarser stopping level in \eqref{eq:paper-stops}. Write \(\operatorname{osc}F=\sup F-\inf F\) and \(B_{\rm tot}=\sum_{f\ge m}\operatorname{osc}\mathcal D_f\). The likelihood ratio lies between \(e^{-B_{\rm tot}}\) and \(e^{B_{\rm tot}}\). The marginal of any observable retains these bounds. Upon conditioning on the same record, the likelihood ratio lies between \(e^{-2B_{\rm tot}}\) and \(e^{2B_{\rm tot}}\).

To apply \eqref{eq:geom-one-step}, choose the odd blocking ratio \(\ell^{\mathrm{st}}>4096\) at the initial stage of the comparison, so that \(\tfrac12\ell_f^{1/4}>4\). The second-order extraction fixes the exponent \(1/2\) before the blocking ratio is chosen. The procedure in \cite[Table 3]{armstrong} fixes the stage constants and blocking ratios before increasing the starting and terminal couplings. In the finite prefix before the chosen stage, each comparison interval has length at most \(O_j((1+\log\log b)^2)\). The cumulative logarithmic scale before the chosen stage therefore satisfies \(s_f=o(\log b)\). Extend the chosen stage to obtain stopping levels \(m,m_b\) with \(PD_m\asymp b^{1/24}\) and \(PD_{m_b}\asymp b^{1/12}\). In a fixed stage the admissible period must control \((1+B_f^{\mathrm{cov}}+B_{f+1}^{\mathrm{cov}})^{C_j}e^{-s_f}\). Since \(B_f^{\mathrm{cov}}=\lceil(\log(b+\kappa_Gs_f))^2\rceil\), one can bound these quantities over the finite prefix by an integer \(M_{\rm per}\) that is at most a fixed power of \(1+\log b\). Define the three level-zero periods by
\[
 p_-=4(M_{\rm per}+1),\qquad P=8p_-,\qquad p_+=2P .
\]
For each subsequent stage, first choose both available blocking ratios to be odd integers satisfying \(\tfrac12\ell_f^{1/4}>4\). Then delay the starting level of each stage so that its local constants are bounded by positive powers of \(n_f\). These choices preserve the asymptotic scale conditions of the construction in \cite[Lemmas 14.2 and 14.4]{armstrong}.

For the levels \(n_f\ge n_m\), the fixed-stage polynomial satisfies \(\log\Pi_j(1+n_f+\widehat\beta_f)=O_j(\log n_f)\). By \eqref{eq:geom-seed}, the one-step shape discrepancy can be made smaller than \(e^{-8c_*n_f^{1/4}}\). Since \(n_{f+1}=\ell_f n_f\), if \(d_{f+1}\le e^{-c_*n_{f+1}^{1/4}}\), taking the square root leaves a decay exponent at least \(4c_*n_f^{1/4}\). The exponent \(c_*n_f^{1/4}\) remains after accounting for the polynomial and the sum. Delaying the starting levels of later stages preserves the same \(c_*>0\). At level \(N\) the two geometries have the same local bare Wilson action, so \(d_N=0\). Thus
\begin{equation}
 d_f+\max_{a<c}|d_{ac,f}|+|\delta c_f^{\mathrm{norm}}|
       \le C e^{-c_*n_f^{1/4}}\qquad(f\ge m).
 \label{eq:geom-uniform}
\end{equation}
Tuning the bare coupling to the prescribed terminal value \(v\) and summing the one-step scalar drifts give
\[
 \big|\beta_{N,P,f}(v)-v-\kappa_G\log D_f\big|
       \le\tfrac14\kappa_G\log D_f,\qquad
 \beta_{N,P,m}(v)=v+O(\log b).
\]
At each step the coupling has a fixed relative margin from the boundary of the real admissibility interval, uniformly for \(|v-b|\le c_0 b^{31/32}\). To obtain convergence at a fixed level uniformly over the band, the result of \cite[Lemma 13.24]{armstrong} gives \(X_N(v_N)\to X_\infty(v)\) whenever \(v_N\to v\), where \(X_N(v)\) denotes the object being compared at the fixed level. If \(X_\infty\) were discontinuous on the closed band, one could take \(v_j\to v\) and then choose \(N_j\) pointwise sufficiently large that \(X_{N_j}(v_j)\) is close to \(X_\infty(v_j)\), contradicting the stated sequential convergence. Once the limit is continuous, failure of uniform convergence would likewise furnish a sequence \(v_j\) on which the error stays bounded below. Compactness of the band gives a convergent subsequence and the same contradiction. In this argument \(b,P,f\) are fixed. The level-dependent quantitative comparison uses \eqref{eq:geom-uniform}.

Since \(SU(2)\) is compact, the clover field strength is bounded on the full configuration space. For time period \(kL_0\) there are \(kn_f^4\) cells, so \(\operatorname{osc}\mathcal D_f\) is at most \(Ckn_f^4\max|d_{ac,f}|\). Summing \eqref{eq:geom-uniform} over scales yields
\begin{equation}
 B_{\rm tot}\le Ck e^{-c n_m^{1/4}} .
 \label{eq:geom-total-reset}
\end{equation}
This likelihood bound holds for all records of triggered tests and random tags.

\subsection{Absolute pressures and the short-time heat trace}
\label{lem:geometry-pressure}

At the common grid spacing \(a_N=L_0/(P D_N)\), the physical side lengths of the two comparison square tori are \(L_-=L_0p_-/P\) and \(L_+=L_0p_+/P\). Both use the bare coupling \(\beta_{N,P}(v)=\beta_{N,P,N}(v)\) selected on the reference square torus. Retain the field-independent logarithmic normalizations in the comparison. For fixed \(b\), one may stop at any shared, sufficiently fine fixed level \(j_0\). The field-independent increments at the flat configuration follow from \eqref{eq:geom-affine}, \eqref{eq:geom-determinant}, and the four Haar normalization denominators per cell. The identity coarse field has no triggered tests on the coarse lattice, and the affine source vanishes when no test is selected. If a test is selected, it can only come from an actual fine plaquette or a coarse-link mismatch. Changes in the number of frozen coordinates, the energy of the affine mean, and determinant differences remain with the connected components that trigger them, and are controlled by the same original exponential test bound. Thus in comparing absolute increments each contribution from the flat windows and each regular logarithmic term at a root is counted once. For the same bare input, the local coefficients for the empty test selection agree per cell across square tori. Differences from the remaining periodic images are controlled by \eqref{eq:geom-seed}, and the actual per-cell differences by \eqref{eq:geom-uniform}. Per cell at the original cutoff, the contribution of level \(f\) carries the volume factor \((D_f/D_N)^4\). The cumulative difference is bounded by
\[
 D_N^{-4}\,C\sum_{f\ge j_0}D_f^4e^{-c n_f^{1/4}}
             \le C_{b,p_-,p_+}D_N^{-4}.
\]
On the same square torus, \eqref{eq:geom-likelihood-recurrence} controls the difference between the joint law used in the comparison and the actual joint law at the prescribed bare coupling. The normalized conditional blocking kernel preserves the exact partition function. At the stopping level, absolute sitewise sums of local actions and grouping of tests give a finite upper bound on the total contribution. Restrict the links to Haar balls of chord radius \(c\widehat\beta_{j_0}^{-1/2}\). No test is then triggered, the action is at most a constant times the number of lattice sites, and the logarithm of the Haar volume is at least \(-C n_{j_0}^4(1+\log\widehat\beta_{j_0})\). Since \(b,j_0\) are fixed, the difference in the remaining logarithmic normalizations at the stopping level, per cell at the original cutoff, is at most \(C_{b,v}a_N^4\).

Let \(Z_{a,\beta}(\boldsymbol L)\) be the unnormalized Wilson partition function on the four-dimensional Euclidean torus of side lengths \(\boldsymbol L=(L_1,L_2,L_3,L_4)\), lattice spacing \(a\), and bare coupling \(\beta\). The pressure per lattice cell is \(\mathfrak p_a(\boldsymbol L)=\log Z_{a,\beta}(\boldsymbol L)/\prod_i(L_i/a)\). Along any coordinate axis, the partition function is the trace of an integer power of the corresponding positive transfer operator. The Schatten norm of a positive operator is nonincreasing in its exponent, so \(\mathfrak p_a\) is nonincreasing in each period. Write \(\boldsymbol L_\pm=(L_\pm,L_\pm,L_\pm,L_\pm)\). The comparison of normalizations at the same bare coupling gives
\begin{equation}
 0\le\mathfrak p_a(\boldsymbol L_-)-\mathfrak p_a(\boldsymbol L_+)\le C_{b,v}a^4 .
 \label{eq:geom-pressure}
\end{equation}
In the common physical units, the two tori with spatial side \(L_0\) and time periods \(s_*=L_0/4\) and \(2s_*\) have all their side lengths between those of the comparison tori. Let \(Z_N(s)\) denote the unshifted heat trace of period \(s\) in this spatial box, and \(S_N(s)\) the vacuum-normalized semigroup. For \(r=s_*/a_N\), the inequality \(\Tr K_N^{2r}\le\|K_N\|^r\Tr K_N^r\) gives the first bound below, and the pressure difference at the same bare coupling gives the second.
\[
 \Tr S_N(s_*)\le
    \frac{Z_N(s_*)^2}{Z_N(2s_*)}
 \le \exp(2L_0^3s_* C_{b,v}) .
\]
For fixed \(b,v\), this constant is uniform in the cutoff. The short-trace estimate with explicit dependence on \(b\) in the main text instead uses a conditional comparison in neighborhoods of the end faces.

\subsection{Time folding and the small-field density on the long cylinder}
\label{lem:geometry-fold}
\label{prop:geometry-density}

At stopping level \(m\), the spatial side has \(n=n_m\) lattice steps and the time period has \(q=kn\) lattice steps, where the integer \(k\) is set by \eqref{eq:paper-longtime}. Truncate the regular activities in their dependence radius and weighted degree. The extended-support representation of the square torus supplies a common family of local functions \(f_{x,t}\) that can be evaluated in either geometry. Here \(x,t\) are the spatial and temporal coordinates of the root. Subtract the value of each local function at the identity field. The total action is invariant under time reflection, so the rooted functions may be averaged with their reflections without changing the total on either torus. The time dependence of each function lies in \([t-r_{\mathrm t},t+r_{\mathrm t}]\) for an integer \(r_{\mathrm t}\le Cn^{1/3}\). With \(\Lambda_n=(\mathbb Z/n\mathbb Z)^4\), define
\[
 f_t(U)=\sum_x f_{x,t}(U),\qquad
 \bar f(U)=q^{-1}\sum_t f_t(U),\qquad
 E_n(V)=\sum_{(x,t)\in\Lambda_n}f_{x,t}(V).
\]
By \eqref{eq:geom-seed} and \eqref{eq:geom-uniform}, the discarded or unmatched terms contribute at most \(Cn^4e^{-c n^{1/4}}\) and \(Ckn^4e^{-c n^{1/4}}\) to the total actions on the two geometries, respectively, on configurations free of large-field flags. The estimate uses weighted sums at a site for the shared local terms in the extended-support representation. Equation \eqref{eq:geom-likelihood-recurrence} compares the extended-support and original representations at the same bare coupling on the same square torus. On the small-field set the log likelihood ratio has absolute value at most \(Ce^{-c n^{1/4}}\). The absolute sums of the original square-torus representation give
\begin{equation}
 \sup_{V\ \text{free of large-field flags}}|E_n(V)|
 \le Cn^4(\widehat\beta_m^{-11/10}+e^{-c n^{1/4}}) .
 \label{eq:geom-folded-box}
\end{equation}

For a starting time \(a\in\mathbb Z/q\mathbb Z\) and length \(\ell\in\{\lfloor n/8\rfloor,\lfloor n/4\rfloor\}\), define the folded square-torus field \(V_{a,\ell}\) of period \(n\). It remains at \(a\) for \(\lfloor n/8\rfloor\) steps, traverses the original time segment \([a,a+\ell]\) forward, remains at the endpoint for \(n-2\ell-\lfloor n/8\rfloor\) steps, and traverses the same segment backward. The spatial links are taken from the original field at the corresponding times. The temporal links are the original links on the forward segment, their inverses on the backward segment, and identities on the stationary segments. If the field \(U\) on the long cylinder is free of large-field flags, each spatial plaquette of the folded torus comes from an original spatial plaquette. Each spacetime plaquette equals an original spacetime plaquette, its inverse, a conjugate, or the identity. Chord lengths do not increase, so \(V_{a,\ell}\) is also free of flags. Write \(V_{t,*}\) for the square-torus field held fixed at the original time \(t\) throughout, and set
\[
 \bar f^*=q^{-1}\sum_t E_n(V_{t,*})/n,\qquad
 A_\ell=q^{-1}\sum_a E_n(V_{a,\ell}).
\]
All four segments are longer than \(2r_{\mathrm t}+1\). After averaging over \(a\), roots in the traversed segments at distance more than \(r_{\mathrm t}\) from a junction contribute \(2(\ell-2r_{\mathrm t}-1)\bar f\). Roots in the corresponding interiors of the stationary segments contribute \((n-2\ell-4r_{\mathrm t}-2)\bar f^*\). The \(2r_{\mathrm t}+1\) time roots near a junction depend only on the stationary endpoint and the \(2r_{\mathrm t}\) steps on the adjacent traversed segment. Both endpoints range over the full time circle as \(a\) varies, so the sum of the averaged contributions from the four junctions is independent of \(\ell\). Subtracting the equations for the two folding lengths gives
\begin{equation}
 \bar f=\bar f^*+
   \frac{A_{\lfloor n/4\rfloor}-A_{\lfloor n/8\rfloor}}
        {2(\lfloor n/4\rfloor-\lfloor n/8\rfloor)} .
 \label{eq:geom-folding-identity}
\end{equation}
If \(n\ge64(r_{\mathrm t}+1)\) and \(|E_n|\le n^4\epsilon\) on the square torus, then \(|\bar f^*|\le n^3\epsilon\). The two integer lengths differ by at least \(n/16\), so the difference quotient is bounded by \(16n^3\epsilon\). Multiplication by the number \(q\) of time steps gives \(17kn^4\epsilon\).

Let \(\mathcal G_{\mathrm{cyl}}\) be the set of configurations of the actual coarse long cylinder free of large-field flags. Let \(\mathbb Q\) denote the probability obtained by evaluating the short torus's common scalar coefficient and analytic path action on the long cylinder using the same lifted paths, then restricting to \(\mathcal G_{\mathrm{cyl}}\) and normalizing. Let \(E_{\mathrm{cyl}}\) be the regular remainder of the action of the actual long cylinder, with its value at the identity subtracted. Equations \eqref{eq:geom-folded-box} and \eqref{eq:geom-folding-identity} give
\begin{equation}
 \sup_{\mathcal G_{\mathrm{cyl}}}|E_{\mathrm{cyl}}|
 \le Ckn^4(\widehat\beta_m^{-11/10}+e^{-c n^{1/4}})
                  =O(b^{-3/5+o(1)}),
 \label{eq:geom-folded-cylinder}
\end{equation}
where \(k\asymp b^{1/3}\) and \(n\asymp b^{1/24}\). In the absence of large-field flags, the expansion factor for triggered tests is identically one. The common field-independent factor cancels. The actual conditional density on this set relative to \(\mathbb Q\) is \(e^{-E_{\mathrm{cyl}}}/\mathbb E_{\mathbb Q}e^{-E_{\mathrm{cyl}}}\). If \(\operatorname{osc}E_{\mathrm{cyl}}\le\zeta\), the total variation distance is at most \((e^\zeta-1)/2\). To restore the actual weights in the six coordinate planes, use \eqref{eq:geom-total-reset}. The reflection-positive chessboard estimate before \eqref{eq:paper-backward} bounds the probability of a flag at the stopping level of the actual long cylinder by \(O_A(b^{-A})\). The total variation distance of the actual marginal from \(\mathbb Q\) is therefore \(O(b^{-3/5+o(1)})\), and the Hellinger distance with our convention is \(O(b^{-3/10+o(1)})\). Applying a common conditional kernel to adjoin observables on both sides preserves this bound.

\section{Boundary Decoupling and Observable Kernels on the Cylinder}
\label{cam-bound-section}

We use the reference period \(P\), cumulative blocking factor \(D_f\), and cutoff level \(N\) from Section~\ref{sec:main}. At level \(f\), the spatial period has \(n_f=PD_f\) lattice steps and the lattice spacing is \(a_f=L_0/n_f\). Write \(\lambda_f=\widehat\beta_f\) for the reference blocking coupling, \(\beta_f=\beta_{N,P,f}(v)\) for the actual coupling, and \(B_f^{\rm cov}\) for the lattice ball radius used to define support weights. The time period is \(T=kL_0\) with \(k\) a positive integer. We take \(k=1\) on the short square torus and \(k\asymp b^{1/3}\) on the weak-coupling long cylinder. Fix a sufficiently large \(A_0\) and define the auxiliary quantity \(M_f\) and the curvature thresholds \(t_f,\tau_f\) by
\[
 \begin{gathered}
 M_f=\exp((\log\log\lambda_f)^2),\qquad B_f^{\rm cov}\asymp(1+\log\lambda_f)^2,\\
 t_f=M_f/\sqrt{\lambda_f},\qquad \tau_f=t_f(1+\log\lambda_f)^{-A_0}.
 \end{gathered}
\]
The fields, predictions, blocking labels, and noise variables at each level are sampled under the Wilson measure and its positive blocking kernels. Probabilities and heat traces are first computed on the full Gauss-law space. The reference measure on a slice is the pushforward of product link Haar measure to the gauge quotient. For densities \(p,q\) with respect to a common reference measure, total variation distance is \(d_{\rm TV}(p,q)=\tfrac12\|p-q\|_1\). Throughout this appendix, asymptotic bounds refer to the cutoff sequence chosen in Theorem~\ref{thm:main}.

\begin{lemma}[Short-time heat trace at a common bare coupling]\label{cam-bound-pressure}
Let \(F_N(R)\) be the log partition function per lattice site on a rectangular torus \(R\). At a common bare coupling, the reference square tori \((L_0/8)^4\) and \((2L_0)^4\) satisfy the bound from Appendix~\ref{lem:geometry-pressure},
\[
 0\le F_N((L_0/8)^4)-F_N((2L_0)^4)\le C_{b,v}a_N^4 .
\]
Thus, at spatial period \(L_0\) and \(s_*=L_0/4\),
\[
 1\le Z_N^0(s_*)\le
 Z_N(s_*)^2/Z_N(2s_*)\le \exp(2C_{b,v}s_*L_0^3).
\]
Here \(Z_N(s)\) is the heat trace at time period \(s\) before vacuum normalization, and \(Z_N^0(s)\) is the vacuum-normalized heat trace.
\end{lemma}
\begin{proof}
Take the positive Wilson transfer operator \(K\) along any coordinate axis. Its Schatten norm \((\operatorname{Tr}K^p)^{1/p}\) is nonincreasing in the integer \(p\). It follows that \(F_N(R)\) is nonincreasing in each side length, so both \(F_N(L_0^3,s_*)\) and \(F_N(L_0^3,2s_*)\) lie between the pressures on the reference tori. At spatial period \(L_0\), the transfer operator satisfies \(\operatorname{Tr}K^{2r}\le\|K\|^r\operatorname{Tr}K^r\) for \(r=s_*/a_N\), giving the upper bound on the vacuum-normalized trace. The two rectangular tori have lattice-site counts in the ratio \(1:2\). The log partition-function ratio is therefore bounded by \(2s_*L_0^3C_{b,v}\). All comparisons use the same bare plaquette weight.
\end{proof}

\subsection{Threshold Exceedances on the Short Square Torus and Long Cylinder}

On the square torus, \cite[Definitions 5.1--5.2]{armstrong} gives the Wilson blocking kernel and its pushforward density. After removal of a field-independent scalar, the density is \(e^{\Phi_f}\mathcal F_f\). The exponent \(\Phi_f\) contains direct plaquette terms, squares of path combinations, and regular interactions. The factor \(\mathcal F_f\) is the sum of signed weights over all compatible exact covers. A vertex carries a large-field flag when it is adjacent to a plaquette of chord length greater than \(t_f\). If there are no flags, only the empty cover contributes and \(\mathcal F_f=1\) \cite[Eq.~(13.13)]{armstrong}. The full density is strictly positive. We take absolute values of the individual cover weights only when estimating upper bounds.

\begin{lemma}[Joint plaquette exceedances under the original Wilson measure]\label{cam-bound-flags}
On the square torus with time period \(L_0\), the probability that each of \(j\) distinct prescribed plaquettes has chord length greater than \(\tau_f\) is at most \(e^{-j\Gamma_f^\square}\). The level \(f\) is sufficiently fine that \(D_f\) exceeds the cell size required in the proof. There is a tiling cell side \(S_f^{\rm cell}\), measured in level-\(f\) lattice steps, that is a product of successive blocking factors and is bounded by a fixed power of \(1+\log\lambda_f\), such that
\[
 \Gamma_f^{\square}\ge
 c\left[\frac{M_f^{1/5}}
 {(S_f^{\rm cell})^4(1+\log\lambda_f)^{2A_0}}
 -C(1+\log\lambda_f)\right].
\]
At the same bare coupling, the rate \(\Gamma_f^{\rm rect}\) on the long cylinder with time period \(T=kL_0\) satisfies the same bound with an additional subtraction of \(C\log Z_N^0(L_0)/n_f^4\) inside the brackets. Write \(\Gamma_f\) for the rate in the geometry under consideration. In a collection of sufficiently separated cells, the joint probability that the local protective-ball conditions fail satisfies the same exponential estimate, with the exponent multiplied by a fixed positive constant.
\end{lemma}
\begin{proof}
In a configuration with \(j\) prescribed threshold exceedances, the direct or linearly truncated action term for each prescribed plaquette contributes at least
\[
 c\,\beta_f c_{{\rm lin},J_f}\tau_f^2
 \ge c'\frac{M_f^{1/5}}{(1+\log\lambda_f)^{2A_0}}.
\]
Here \(J_f\) is the range parameter for the level-\(f\) blocking prediction, and \(c_{{\rm lin},J_f}\) is the coefficient of the linear truncation of the direct term. For a chord length at most \(t_f\), we use the square term. Above \(t_f\), we use the lower bound from the linear term at \(t_f\). The exponential constants in \cite[Lemmas 11.12, 11.14, 13.5, 13.8, and 13.11--13.12]{armstrong} are chosen stage by stage. The constant denoted there by \(\kappa_B\) is at most \(c_{{\rm lin},J_f}/16\). Write \(\kappa_{\rm st}^{\rm pre}\) for the smaller exponential constant at that stage. Then \(\kappa_{\rm fus}=\kappa_{\rm st}^{\rm pre}M_f^2/2\) is at least \(M_{f-1}^{1/5}\), as required at the adjacent level. The comparability of the adjacent \(M_f,M_{f-1}\) and the bound \(\beta_f\asymp\lambda_f\) give the second inequality. At the bare level, we use the quadratic lower bound on the Wilson plaquette action directly.

In the numerator of the event probability, we may drop the remaining nonnegative squares. The length moments of the path weights control the absolute sums per site of the cubic, quartic, and remaining path terms \cite[Lemma 12.3]{armstrong}. Assign each signed weight \(K_Y\) to one vertex \(o\) in its support \(Y\), and relax the restriction that distinct supports be disjoint. The resulting sum of absolute weights is at most \(\prod_o(1+\sum_{Y\ni o}|K_Y|)\), yielding the bound \(e^{Cn_f^4}\). For the lower bound, require every link to lie in the ball of radius \(c\lambda_f^{-C}\) about the identity. On this event, the direct plaquette terms, path combinations, and regular terms are uniformly bounded at each site, and the cover factor is one. The product Haar volume gives the lower bound \(e^{-Cn_f^4(1+\log\lambda_f)}\) on the normalizing integral. The probability of simultaneous exceedances at the prescribed \(j\) plaquettes is thus at most
\[
 \exp\!\left[Cn_f^4(1+\log\lambda_f)
    -c j M_f^{1/5}(1+\log\lambda_f)^{-2A_0}\right].
\]

The event can be expressed as a local test of the bare Wilson field and the blocking random variables. The level-\(i+1\) prediction depends on the level-\(i\) field within radius \((J_{i+1}+3)a_i\), where \(J_{i+1}+3\lesssim B_i^{\rm cov}\). The ratio of adjacent covering radii is at most \(\sqrt2\), while the blocking ratio is at least \(3\). Summing the dependence radii over the levels bounds the radius of the bare data needed for a single level-\(f\) plaquette test by \(C(1+B_f^{\rm cov})a_f\). Conditional on the bare field, tests with disjoint dependence regions use independent auxiliary random variables. The prediction rule is equivariant under link reflections \cite[Definition 5.1 and Lemmas 11.4 and 14.29]{armstrong}.

Starting at level \(f\) and moving toward coarser levels, multiply the successive blocking factors until the product first exceeds \(16(C+1)(1+B_f^{\rm cov})\). Write the resulting product as \(S_f^{\rm cell}=D_f/D_{j(f)}\), where \(j(f)\) is the level at which the multiplication stops. The weak-coupling proof uses \(f\ge m\) with \(n_m\asymp b^\alpha\), fixed \(\alpha>0\), and polylogarithmic \(P\). At fixed coupling, we use only levels above a sufficiently fine auxiliary level. In both ranges the required product exists. It is bounded by a fixed power of \(1+\log\lambda_f\), and \(PD_{j(f)}\) is even. Averaging over translations of the tiling, a fixed fraction of the prescribed plaquettes have their entire test dependence regions inside a single cell. Each cell contains at most \(C(S_f^{\rm cell})^4\) prescribed plaquettes, so retain one test in each occupied cell. The link-reflection chessboard estimate \cite[Eq.~(14.19)]{armstrong} disseminates a single-cell test into a simultaneous event on \(n_f^4/(S_f^{\rm cell})^4\) distinct plaquettes. Taking the root of the probability bound with exponent equal to the reciprocal number of cells, and then multiplying over the retained cells, gives the square-torus bound.

On the long cylinder, first disseminate the same test along the three spatial directions. Let \(G\) be the transfer operator for a time slab of physical thickness \(s=S_f^{\rm cell}a_f\) with the test inserted. The neighboring slab carries the reflected test and has operator \(G^*\). The field variables and auxiliary random variables used by a test all lie inside its slab, so the short square torus and the long cylinder use the same \(G\). After dissemination, the numerators are \(\operatorname{Tr}\bigl[(G^*G)^{L_0/(2s)}\bigr]\) and \(\operatorname{Tr}\bigl[(G^*G)^{kL_0/(2s)}\bigr]\). Schatten monotonicity bounds the long-cylinder numerator by the \(k\)th power of the short-torus numerator, while the long-cylinder denominator is at least \(\|K_N\|^{kL_0/a_N}\). After taking the corresponding cell root, the logarithmic bound increases by \(\log Z_N^0(L_0)/(L_0/s)^4\) per cell. Dividing by the number of candidate plaquettes per cell produces the last term of \(\Gamma_f^{\rm rect}\).

Given the fine field and the blocking labels at this step, the output links are conditionally independent at the final noise update. The multiplicative error \(W\) relative to the prediction has density \(Z_\gamma^{-1}e^{-\gamma(1-\tfrac12\operatorname{Re}\operatorname{Tr}W)}\), where \(\gamma\asymp\lambda_f\) is the noise parameter and \(Z_\gamma\) is the normalization integral for Haar measure. A Haar ball of radius \((1+\gamma)^{-1/2}\) gives \(Z_\gamma\ge c(1+\gamma)^{-3/2}\). Thus the probability of chord length greater than \(\tau_f\) is at most \(C(1+\gamma)^{3/2}e^{-c\gamma\tau_f^2}\). Under the stated conditioning, these bounds multiply over distinct output links.

Fix the block spanning trees used by the prediction. An off-tree link inside a block is compared with the coarse-field representative using the loop it forms with the tree path. For a link between blocks, make the comparison by moving the path between the roots across successive plaquettes. The number of fine and coarse plaquettes and the number of links involved are bounded by fixed powers of \(1+\log\lambda_f\). The chord length is subadditive under group multiplication, and the derivatives of the local inverse associated with the small-curvature prediction are also bounded by fixed logarithmic powers. Thus, for sufficiently large \(A_0\), a local link coordinate can leave its protective ball only if some fine plaquette, coarse plaquette, or noise error exceeds \(\tau_f\). If the protective-ball tests fail in a collection of separated cells, at least a third of them have an exceedance of the same type. Retain a fixed fraction at distinct locations and apply the joint probability estimate for that type. The number of candidate locations per cell is bounded by a fixed logarithmic power and is absorbed by the exponential bound.
\end{proof}

\subsection{Finite-Range Conditional Laws and Spatial Mixing Across Buffers}

Let \(U=U_f\) be the finer field, \(V=U_{f-1}\) the next coarser field, and \(\mathcal K_f(U,dV,d\xi)\) the blocking kernel with label \(\xi\). By \eqref{eq:geom-history-resets}, the six directional curvature corrections between the rectangle and square torus reweight the joint blocking law by \(\exp(\sum_f\mathcal D_f(U_f))\). Equation~\eqref{eq:geom-total-reset} gives
\(B_{\rm tot}=\sum_f\operatorname{osc}\mathcal D_f=O_A(b^{-A})\).
Write \(\rho_f^{\rm pre}\) for the positive density incorporating the corrections from the bare level down to level \(f\). Integrating the corrections from level \(f-1\) onward gives a positive function \(w(V)\). The joint measure on the adjacent levels and the label is therefore proportional to
\begin{equation}
 \rho_f^{\rm pre}(U)\,dU\,\mathcal K_f(U,dV,d\xi)\,w(V).
 \label{cam-bound-prefix}
\end{equation}
Here \(dU\) is product Haar measure at level \(f\). The factor \(w(V)\) cancels on conditioning on \(V\). The density ratio of the two full blocking laws lies between \(e^{-B_{\rm tot}}\) and \(e^{B_{\rm tot}}\). On the square torus, take \(\mathcal D_f=0\).

\begin{lemma}[Finite-range approximation by positive densities]\label{cam-bound-positive}
Fix a sufficiently small \(\zeta>0\). The conditional law of \(U,\xi\) given \(V\) can be approximated by replacing the positive density \(\rho_f^{\rm pre}\) in \eqref{cam-bound-prefix}. If \(n_f\le\lambda_f^K\), the approximating density has interaction range \(b^{\zeta+o(1)}\) in level-\(f\) lattice units, and its total variation error, averaged over the actual \(V\)-marginal, is \(O_A(b^{-A})\). If \(n_f>\lambda_f^K\), the range can be taken to be
\[
 R_f=C(1+\lambda_f)^D(1+\log n_f),
\]
with average error \(O_p(n_f^{-p})\) for any prescribed \(p>0\). First take the fixed exponent \(D\) sufficiently large and then take \(K>D+2\). On the long cylinder, the estimate requires \(\Gamma_f^{\rm rect}\) in Lemma~\ref{cam-bound-flags} to exceed every fixed power of \(1+\log\lambda_f\) at the levels used.
\end{lemma}
\begin{proof}
When \(n_f\le\lambda_f^K\), the trajectory bound \(\lambda_f\asymp b+\log n_f\) gives \(\lambda_f=O_K(b)\), and there are at most \(O_K(\log b)\) such levels. The number of lattice sites is bounded by a fixed power of \(b\). By Lemma~\ref{cam-bound-flags}, the probability that any required local test fails at this level or the adjacent coarser level is \(O_A(b^{-A})\). On the remaining configurations, truncate each path combination at path length \(b^\zeta\) before squaring. Truncate the cubic, quartic, and quadratic loop terms by total path length, and the regular terms by the radius of their complete dependence regions. Exponential weights in path length and interaction support give the action error \cite[Lemma 12.3 and Definition 5.2]{armstrong}
\[
 C k n_f^4\lambda_f^C\exp[-c b^\zeta/(1+\log\lambda_f)^C].
\]
On the event with no flags, the full cover factor is one. Multiplying the truncated exponential density by the indicator of that event gives the nonnegative approximating density.

When \(n_f>\lambda_f^K\), take a cell side scale \(l=\lceil(1+\lambda_f)^2\rceil\), a connection distance \(r=(1+\lambda_f)^D\) for flagged vertices, and a count bound \(j_*=1+\lceil C\log n_f\rceil\). Set the auxiliary threshold to \(h_{\rm inv}=t_f/(2C_*^{L_{\rm col}})\), where \(C_*\) is the curvature bound in the filling construction below and \(L_{\rm col}\) is its coloring number. Both depend only on dimension and the gauge group. Let \(\mathcal E_f\) be the event that every lattice ball of radius \(6l\) contains at most \(j_*\) plaquettes of chord length greater than \(h_{\rm inv}\), and every connected component of the graph joining flagged vertices at lattice distance at most \(r\) contains at most \(j_*\) vertices.

In enumerating prescribed plaquettes in a ball or rooted trees from a flagged vertex, there are at most \(C(1+\lambda_f)^8\) or \(C(1+r)^4\) choices per step, respectively. A plaquette can flag only a bounded number of neighbors. An oversized component contains a connected subset with exactly \(j_*\) vertices. The probability of the exceptional event is therefore at most
\[
 Ckn_f^4\exp[-c(\Gamma_f-C\log(1+\lambda_f))j_*]
 =O_p(n_f^{-p}).
\]
Here \(k\lesssim\lambda_f^{1/3}\). Both the flag threshold \(t_f\) and the auxiliary threshold \(h_{\rm inv}\) exceed \(\tau_f\), so Lemma~\ref{cam-bound-flags} applies to both exceptional events under the original blocking law.
Consider a four-dimensional lattice cube whose three-dimensional boundary has all plaquette chord lengths at most \(h\). It has a filling that fixes every boundary link and bounds every interior plaquette chord length by \(C_*h\). If \(h=0\), the boundary is simply connected, and a spanning-tree gauge gives a flat filling. Suppose \(h>0\) and the cube has side \(r_0\). If \(r_0\sqrt h\) exceeds a fixed small constant, use the same integer subdivision of length \(\asymp h^{-1/2}\) along each coordinate axis to obtain a coarser cell decomposition of the boundary. Choose gauges at its vertices. Along each coarse edge, distribute the discrepancy at the endpoint of the original link transport in equal increments along a shortest group geodesic. Every fine link on that edge then has chord length \(O(\sqrt h)\). If \(r_0\sqrt h\) is smaller, use the eight three-dimensional facets of the cube and fix a spanning tree of its coarse edge graph to set the gauges at the corners. Each coarse plaquette is tiled by at most \(r_0^2\) fine plaquettes, and each tree loop can be filled with a fixed number of coarse plaquettes. The discrepancy at a coarse endpoint is therefore \(O(r_0^2h)\). Equal subdivision makes every fine link on the coarse edge \(O(r_0h)\) in chord length.

In each two-dimensional cell of the decomposition and then each three-dimensional cell, take a local gauge along ordered coordinate paths. The paths to neighboring endpoints differ by moving one coordinate step across the other segments. The strip between the paths contains \(O(l_0)\) fine plaquettes, where \(l_0\) is the coarse edge length in lattice units. Every fine link in this local gauge has chord length at most \(Cl_0h\). After scaling the cell to unit size, the transition map between the local gauge and the gauge already fixed on the lower-dimensional skeleton has a uniform Lipschitz bound. In the small-\(r_0\sqrt h\) case, that bound is \(Cr_0^2h\). Interpolate first on the small boundary simplices and normalize to unit quaternions. For \(d-1\le2\), the volume of the \(\eta\)-neighborhood in \(S^3\) of the image of a \((d-1)\)-dimensional map with Lipschitz constant at most \(L\) is bounded by \(C(1+L)^{d-1}\eta^{3-(d-1)}\). Choose a pole a positive distance from the image. Use stereographic projection to extend each real coordinate as a Lipschitz function, and then apply the inverse projection. In the small-scale case, the same extension in a local logarithm chart preserves the bound \(Cr_0^2h\). Extending the gauge in order of dimension thus bounds the logarithm norm \(b_0\) of every boundary link by \(C\sqrt h\) or \(Cr_0h\), respectively.

Pull the boundary field back to an inner collar of width \(w\asymp\min(r_0,h^{-1/2})\). Center the cube at the origin and round the image of \(x\mapsto (r_0/2)x/\|x\|_\infty\) to a boundary lattice point, fixing points already on the boundary. The images of neighboring lattice points can be joined by a boundary path of length \(O(1)\). The boundary loop corresponding to a lattice face can be filled by \(O(1)\) boundary plaquettes. Indeed, the loop lies in a uniformly bounded lattice box and meets at most one exterior facet for each coordinate direction. Successively moving along coordinate directions toward the common corner of those facets fills the loop. The plaquette chord lengths of the pullback are thus \(O(h)\). Multiply each link logarithm of the pullback by a function of the depth from the boundary that decreases from one to zero and varies by \(O(1/w)\) on adjacent links. The quadratic remainder in the Baker--Campbell--Hausdorff formula on each plaquette gives chord length \(O(h+b_0^2+b_0/w)=O(h)\). Set the remaining interior links to the identity and apply the inverse gauge transformation. When the side length is bounded and small, one may extend the boundary spanning-tree gauge directly to the interior, since every relevant loop can be filled by a bounded number of plaquettes.

Tile the four-torus by cells of side between \(l\) and \(2l\). About the center of each cell, consider lattice cubes of integer radii between \(3l\) and \(5l\). A single above-threshold plaquette can rule out at most a fixed number of candidate radii. On \(\mathcal E_f\), the inequality \(j_*\ll l\) allows a cube to be chosen in every cell whose tangential boundary plaquettes all have chord length at most \(h_{\rm inv}\). Use the largest candidate cubes to define the adjacency graph. Its degree is uniformly bounded, so \(L_{\rm col}\) colors suffice. Apply the filling only in cells containing a plaquette above the flag threshold \(t_f\), in the order of the colors. Each filling increases the relevant curvature bound by at most \(C_*\), and \(C_*^{L_{\rm col}}h_{\rm inv}<t_f\). The specified large-field flags can therefore be removed without creating new ones. Each component requires changes to at most \(N_f^{\rm edit}\le Cj_*l^4=O((1+\lambda_f)^9)\) links. Distinct components are separated by more than \(r\gg l\), so they can be modified separately. If a measurable choice of filling is needed, minimize the coordinates successively in a fixed order over the compact set of configurations satisfying the closed curvature constraints.

A change to one link changes \(\log\mathcal F_f\) by at most \(C_f^{\rm edge}\le C(1+\lambda_f)^C\). To see this, condition on the finer field and its blocking labels. The ratio of the two noise densities on that link lies between \(e^{-2\gamma}\) and \(e^{2\gamma}\). Integration against the positive fine-field density preserves that ratio bound. The corrections at finer levels stay inside this positive integral, and the correction at the present level adds at most its oscillation. At the finest, bare level, a link affects only a bounded number of bounded Wilson plaquette terms. To remove \(e^{\Phi_f}\), factor the difference of squares for the change in each squared path combination and sum with the length weights over paths containing the affected link. For the regular interactions, sum only the absolute values of terms that depend on this link, with their support weights. The total of these contributions and the finitely many adjacent direct terms is at most \(C(1+\lambda_f)^C\).

Let \(I_i\) be the \(i\)th connected component of large-field flags and \(\Omega_i\) its \(r/4\)-neighborhood. Define \(Z_i\) as the local cover sum restricted to supports lying entirely in \(\Omega_i\) and containing a flagged vertex in \(I_i\), with their original admissibility and compatibility conditions. Remove the other components to obtain \(U^i\), and then remove \(I_i\) to obtain the flag-free configuration \(U^0\). The change from \(U\) to \(U^i\) leaves \(Z_i\) unchanged, and \(\mathcal F_f(U^0)=1\) gives \(\mathcal F_f(U^i)\ge e^{-C_f^{\rm edge}N_f^{\rm edit}}\). Let \(\epsilon_{\rm for}\) be the upper bound on the exponentially weighted absolute sum at a given lattice site of the level-\(f\) forced-cover weights. The sum of the absolute weights of eligible terms involving only this component's flags is at most \(j_*\epsilon_{\rm for}\). A support extending beyond \(\Omega_i\) needs at least \(cr/B_f^{\rm cov}\) covering balls. Reserving half of the exponential support weight bounds the absolute sum over covers containing such a nonlocal term by
\[
 j_*\epsilon_{\rm for}
       \exp(j_*\epsilon_{\rm for}-c r/B_f^{\rm cov}).
\]
Hence \(Z_i\ge \frac12e^{-C_f^{\rm edge}N_f^{\rm edit}}\).

Given a compatible family of nonlocal cover terms, sum the rest separately in the \(\Omega_i\). A neighborhood \(\Omega_i\) that no nonlocal support meets contributes the full \(Z_i\). All other local sums have absolute value at most \(e^{j_*\epsilon_{\rm for}}\). Let \(d_Y\) be the smallest number of lattice balls of radius \(B_f^{\rm cov}\) that cover \(Y\) and have connected union. A tree through their centers has length at most \(CB_f^{\rm cov}d_Y\). Select representative points every \(r\) steps along the tree. The separation between the components then shows that \(Y\) can meet at most \(C(1+B_f^{\rm cov}d_Y/r)\) of the \(\Omega_i\). Each affected local sum divided by \(Z_i\) has absolute value at most \(2\exp(C_f^{\rm edge}N_f^{\rm edit}+j_*\epsilon_{\rm for})\). Take \(D\) large enough that \(r/B_f^{\rm cov}\) dominates this exponent and \(\log(kn_f^4)\). Each nonlocal term then uses at most half of its exponential support weight. Bounding the absolute sum over compatible nonempty families by the sum over arbitrary families gives
\[
 \mathcal F_f=\left(\prod_i Z_i\right)
             (1+O(kn_f^4e^{-cr/B_f^{\rm cov}}))
 \quad\text{on }\mathcal E_f .
\]
This is a relative error for the full signed cover sum. Each term is still evaluated with its original support and admissibility conditions.

Truncate the regular interactions and path terms at \(r\) and multiply by the indicator of \(\mathcal E_f\). For each flagged vertex, explore its component \(I_i\) up to the count bound \(j_*\) and use the local factor
\(\bigl(\max(Z_i,\tfrac14 e^{-C_f^{\rm edge}N_f^{\rm edit}})\bigr)^{1/|I_i|}\).
A radius \(Cr(1+j_*)\) suffices to determine the component and verify the required local conditions. Use the pointwise relative error on \(\mathcal E_f\) and the original probability of the exceptional event on its complement. Combining the two bounds and normalizing gives \(O_p(n_f^{-p})\) in total variation. Application of the probability kernel \(\mathcal K_f\) is a contraction in total variation. If two joint laws differ by \(\eta\), the total variation distance of the fine-field conditional laws, averaged over the first coarse-field marginal, is at most \(2\eta\). This follows by first replacing the marginal and then applying the triangle inequality. Equation~\eqref{cam-bound-prefix} permits this average under the original coarse-field marginal.
\end{proof}

\begin{lemma}[Conditional mixing across a common buffer]\label{cam-bound-cross}
Fix the coarse field \(V\). In the finite-range conditional law of Lemma~\ref{cam-bound-positive}, take a separating shell with a protective layer on each side. Suppose the shell and both protective layers satisfy the small-curvature conditions and that the interior variables are integrated over a common product of linkwise balls. Let \(x,y\) be the two boundary values, \(K(x,y)\) the resulting kernel, and \(d_{\rm shell}\) the distance between the two sides in level-\(f\) lattice units. Then
\[
 \left|\log\frac{K(x,y)K(x',y')}{K(x',y)K(x,y')}\right|
 \le C(kn_f^4\lambda_f)^C
             \exp[-c\,d_{\rm shell}/\lambda_f^{o(1)}].
\]
If two coarse fields \(V\) differ only outside a common window, the conclusion applies to shells whose full dependence regions lie inside that window. The exterior may contain large-field components and arbitrary exterior weights.
\end{lemma}
\begin{proof}
Within each original block, fix the axial spanning-tree gauge with respect to its root. The tree variables give independent Haar factors, which can be restored when taking endpoint observables. Let \(U^{\rm rep}=\operatorname{rep}V\) be the blockwise representative of \(V\). Write the remaining fine links as \(U_e=e^{X_e}U_e^{\rm rep}\), and use Euclidean coordinates \(z_e=\sqrt{\beta_f}X_e\). Let \(r_{\rm ball}\) be the radius of the range of each \(X_e\) and \(h_0\) the bound on coarse-plaquette chord lengths in the protective layer. Choose \(Cr_{\rm ball}+h_0<h_{\rm inv}\) and take the protective layer to contain every link needed by the action terms that depend on the interior variables. Any fine plaquette that varies with an interior variable remains below \(h_{\rm inv}\). The large-field flags, count constraints, and plaquette indicators are constant on the common integration domain. A large-field cover term depending on an interior variable would have to have a flagged vertex in the protective layer, which has none. Pure exterior factors remain on their respective sides.

The blocking prediction in this region uses the small-curvature branch of \cite[Definitions 11.2--11.3 and Lemma 11.4]{armstrong}. In the local tree gauge, comparing two paths of transport between block roots requires exchanges across at most a fixed logarithmic power of coarse and fine plaquettes. For derivatives, use the frame equation in quaternion form. Write \(T_{rs;o}\) for transport along the \(o\)th template path between block roots \(r,s\), and let \(\mathcal P(q)=q/|q|\) be radial projection to the unit quaternions. Up to a common left multiplication, the frame \(h_r\) satisfies \(h_r=\mathcal P(F_r)\), where \(F_r=\operatorname{av}_{s,o}h_sT_{rs;o}^{-1}\) averages over template roots and paths. Normalization at root \(0\) gives the fixed-point map \(\mathcal P(F_0)^{-1}\mathcal P(F_r)\). At a flat field its first derivatives in the \(h_s\) cancel. Within distance \(\delta_{\rm fr}\) of a flat frame, the derivative has block-row norm at most \(C\delta_{\rm fr}\). Take this norm below \(1/2\) and differentiate the fixed-point equation successively, using an inverse matrix of norm at most two. The lengths of the input paths and the first three derivatives of the output prediction are all bounded by fixed logarithmic powers. The label weights on the small-curvature branch are constant, and the prediction is independent of these labels, so integrating them out preserves the prediction.

Let \(d_0,d_1\) denote the lattice gradient and curl on the flat template. Decompose a fine-link tangent vector \(u\) as \(d_0\phi+u_{\rm rem}\), where \(\phi\) is the sum of \(u\) along the template spanning tree. The number of lattice faces needed to pass between adjacent tree paths is bounded by a fixed power of the template size \(R\), so
\(\|u_{\rm rem}\|_{\infty}\le CR^C\max_{\rm patch}|d_1u|\).
The tangent vector vanishes on fine tree links, and \(\phi\) inside a block differs from its value at the block root by at most \(CR^C\max|d_1u|\). Let \(Q_E(U)\) be the prediction at coarse link \(E\) and \(Q_0\) its derivative at a flat field. Gauge covariance gives
\((Q_0d_0\phi)_E=\phi(\text{terminal root})-\phi(\text{initial root})\).
The predicted derivative of the remainder has supremum norm at most \(CR^C\|u_{\rm rem}\|_\infty\). Thus the tangent vector on a fine link between blocks is controlled by \(|(Q_0u)_E|+CR^C\max|d_1u|\). A fine link within a block requires only the curvature term. In a coarse axial gauge chosen separately on each template, local transport differs from its flat value by \(O(R^C(h_0+r_{\rm ball}))\). This estimate also holds when the coarse field has nontrivial flat holonomy. The difference between the derivatives of the actual defects \(U_p-1\) and \(V_EQ_E(U)^{-1}-1\) and their flat counterparts is bounded by this small quantity times \(\max_{\rm patch}|u|\) and a fixed logarithmic power. Square and sum over the links. The template size and overlap multiplicities are bounded by fixed logarithmic powers, so the derivative error can be absorbed in the quadratic-form lower bound for all interior directions.

The direct plaquette terms and the prediction-error terms are squares of defects, so their Hessians contain positive contributions of the form \(|D(\text{defect})u|^2\). The preceding lower bound, \(\beta_f\asymp\lambda_f\), and the lower bound on the direct coefficients give the modulus \(c_f^{\rm cv}\ge cM_f^{-2}(1+\log\lambda_f)^{-C}\) in \(z\)-coordinates. The squared path combinations contribute nonnegative derivative squares as well. For a path combination \(\mathcal L\) near the interior variables, small curvature along short paths and the exponential bound on long paths give \(|\mathcal L|\le C(1+\log\lambda_f)^C(t_f+\lambda_f^{-4})\). This controls the possibly negative term \(\mathcal L D^2\mathcal L\). A more distant combination can contribute to the Hessian only if at least one of its paths passes through the relevant link. The exponential length weights absorb the powers of path length introduced by two derivatives and the multiplicity of the anchors \cite[Lemma 12.3]{armstrong}. The cubic and quartic path terms retain at least one small curvature factor after two derivatives. The contributions of the quadratic loop terms, regular interactions, and Haar Jacobian are controlled by
\[
 C(1+\log\lambda_f)^C\!
 \left[
 \frac{\Lambda_f^{\rm loop}}{\beta_f}
 +\frac{\epsilon_f^{\rm wide}}{\beta_f t_f^2}
 +\frac1{\beta_f}\right].
\]
Here \(\Lambda_f^{\rm loop}\) is the row-sum bound for the quadratic loop terms and \(\epsilon_f^{\rm wide}\) is the weighted norm bound on the regular interactions on the rectangle from Appendix~\ref{app:geometry}. The second ordered left derivatives of a regular term cost \(t_f^{-2}\), and passing to derivatives in \(z\) divides by \(\beta_f\) \cite[Definitions 5.3--5.4]{armstrong}. For a support \(Y\) with covering number \(d_Y\), taking the second-derivative row sum at a fixed link introduces at most \(C(B_f^{\rm cov})^8(1+d_Y)^5\), absorbed by the exponential support weight. The estimates continue to hold with weight \(e^{d(e,e')/L_f}\) on each Hessian entry, where \(d(e,e')\) is link distance and the decay length \(L_f\) is bounded by a fixed power of \(1+\log\lambda_f\). Only a small fraction of the original exponential weights on paths and regular terms is needed. The bounds on rectangular regular and loop terms make the total negative remainder at most \(\lambda_f^{-1/512+o(1)}<c_f^{\rm cv}/2\). The path-length moments and the derivative bounds for template predictions also give a weighted row-sum bound \(\lambda_f^{o(1)}\) on the positive part of the square terms. The integration potential is therefore uniformly strongly convex in the \(z\)-coordinates at each level. The weighted Schur norm, the decay length \(L_f\), and the reciprocal convexity modulus \((c_f^{\rm cv})^{-1}\) are all \(\lambda_f^{o(1)}\).

Let \(\Psi_f(x,y,z)\) be the negative log density on the common product of balls. With \(x,y\) fixed, take the Neumann operator \(\mathcal L_f=-\Delta_z+\nabla_z\Psi_f\cdot\nabla_z\) on the interior variables \(z\). Integration by parts on gradients gives the vector operator \(\mathcal L_f+\nabla_z^2\Psi_f\), with a nonnegative second fundamental form at the boundary. Take the weight \(m_e=e^{\eta_f d(e,\partial_x)}\), where \(\partial_x\) is the side with boundary value \(x\) and \(d(e,\partial_x)\) is link distance to that side. The weight is constant on each link coordinate block and therefore preserves the corresponding vector boundary conditions. The change in the conjugated Hessian is bounded by the original exponential row weights and
\(\big|e^{\eta_f(d(e,\partial_x)-d(e',\partial_x))}-1\big|\le C\eta_f d(e,e')e^{\eta_f d(e,e')}\).
Take \(\eta_f=\lambda_f^{-o(1)}\) so that this change is less than \(c_f^{\rm cv}/4\). The resulting weighted norm bound on the inverse vector operator is \(4/c_f^{\rm cv}\).

Given \(x,y\), apply the covariance representation to the two boundary derivatives \(\partial_x\Psi_f,\partial_y\Psi_f\) in the normalized interior integral. Their covariance is
\[
 \left\langle\nabla_z\partial_x\Psi_f,
 (\mathcal L_f+\nabla_z^2\Psi_f)^{-1}\nabla_z\partial_y\Psi_f\right\rangle .
\]
The representation follows from integration by parts for the smooth Neumann variational problem and extends by \(L^2\) closure. The weighted inverse bound gives the exponential factor \(e^{-\eta_f d_{\rm shell}}\), and the same exponential weight controls the direct mixed derivative. A factor depending only on one boundary cancels from the cross-ratio. The mixed derivative of the log kernel is the stated covariance minus the mean of \(\partial_x\partial_y\Psi_f\). Interpolate linearly between the two prescribed boundary values on each side inside the common balls. Each coordinate in the \(z\)-scaling traverses a distance at most \(O(M_f)\). Integrate in the two boundary coordinates to obtain the cross-ratio bound.
\end{proof}
\begin{lemma}[Scale-by-scale factorization of the two-end conditional laws]\label{cam-bound-backward}
Take an initial observable level \(m\) with \(n_m\asymp b^\alpha\), and let \(h\asymp L_0b^{-\alpha+\delta}\) be the width scale of the two endpoint neighborhoods, where \(0<\delta<\alpha\le1/12\). Suppose the two slices are separated by at least \(L_0/2\). Let \(e_f\) be the average total variation distance at level \(f\) between the two-end conditional law and the product of the two marginals conditioned on their respective coarser endpoint observables. On the square torus, for any fixed polynomial weight \(B_f'\le C(n_f+\lambda_f)^C\),
\[
 \sup_N\sum_{f>m}B_f'e_f=O_A(b^{-A}) .
\]
The same conclusion holds on the long cylinder if the short trace on the square torus has the subpower bound of Lemma~\ref{cam-bound-entropy}. On the square torus at fixed \(b\), take an auxiliary observable level \(j\) and let \(h_j=L_0n_j^{-1/4}\). The sum of errors beyond \(j\) tends to zero as \(j\to\infty\), uniformly in finer cutoffs.
\end{lemma}
\begin{proof}
At level \(f\), reserve a physical buffer of thickness \(h_f^{\rm buf}=ch\sqrt{n_m/n_f}\) on each side. The blocking ratio is at least three, so \(\sum_{f>m}h_f^{\rm buf}<h/4\) for sufficiently small \(c>0\). At each level, shrink the observable neighborhood by \(h_f^{\rm buf}\), retaining the central slice. When \(n_f\le\lambda_f^K\), we have \(h_f^{\rm buf}/a_f\gtrsim b^\delta\sqrt{n_f/n_m}\). Take \(\zeta<\delta/4\). Both the fine-field dependence radius of the blocked observables and the range \(b^{\zeta+o(1)}\) of the finite-range approximation are then much smaller than this ratio. Under the original joint law, first discard configurations in which any local protective-ball condition fails. The remaining two conditional laws integrate the same product of balls in the common window. Lemma~\ref{cam-bound-cross} gives an average error \(O_A(b^{-A})\) per level for any prescribed \(A>0\).

When \(n_f>\lambda_f^K\), start from a given exterior region and test successive adjacent cells to determine whether both conditional samples satisfy the linkwise protective-ball and local coarse-plaquette conditions. Each cell includes the neighboring auxiliary random variables and is large enough to cover the dependence radius \(R_f\) of the approximating density. Test only cells within a fixed number of cell steps of the region already revealed. If both samples pass, retain only the two sets of common coordinate constraints. Otherwise, reveal all variables involved in the test in both samples, and add the cell to the exterior. The test order depends only on the values already revealed. A successful test is the Cartesian product of two coordinate-product sets, and a failed test is determined by the values just revealed. By successive conditioning, at the end of the procedure the two remaining conditional laws are subject only to the common ball constraints and the revealed boundary values.

If the procedure has not reached the reserved boundary of the target interior region, there is a shell satisfying the small-curvature conditions between the revealed exterior and the unrevealed interior. On each side retain a protective layer of width at least \(R_f\), and integrate the interior of the shell over the common product of balls. The two coarse fields can differ outside, while the finite-range conditional densities in their common window follow the same rule. Thus the cross-ratio bound \(\eta\) in Lemma~\ref{cam-bound-cross} controls the oscillation of the unnormalized likelihood ratio under a change of boundary values. After normalization, the total variation distance of the interior conditional laws is at most \(e^\eta-1\).

Set the cell side length in the local tests, and hence the thickness scale of the shell, to
\[
 L_f^{\rm search}
  =C_p R_f\lambda_f^{\epsilon}(1+\log n_f),
\]
in lattice units, with \(\epsilon>0\) fixed and \(C_p\) sufficiently large. Define adjacency in the cell graph using a fixed multiple of this length. If the procedure reaches the reserved boundary of the target region, trace each cell back to the adjacent cell that first caused it to be tested. This gives a chain of failed cells of length at least \(c h_f^{\rm buf}/(a_f L_f^{\rm search})\). Since \(\log n_f\lesssim\lambda_f\) and \(n_m\asymp b^\alpha\), choosing \(K>4(D+4)\) and increasing it as needed for \(\epsilon\) makes the chain length at least \(c n_f^{1/4}\).

At least half the cells on any given chain fail in the same sample. Retain a fixed fraction of separated cells and apply Lemma~\ref{cam-bound-flags}. A prescribed chain of length \(j\) then has probability at most \(2^je^{-c\Gamma_f j}\). There are at most \(Ckn_f^4\) starting points and at most \(\lambda_f^C\) choices at each step. In the one-sample marginal under the original measure, first take the union over all such chains. Then use the total variation bound of Lemma~\ref{cam-bound-positive} for the probability of this single union event. The probability of failure is \(O_p(n_f^{-p})\). If the procedure stops before the reserved boundary, taking \(C_p\) sufficiently large gives the same order for the cross-ratio bound. Choose \(p\) large enough to control the prescribed polynomial weight, and sum using geometric growth of \(n_f\).

Let \(\mathcal X,\mathcal Y\) be the \(\sigma\)-algebras generated by the coarser observables at the two ends. Given the full coarse field \(V\), let \(J_f(V)\) be the conditional joint law of the two new observables and \(J_f^\pm(V)\) its two marginals. The one-end conditional laws are \(q_f^-=\mathbb E[J_f^-\mid\mathcal X]\) and \(q_f^+=\mathbb E[J_f^+\mid\mathcal Y]\). Jensen's inequality gives
\[
\begin{split}
 e_f\le{}&
  \mathbb E\|J_f-J_f^-\otimes J_f^+\|_{\rm TV}
  +\mathbb E\|J_f^--q_f^-\|_{\rm TV}
  +\mathbb E\|J_f^+-q_f^+\|_{\rm TV}.
\end{split}
\]
For the first term, condition on \(V\) and draw two independent copies of the new far-end observable. Compare the near-end kernels conditional on these copies. For the second and third terms, condition on the coarser observable at the respective end and sample two independent copies of the full coarse field \(V\). In each case, integrating out the conditioning restores the original sample marginal, so the chain probability bound applies to all three terms.

Replace the two-end update kernels by the products of the one-end conditional kernels in order from finer to coarser levels. At each replacement, the input is still generated by the original blocking kernels that have been retained, and the probability kernels already replaced downstream contract total variation. The total error is therefore at most the sum of the layerwise average errors. After marginalizing the intermediate observables, if the resulting one-end kernel is only a candidate conditional law, take the original one-end marginal and condition again. Jensen's inequality bounds the extra error by one copy of the original error at each end. On a fixed square torus, substituting the auxiliary level \(j\) and \(h_j=L_0n_j^{-1/4}\) gives \(h_f^{\rm buf}/a_f\asymp n_j^{1/4}n_f^{1/2}\), while \(\lambda_f=O_b(1+\log n_f)\). For sufficiently large \(j\), all subsequent levels have \(n_f>\lambda_f^K\). For each \(p>0\), the preceding estimates give a weighted tail \(O_p(n_j^{-p})\) uniformly in finer cutoffs.
\end{proof}

\subsection{Square-Torus Entropy Bound and the Square-Root Representation of the Heat Kernel}

\begin{lemma}[Short trace on the weak-coupling square torus]\label{cam-bound-entropy}
For the blocking schedule and square torus of Theorem~\ref{thm:main}, each fixed \(\varepsilon>0\) satisfies
\[
 \log Z_N^0(L_0)\le C_\varepsilon b^\varepsilon
\]
uniformly in sufficiently fine cutoffs. The proof uses the threshold exceedance estimates and the scale-by-scale conditional factorization on the short square torus.
\end{lemma}
\begin{proof}
Fix \(0<\delta<\alpha\le1/12\) and choose \(n_m\asymp b^\alpha\) along the same blocking schedule. Let \(X_m,Y_m\) be the level-\(m\) group-valued link observables within half-width \(2h\asymp L_0b^{-\alpha+\delta}\) of the two opposite slices. Each comprises at most \(C n_m^3(1+h/a_m)\) links. Let \(F\) consist of the finer field data and prediction labels on which the two ends depend. Given \(F\), the output links are conditionally independent, and each density relative to Haar measure is bounded by \(C(1+\gamma)^{3/2}\), with the noise parameter \(\gamma\) bounded by a fixed power of \(b\). Write \(\mathrm{KL}\) for Kullback--Leibler relative entropy and \(I\) for mutual information. The chain rule for relative entropy and the data-processing inequality give
\[
 I(X_m:Y_m)\le I(X_m:F)
 \le \mathbb E_F\mathrm{KL}(\mu_{X_m\mid F}\|\mathrm{Haar}^{\otimes \#X_m})
 \le C_{\alpha,\delta} b^{3\alpha+\delta}\log(2+b).
\]
Only the group-valued links are needed as observables in this estimate. The prediction labels can be integrated out under their conditional law.

Let \(\mu_f\) be the joint law of \((X_f,Y_f,X_{f-1},Y_{f-1})\) at successive levels. Let \(\widetilde\mu_f\) be obtained by first sampling the original coarser two-end marginal, then drawing the new observables independently from \(\mu_{X_f\mid X_{f-1}}\) and \(\mu_{Y_f\mid Y_{f-1}}\). The definition of relative entropy gives
\[
 \mathrm{KL}(\mu_f\|\widetilde\mu_f)=
 I((X_f,X_{f-1}):(Y_f,Y_{f-1}))-I(X_{f-1}:Y_{f-1}).
\]
The left side is therefore at least \(I(X_f:Y_f)-I(X_{f-1}:Y_{f-1})\). For each output link, the ratio of upper and lower bounds on its density relative to Haar measure is at most \(e^{4\gamma}\). Apply this bound at the finer and coarser of the two successive levels, and then integrate over the original finer input. The absolute log density of \(\mu_f\) is bounded by \(Cn_f^4(\lambda_f+\lambda_{f+1})\). The same order bounds the conditional log densities, each given by the joint density divided by the corresponding marginal density. At the bare level \(f=N\), bound the bare log density by \(\beta_{N,P,N}(v)\) times the number of Wilson plaquettes in place of the term involving \(\lambda_{f+1}\). Thus
\[
 \left|\log(d\mu_f/d\widetilde\mu_f)\right|
 \le B_f^{\mathrm{lik}},\qquad
 B_f^{\mathrm{lik}}\le
 \begin{cases}
 C n_f^4(1+\lambda_f+\lambda_{f+1}),& f<N,\\
 C n_N^4(1+\lambda_N),& f=N .
 \end{cases}
\]
Integrating out extra coordinates does not enlarge the ratio of the upper and lower density bounds, so the estimate also holds on thinner endpoint bands. The symmetrized relative entropy equals \(\int\log(d\mu_f/d\widetilde\mu_f)(d\mu_f-d\widetilde\mu_f)\), so \(\mathrm{KL}(\mu_f\|\widetilde\mu_f)\le2B_f^{\mathrm{lik}}\DTV(\mu_f,\widetilde\mu_f)\). Lemma~\ref{cam-bound-backward} makes this upper bound summable over the levels. Pushing the observables forward to the original opposite slices gives an upper bound of the same order on \(I(X_N:Y_N)\), where \(X_N,Y_N\) now denote the two original slices.

Let \(u_N,v_N\) be the square roots of the two original slice marginal densities, so \(\|u_N\|=\|v_N\|=1\). The square root of the joint density of the opposite slices is the nonnegative integral kernel of
\(A_N(L_0)=S_N(L_0/2)/\sqrt{Z_N^0(L_0)}\).
The square-root affinity between the joint law and the product of its marginals is
\(\langle u_N,A_N(L_0)v_N\rangle\le1/\sqrt{Z_N^0(L_0)}\).
Jensen's inequality \(\mathrm{KL}(p\|q)\ge-2\log\int\sqrt{pq}\) yields
\(\log Z_N^0(L_0)\le I(X_N:Y_N)\).
Take \(3\alpha+\delta<\varepsilon\). This choice concerns only the observables at the different levels, and every heat trace uses the same bare coupling.
\end{proof}

On the long cylinder starting at \(n_m\asymp b^{1/24}\), Lemma~\ref{cam-bound-entropy} and \(n_f\ge n_m\) give \(\log Z_N^0(L_0)/n_f^4=o(1)\). Hence \(\Gamma_f^{\rm rect}\) is large enough for the long-cylinder conditional factorization, and Lemma~\ref{cam-bound-backward} gives an error bounded by any prescribed inverse power of \(b\) at the same bare coupling.

\begin{lemma}[Forward and reverse square-root comparisons for a common joint law]\label{cam-bound-arms}
Let \(\mu\) be the joint law of the two opposite original slices \(x,y\) on the long cylinder and their stopping-scale coarse-grained observables \(u_1,u_2\). Identify laws with their densities after choosing a common reference measure. Each of the square roots
\[
 (\mu_{xy}r_x r_y)^{1/2},\qquad
 (\mu_{u_1u_2}s_{u_1}s_{u_2})^{1/2}
\]
is at \(L^2\) distance \(O(b^{-3/16+\delta/2+o(1)})\) from \(\sqrt\mu\). Here \(\mu_{xy}\) and \(\mu_{u_1u_2}\) are the respective marginals. The forward kernel \(r\), given an original slice, generates the observable on its two-sided neighborhood from ground-state Doob paths. The reverse kernel \(s\) is the original conditional law of the slice given the coarse-grained observable at that end. These square roots give the isometric zero extensions, to a common endpoint space, of \(S_N(T/2)/\sqrt{Z_N^0(T)}\) and the square-root kernel of the joint density of the coarse-grained observables.
\end{lemma}
\begin{proof}
Lemma~\ref{cam-bound-backward} gives total variation error \(O_A(b^{-A})\) for the second density. By \(\|\sqrt p-\sqrt q\|_2^2\le 2d_{\rm TV}(p,q)\), the square-root distance is bounded by any prescribed inverse power. For the first density, the four-segment Doob path estimate of Appendix~\ref{app:cam-bridge} gives the square-root distance bound \(\sqrt{2[Z_N^0(T-4h)/Z_N^0(T)-1]}\). Since \(T-4h>L_0\), the secant bound for the decreasing convex function \(\log Z_N^0\) gives
\[
 0\le \log\frac{Z_N^0(T-4h)}{Z_N^0(T)}
 \le \frac{4h}{T-4h-L_0}\log Z_N^0(L_0).
\]
Using \(T\asymp L_0b^{1/3}\), \(h\asymp L_0b^{-1/24+\delta}\), and Lemma~\ref{cam-bound-entropy} gives the stated power, uniformly on the compact interval of dimensionless times chosen in the main text.

The two isometric embeddings come from normalized conditional densities, for example \(f(x)\mapsto f(x)\sqrt{r_x(u_1)}\). Each conditional density integrates to one, so the kernels of the zero extensions are the square roots of the corresponding joint densities. The coarse-grained observables depend only on the original field and auxiliary random variables near the slices. We may therefore first construct the square-root operators on the product space of these variables and then push forward to the desired observables. Both comparisons preserve equivariance under inversion, center transformations, and translations of the blocked lattice. A center transformation within a block can be moved to the interblock links by a gauge transformation fixing the root, so the conditional kernels remain center-equivariant. Compressing to the subspaces of these shared symmetries does not increase Hilbert--Schmidt norms. Section~\ref{sec:soft-spectrum} identifies the momentum sector for translations on the fine lattice.
\end{proof}

\begin{lemma}[Reverse comparison with a fixed observable]\label{cam-bound-marked}
Fix \(b,v,L_0\). Between the opposite original slices \(x,y\) on the short square torus, take a coarse-grained observable \(\mathcal O\) at a fixed level whose entire dependence region is a positive distance from both slices. At an auxiliary level \(j\), record \(u_1,u_2\) in the two slice neighborhoods of half-width \(h_j=L_0n_j^{-1/4}\). As \(j\to\infty\), the total variation distance between the original joint law of \((x,u_1,y,u_2,\mathcal O)\) and
\[
 \mu_{u_1u_2\mathcal O}\,\mu_{x\mid u_1}\mu_{y\mid u_2}
\]
tends to zero uniformly in finer cutoffs.
\end{lemma}
\begin{proof}
Given the full level-\((f-1)\) field \(V\), the observable \(\mathcal O\) is a function only of \(V\) and the subsequent coarsening variables. It is therefore conditionally independent of the update to the finer level. The three comparisons of Lemma~\ref{cam-bound-backward} still apply. The two coarse-field samples are drawn respectively conditional on all available observables including \(\mathcal O\), and conditional only on the available observable at the given end. Integrating over the conditioning recovers the original coarse-field marginal in both cases. The far-end data and the dependence region of \(\mathcal O\) lie outside the local window, with a buffer of thickness \(h_f^{\rm buf}\) between them and the local target. The original bounds for chain probabilities, shell cross-ratios, and the total variation cost of the finite-range approximation therefore remain valid. Replace the fine-level update kernels at both ends, one level at a time, keeping \(\mathcal O\) in the initial joint law. Each one-end kernel depends only on its own coarse-grained observable. Integrating out the intermediate observables proves the claim.
\end{proof}

For fixed \(b\), Lemma~\ref{cam-bound-pressure} gives the short-time heat-trace bound. Combining Lemma~\ref{cam-bound-marked} with the forward comparison with an observable in Appendix~\ref{app:cam-bridge} constructs, from the same joint law, the Hilbert--Schmidt kernel with any bounded real observable inserted. First fix the auxiliary level and remove the original cutoff, and then send the auxiliary level to finer scales. The error bounds are uniform in finer cutoffs and apply simultaneously to all the variables on which the observable depends and to the auxiliary blocking variables. The fixed-volume spectral limit and the continuous-time observable representation in the main text use this order of limits.
\section{Coulomb coordinates and effective heat-kernel comparison}
\label{app:cam-soft}

We compare Wilson observables at two endpoints at the renormalization stopping scale with the heat kernel of the nine-dimensional quartic model. We retain the endpoint parameter \(v\) and the weak-coupling parameter \(g=(4/v)^{1/2}\) from the main text. The actual couplings at the stopping scale are \(\beta_m=\beta_{N,P,m}(v)\) and \(g_m=(4/\beta_m)^{1/2}\). Equation \eqref{eq:paper-coupling} gives \(g_m/g=1+O((\log b)/b)\) uniformly over the chosen real endpoint band. We express asymptotic bounds in \(g\) and use the actual \(g_m\) in the action and time scales. The spatial lattice is \(\Lambda_n=(\mathbb Z/n\mathbb Z)^3\), with side length \(n=PD_m\asymp g^{-1/12}\) in lattice units and spacing \(a_m=L_0/n\). There are \(q=T/a_m\asymp ng^{-2/3}\) time steps in the full trace period \(T\). The small-curvature threshold \(t=t_m\) satisfies \(t/g_m\asymp \exp((\log\log \widehat\beta_m)^2)\) and \(t=g^{1-o(1)}\). For the blocked curvature-path coefficients and the comparison of positive measures after extraction of local polymer factors, we use the rectangular estimates above. The analytic action is the real action obtained by replacing the cutoff functions by their constant values on the small-curvature domain. All Gaussian normalization constants refer to the same original product Haar measure.

\subsection{Curvature-path symbols and Coulomb coordinates}

At the identity background connection let \(p\in(\mathbb R/2\pi\mathbb Z)^4\) denote the four-dimensional Fourier variable and let \(Q(p)\) be the Hermitian symbol of the quadratic curvature matrix on the six plaquette directions. Lemma 12.3 of the published blocking construction gives exponential length moments for the curvature-matrix coefficients and the path coefficients of degrees two through four. Lemma 12.4 gives
\begin{equation}
 Q(0)=I,\qquad Q(p)\succeq \bigl(c_{\mathrm{dir}}-O(g^{2-o(1)})\bigr)I .
 \label{eq:cam-soft-path}
\end{equation}
The second estimate uses normalization by \(\beta_m\). The principal quadratic symbol is \(c_{\mathrm{dir}}I+\widehat{R}_{\rm curv}(p)^*\widehat{R}_{\rm curv}(p)\), where \(\widehat{R}_{\rm curv}\) is the symbol of the linearized renormalized curvature. In this normalization, the quadratic action inherited from earlier scales is a perturbation of operator norm \(O(g^{2-o(1)})\). Both its zeroth and first displacement moments vanish, so its symbol vanishes to second order at \(p=0\). At the fixed stage chosen here, \(c_{\mathrm{dir}}>0\). We write \(C_g=g^{-o(1)}\) for fixed-order derivative bounds and exponentially weighted row-sum bounds, allowing the constant to increase from one occurrence to the next. The full blocked symbol is retained at general frequencies, and the zero-frequency value has the Wilson normalization. Parallel transports along paths are computed using their lifts on the actual periodic lattice. \cite[Lemmas 12.3--12.4]{armstrong}

Set \(t_a=-{\rm i}\sigma_a/2\), where the \(\sigma_a\) are the Pauli matrices. We identify a Lie algebra element with its real coefficient vector in the basis \((t_a)\). We use \(|\cdot|\) for the Euclidean norm of coefficients and use the Frobenius norm for matrices. The tangent norm for unit quaternions is \(|\cdot|_{\rm grp}\) and satisfies \(\lvert\sum_a y_a t_a\rvert_{\rm grp}=|y|/2\). Use three closed center-valued cochains to send a chosen center to the identity. On each time slice let \(P_i\) denote the three spatial period holonomies based at the common site \(x_0\). For sufficiently small \(\varepsilon_0>0\), the holonomy chamber is \(|\log P_i|_{\rm grp}\le \pi/2+\varepsilon_0\). We take the local Coulomb chart on a slightly larger chamber of the same form. In a flat background, the charged adjoint phase \(\phi_i\) has absolute value at most \(\pi+2\varepsilon_0\) on the given chamber and remains a fixed distance below \(2\pi\) on the larger one. Let \(\mathcal R_1\subset\mathbb R^{3\times3}\) be the cone of real matrices of rank at most one. Write \(d_0\) for the forward difference in the spatial directions and \(d_0^*\) for its adjoint with respect to unweighted lattice summation.

\begin{lemma}[Ordinary Coulomb coordinates in a holonomy chamber]\label{lem:cam-soft-chart}
For a small-curvature configuration in the chamber above, the based gauge orbit has a unique small representative
\[
  U_i(x)=\exp(c_i/n+w_i(x)),\qquad
  \sum_x w_i(x)=0,\qquad d_0^*w=0 .
\]
Let \(C_i=\log P_i\) be the principal logarithm and form the matrix \(C\) from the three columns. Then
\begin{equation}
 \|w\|_\infty\le Cn^2t(1+\log n),\qquad
 |c-C|\le Cn^3t(1+\log n).
 \label{eq:cam-soft-chart-bound}
\end{equation}
Parallel transport the color frame at the base point along the time direction, so that the residual color transformation occurs only at the periodic seam. Denote this seam holonomy by \(\hper\) and let \(c_j\) be the full angle matrix at time slice \(j\). On the lift \(c_{j+q}=\operatorname{Ad}_\hper  c_j\), define \(\Delta_\hper  c_j=c_{j+1}-c_j\). The nonzero spatial modes of the time-link logarithm \(z\) and the angle increments satisfy
\begin{equation}
 \|z\|_\infty+\sup_j |\Delta_\hper  c_j|
       \le C n^3t(1+\log n).
 \label{eq:cam-soft-time-links}
\end{equation}
The chart uses a common color frame across time slices.
\end{lemma}

\begin{proof}
A spanning tree with fixed axial ordering puts the links in a based gauge representative. For an interior edge the loop comparing two root paths is tiled by \(O(n)\) plaquettes of the original lattice. A periodic edge requires \(O(n^2)\) additional plaquettes and leaves the based period holonomy \(P_i\). Telescoping the products in the bi-invariant metric gives errors \(Cnt\) and \(Cn^2t\), respectively. The \(n^2\) plaquettes in the periodic plane give \(\|P_iP_jP_i^{-1}P_j^{-1}-1\|\le Cn^2t\). In quaternion notation write \(P=\cos r+\widehat a\sin r\) and \(P'=\cos s+\widehat b\sin s\), where \(r,s\) are the group angles. Then \(PP'-P'P=2\sin r\sin s(\widehat a\times\widehat b)\). Since \(\sin r/r\) is uniformly bounded below in the chamber, the Lie brackets satisfy \(|[C_i,C_j]|\le Cn^2t\). We use
\[
 G(x)=\prod_{i=1}^3\exp(-x_iC_i/n),\qquad 0\le x_i<n,
\]
as a gauge transformation on the actual finite vertex set. For bounded real \(a,b\), differentiating twice and integrating twice gives
\(\|e^{aC_i}e^{bC_j}-e^{bC_j}e^{aC_i}\|\le C|ab|\,\|[C_i,C_j]\|\).
On each edge, move the relevant fractional powers into the same order. On a periodic edge, \(P_i\) cancels \(e^{-C_i}\). Each transformed edge is within \(d_g=Cn^2t\) of \(e^{C_i/n}\).

Let \(D_C\) be the derivative of the gauge action in exponential coordinates at the constant links \(e^{C/n}\). Restrict \(M_C=d_0^*D_C\) to the gauge algebra with zero spatial mean. Take centered representatives for the spatial Fourier indices \(k\in(\mathbb Z/n\mathbb Z)^3\) and set \(\omega_k=\big(4\sum_i\sin^2(\pi k_i/n)\big)^{1/2}\asymp\min(1,|k|/n)\). First expand at a rank-one matrix \(C_*\). Set
\(p_{0,i}=e^{2\pi {\rm i} k_i/n}-1\),
\(p_{\phi,i}=e^{{\rm i}(2\pi k_i+\phi_i)/n}-1\).
In the enlarged chamber, for every nonzero coordinate frequency,
\[
 \operatorname{Re}(\overline p_{0,i}p_{\phi,i})
 =4\cos(\phi_i/(2n))\sin(\pi k_i/n)
                 \sin((2\pi k_i+\phi_i)/(2n))
 \ge c|p_{0,i}|^2 .
\]
The symbol of \(M_{C_*}\) differs by the logarithmic derivative factor \(I+O(n^{-1})\). If \(s_1\ge s_2\ge s_3\) are the singular values of \(C\), then
\begin{equation}
 V_4(C)=\sum_{i<j}s_i^2s_j^2,\qquad
 \operatorname{dist}(C,\mathcal R_1)^2=s_2^2+s_3^2
 \le C\sqrt{V_4(C)}.                                     \label{eq:cam-soft-minors-bound}
\end{equation}
The cross-product bound and this identity allow the choice \(|C-C_*|\le C(n^2t)^{1/2}\). The perturbation of \(D_C\) is at most \(C|C-C_*|/n\). Left multiplication by \(d_0^*\) gives \(O(\omega_k|C-C_*|/n)\). Consequently \(\|M_C(k)^{-1}\|\le C\omega_k^{-2}\).

The necessary inverse bounds in the lattice supremum norm are
\begin{equation}
        \|M_C^{-1}d_0^*\|_{\infty\to\infty}\le Cn,\qquad
        \|D_CM_C^{-1}d_0^*\|_{\infty\to\infty}
                                    \le C(1+\log n).     \label{eq:cam-soft-row}
\end{equation}
Apply the inverse estimate just obtained to the individual frequencies \(|k|\le K\) for fixed large \(K\). For the remaining frequencies, on a continuous shell of radius \(r\ge K/n\), the symbol of \(M_C\) differs from that of the ordinary lattice Laplacian by \(O(r/n)\). After localization to a smooth shell, the derivatives of order \(j\) of the two inverse multipliers are bounded by \(Cr^{-1-j}\) and \(Cr^{-j}\), respectively. Integrating by parts four times in the three-dimensional Fourier coefficients gives \(Cr^2(1+r|x|)^{-4}\) and \(Cr^3(1+r|x|)^{-4}\). The lattice kernels are obtained by summing over periodic images, with absolute sums at most \(C/r\) and \(C\) per shell. Summing over dyadic shells gives \eqref{eq:cam-soft-row}.

If the initial logarithm is \(C/n+w_0\), solve in the mean-zero gauge variable \(\psi\) the equation
\[
 \psi=-M_C^{-1}d_0^*\{w_0+R(w_0,\psi)\},
\]
The second-order Taylor formula in the bounded local group chart gives
\[
 \|R(w_0,\psi)-R(w_0,\psi')\|_\infty
 \le C(d_g+\|\psi\|_\infty+\|\psi'\|_\infty)\|\psi-\psi'\|_\infty .
\]
The contraction factor on the ball \(\|\psi\|_\infty\le Cnd_g\) is \(O(n^2d_g)=o(1)\). The second bound in \eqref{eq:cam-soft-row} bounds the nonconstant error in the corrected links by \(C(1+\log n)(d_g+n^2d_g^2)\). Normalizing the gauge transformation at the base point introduces a common color rotation, which contributes only \(O(d_g)\) on constant links of size \(O(1/n)\). Taking the spatial average gives \eqref{eq:cam-soft-chart-bound}. If two based representatives lie in the small domain, their based holonomies agree and their constant angles differ by at most \(Cn^3t(1+\log n)\). Propagation along the rooted tree of length \(O(n)\) shows that the gauge transformation between them is within \(Cn^3t(1+\log n)\) of the identity. Subtract the mean from its logarithm and apply the same contraction estimate to the difference of the two gauge equations. The operator \(d_0^*D_C\) annihilates the constant gauge directions, and the base-point condition removes the remaining constant freedom. This proves uniqueness.

After transport along the time link at the base point, adjacent time slices have period holonomies that differ by \(O(nt)\), as follows from the \(n\) spacetime plaquettes along the loop. Substituting \eqref{eq:cam-soft-chart-bound} for the two slices along the spatial tree into the actual spacetime plaquettes shows that the normalized time links are within \(Cn^3t(1+\log n)\) of their base-point value. Choose the common color transport \(R_j\) by \(U_0(x,j)=e^{z_{x,j}}R_j\) and \(\sum_x z_{x,j}=0\). Its deviation from the base-point value has the same order. Transporting these color frames successively leaves the time seam holonomy given by the path-ordered product \(\hper=\prod_j R_j\). This gives \eqref{eq:cam-soft-time-links}.
\end{proof}

\subsection{The Haar factors and time-dependent nonzero spatial modes}

In the original Haar integral we change variables jointly along the spatial gauge orbits and the Coulomb slice. Let \(F_{\rm C}=d_0^*\) have mean-zero range and set \(\Delta_0=F_{\rm C}F_{\rm C}^*\). For \(U=\exp u\), normalized Haar measure in exponential coordinates is \(j(u)\,du\). Set \(M_u=F_{\rm C}D_u\iota_0\), where \(\iota_0\psi=\psi-\psi(x_0)\) identifies mean-zero gauge variables with gauge algebra elements vanishing at the base point. The volume factor of this linear map depends only on \(n\) and is absorbed into \(C_n\) below. The coarea formula along the based gauge orbit and uniqueness of the representative in \eqref{eq:cam-soft-chart-bound} give
\begin{equation}
 d\mu_{\rm Haar}^{\rm quotient}(u)
 = C_n\,\frac{|\det M_u|}{\sqrt{\det\Delta_0}}\,
          \prod_{x,i} j(u_i(x))\,d_{\ker F_{\rm C}}u .             \label{eq:cam-soft-coarea}
\end{equation}
In the gauge-variable integral, \(M_u\) is the derivative of the gauge constraint along the based orbit. Applying the coarea formula to the real link coordinates with linear constraint \(F_{\rm C}u=0\) gives \(1/\sqrt{\det\Delta_0}\). Both steps use the same original measure and preserve the Haar integral over the common color. For each time step write \(U_0(x)=e^{z_x}R\) with \(\sum_x z_x=0\). Up to a fixed constant, this change of variables has Haar factor
\begin{equation}
 \left(\prod_xj(z_x)\right)
 \left|\det\left(n^{-3}\sum_x\frac{\operatorname{ad}z_x}
             {1-e^{-\operatorname{ad}z_x}}\right)\right|\,dz\,dR .   \label{eq:cam-soft-time-haar}
\end{equation}
To obtain this determinant, translate the logarithmic derivative using \(\frac{\operatorname{ad}z}{1-e^{-\operatorname{ad}z}}\) and split it into its mean-zero and constant blocks. The linear term averages to zero, so the second factor equals one at \(z=0\) and its logarithm changes by \(O(n^{-3}\sum_x|z_x|^2)\). The factor is independent of \(R\). Consequently, at fixed total seam holonomy \(\hper\), any color link outside the observed region can carry \(\hper\) by a Haar change of variables. The color frames on two disjoint time intervals can thus each be generated using the original Haar parameters.

\begin{lemma}[Hessian lower bound for time-dependent nonzero spatial modes]\label{lem:cam-soft-normal}
Fix the angle path \(c_\bullet\) and seam holonomy \(\hper\) in one chamber. Rescale the spatial Coulomb variables and the time variables at nonzero spatial frequency together as \((w,z)=g_m\xi\). Each nonzero spatial-frequency block \(A_k\) of the action Hessian satisfies
\begin{equation}
 A_k\succeq C_g^{-1}\omega_k^2 I .
 \label{eq:cam-soft-normal}
\end{equation}
The absolute row sums of \(A_k\) obey an exponentially weighted bound in the periodic time distance. Fixed-order derivatives with respect to the constant angles have row bounds of order \(C_g n^{-r}\) at order \(r\), multiplied by the supremum norms of the variation directions.
\end{lemma}

\begin{proof}
First freeze at a flat background \(D\) and set the seam holonomy to the identity. Gauge-transforming the flat connection to the identity on the periodic lift produces the corresponding covariant frequency shift in the curvature matrix. Let \(p_\phi\) be the covariant spatial difference symbol, \(p_\tau=e^{{\rm i}\tau}-1\) the time difference symbol, and \(\mathcal J_D=I+O(n^{-1})\) the derivative from logarithmic coordinates to the left-translated tangent space. For \(k\ne0\), the preceding estimate on \(\operatorname{Re}(p_0^*p_\phi)\) shows that \(\mathcal J_D(\ker p_0^*)\) makes a uniformly positive angle with the line spanned by \(p_\phi\). Writing \(w\) and \(z\) for the variations in the spatial Coulomb and time directions of this block gives
\[
 \|p_\tau\mathcal J_D w-p_\phi z\|^2+
     \|p_\phi\wedge\mathcal J_D w\|^2
 \ge c\bigl[(|p_\tau|^2+\omega_k^2)\|w\|^2+
                      \omega_k^2\|z\|^2\bigr].
\]
Equation \eqref{eq:cam-soft-path} preserves this lower bound. For an angle displacement \(\Delta c\) from the flat background, the background curvature is \(O(|\Delta c|/n^2)\). The linear variation at each slot of an actual path is controlled by \(O(|\Delta c|/n)\|(w,z)\|\) and \(O(|\Delta c|)\) times the covariant curl. Use \(\omega_k\ge c/n\) to absorb the first term into the positive form. The terms with background curvature in the second path variation have row bound \(C_g|\Delta c|n^{-2}\|(w,z)\|^2\). Summing against the original displacement and length moments, the error relative to the static Hessian lower bound is \(C_g|\Delta c|\), which tends to zero in the near-flat chamber of \eqref{eq:cam-soft-chart-bound}.

For a time-dependent \(c_\bullet\), take time localization functions \(\chi_\ell\) with \(\sum_\ell\chi_\ell^2=1\) on supports of width \(R_k=C'_g/\omega_k\) and first differences of size \(O(R_k^{-1})\). Here \(C'_g\) is a sufficiently large constant, or increases more slowly than every negative power of \(g\) as needed. Choose a periodic lift on each support with the seam outside it. The exponential row bounds of the time matrix give an IMS error \(C_g R_k^{-2}\). Choose \(C'_g\) so that this error is only a fixed small fraction of the Hessian lower bound. The derivative bound of a fixed spatial block with respect to a static angle is \(C_g/n\). Freezing at the central time slice of the window therefore gives an error at most \(C_g R_k \sup|\Delta_\hper  c|/n\). By \eqref{eq:cam-soft-time-links}, for the worst frequency \(\omega_k\asymp1/n\), the ratio of this error to the Hessian lower bound is at most
\(C_g n^5t(1+\log n)=o(1)\). The same exponential row bound controls paths extending outside the lifted region. This proves the time-dependent Hessian lower bound. Each angle derivative along a path brings a factor \(1/n\), and the exponential length bounds in \eqref{eq:cam-soft-path} absorb the fixed extra powers of path length.
\end{proof}

\paragraph{Nonlinear integration over the full small-curvature domain.}
Write \(V=n^3q\) and \(\mathcal A=\bigoplus_{k\ne0}A_k\), and extend \(G=\mathcal A^{-1}\) by zero off the normal subspace. Let \(d(j,j_0)\) be the shortest lattice time distance around the period. For each \(k\) conjugate by the weight \(e^{\lambda d(j,j_0)}\). The change in the Hermitian part is controlled by
\(\cosh(\lambda(d(j,j_0)-d(j',j_0)))-1\), so one can take \(\lambda=C_g^{-1}\omega_k\). This gives, block by block,
\begin{equation}
 |G_k(j,j')|\le C_g\omega_k^{-2}e^{-C_g^{-1}\omega_k d(j,j')},
 \quad \|G\|_{\infty\to\infty}\le C_g n^3(1+\log n),
 \quad \sup_eG_{ee}\le C_g .
 \label{eq:cam-soft-fastinv}
\end{equation}
The second bound uses \(\sum_{k\ne0}\omega_k^{-3}\le C n^3(1+\log n)\) and the third uses \(n^{-3}\sum_{k\ne0}\omega_k^{-2}\le C\). We also have \(\|G\|\le C_g n^2\).

On the full small-curvature domain the rescaled variables have magnitude at most \(M_{\mathrm{br}}=C n^3(t/g_m)(1+\log n)\). Use coordinatewise smooth odd truncations that agree with the identity up to \(M_{\mathrm{br}}\) and are bounded by \(2M_{\mathrm{br}}\) outside it. Write the local analytic remainder in all logarithmic coordinates, apply the truncations, and restrict to the normal subspace. Spatial translation invariance makes the first variation vanish in nonzero spatial modes, so the remainder starts at third order. Denote its truncated form by \(R=R_3+R_{\ge4}\), with \(R_3\) still odd. The path-length moments and Taylor's formula give
\begin{equation}
 \begin{split}
 |\partial_eR(\xi)|&\le C_g g\sum_{e'} w_{ee'}|\xi_{e'}|^2,
          \qquad \sup_e\sum_{e'}w_{ee'}\le C,\\
 \|\nabla^2 R\|&\le C_g gM_{\mathrm{br}},\\
 |R_{\ge4}(\xi)|&\le C_g g^2\sum_e|\xi_e|^4 .
 \end{split}                                                     \label{eq:cam-soft-remainder}
\end{equation}
Each path derivative produces a length factor of bounded degree. The weighted arithmetic--geometric mean inequality allocates the two or four local variables to the corresponding sites. The truncations have second derivative \(O(M_{\mathrm{br}}^{-1})\) and preserve the Hessian bound above. Here \(gM_{\mathrm{br}}n^2=g^{7/12-o(1)}\). For the interpolating density \(p_s\propto\exp(-\langle\xi,\mathcal A\xi\rangle/2-sR(\xi))\) with \(0\le s\le1\), the potential Hessian is bounded below by \(\mathcal A/2\). Applying the Brascamp--Lieb inequality to linear tilts of this potential gives sub-Gaussian bounds about the coordinate means, with variance proxy at most \(2G_{ee}\). With \(\mathbf m_s=\mathbb E_{p_s}\xi\), integration by parts gives
\[
 \mathbf m_s=-sG\mathbb E_{p_s}\nabla R,\qquad
 \|\mathbf m_s\|_\infty\le C_g gn^3(1+\log n)(1+\|\mathbf m_s\|_\infty^2).
\]
Since \(\mathbf m_0=0\), continuity in \(s\) and \(gn^3=g^{3/4+o(1)}\) select the small branch of this quadratic inequality, giving \(\|\mathbf m_s\|_\infty=o(1)\) and \(C_g\) bounds on all fixed-order coordinate moments. A further application of Brascamp--Lieb gives
\[
 \operatorname{Var}_{p_s}R\le C_g g^2n^2V,\qquad
 |\mathbb E_{p_0}R|\le C_g g^2V .
\]
In the second estimate the cubic term vanishes exactly by oddness. Integrating the second derivative of \(\log\mathbb E_{p_0}e^{-sR}\) gives a relative logarithmic error \(C_g g^2n^2V\) for the unnormalized integrals. This also bounds the squared Hellinger distance, since the affinity is \(Z(1/2)/\sqrt{Z(1)Z(0)}\). Here \(g^2n^2V\asymp g^{5/6}\).

The same upper bound holds for the original Haar measure remainder in the fast logarithmic variables. Write \(\bar a_i=c_i/n\) for the constant logarithmic background and set \(\Gamma_{\bar a}=M_{\bar a}^{-1}d_0^*\). On the nonzero modes, the absolute row and column sums of \(\Gamma_{\bar a}\) are bounded by \(C_gn\) and their squared-entry sums by \(C_g\). The first estimate uses \eqref{eq:cam-soft-row} and the second uses \(n^{-3}\sum_{k\ne0}\omega_k^{-2}\). For a mean-zero spatial Coulomb variable \(w\), denote the linkwise linear part of \(D_{\bar a+w}-D_{\bar a}\) by \(D_1(w)\). The remainder is \(O(|w|^2)\). It follows that
\[
 \|\Gamma_{\bar a}(D_{\bar a+w}-D_{\bar a})\|\le C_gn\|w\|_\infty,\qquad
 \|\Gamma_{\bar a}(D_{\bar a+w}-D_{\bar a})\|_{\mathrm{HS}}^2\le C_g\sum_e|w_e|^2 .
\]
The trace \(\Tr(\Gamma_{\bar a}D_1(w))\) is a spatially translation-invariant linear functional and hence vanishes on the nonzero spatial modes. Expand \(\log\det\) in its convergent power series and use the squared-entry row and column sums to bound terms of order two and higher. The resulting logarithmic determinant error on the mean-zero gauge variables is at most \(C_g\sum_e|w_e|^2\). When \(F_{\rm C}=d_0^*\), the extra factor for fixing the base point is exactly one, because the gauge tangent of a constant color generator \(\alpha\) is \([\alpha,u]\) and its divergence is \([\alpha,d_0^*u]=0\). The linear Haar correction on the links vanishes by the zero-mean condition. The time factor in \eqref{eq:cam-soft-time-haar} starts at second order. The same strongly convex interpolation controls the exponential moments of the nonlinear Haar weight and gives a logarithmic change of at most \(C_g g^2 V\). Indeed, one may introduce a quadratic tilt of order at most \(C_g g^2 I\) while preserving the Hessian lower bound \(\mathcal A/2\). Uniform sub-Gaussian bounds hold along the interpolation, and the sum of the sitewise second moments is bounded by \(C_gV\). Integrating the first derivative of the logarithm of the normalization constant gives the claim.

Integration over the full domain gives an upper bound on all allowed paths. The lower bound in a center chamber requires that the conditioning set contain a narrower domain. In the logarithmically growing slow-angle window at the center, the magnetic plaquette contributions of the constant angles are at most \(C g^{4/3}M^2/n^2\). The electric contributions are controlled by sub-Gaussian truncation of the actual soft-mode increments. Restrict the normal coordinates to \(|\xi_e|\le \bigl(\exp((\log\log\widehat\beta_m)^2)\bigr)^{1/2}\). This removes probability below every prescribed power of \(g\). The magnetic and electric contributions, together with the normal fluctuations, lie below the plaquette threshold \(t\). At the fixed stage containing scale \(m\) the covariance diagonal is bounded by a fixed constant. The reference paths used in the center lower bound therefore belong to the same original small-curvature set. The full-domain upper bound and the narrow-domain lower bound use the same Gaussian normalization.

\subsection{The determinant at a static background and local freezing in time}

\begin{lemma}[Static zero-point energy relative to a common flat reference]\label{lem:cam-soft-hard}
Use as common normalization the nonzero-spatial-mode Gaussian mass at zero angle and identity seam holonomy, together with the measure factors in \eqref{eq:cam-soft-coarea}--\eqref{eq:cam-soft-time-haar}. Write \(W_g(D)\) for the excess normal-mode energy at a static angle \(D\), per interval of physical duration \(L_0\), after factoring out this scalar. We have \(W_g(0)=0\). If \(D\in\mathcal R_1\) lies in the enlarged holonomy chamber, then
\begin{equation}
 W_g(D)+4\sqrt{c_ec_m}|D|
   \ge c_0|D|-o(\rho),\qquad \rho=g^{5/9}.                        \label{eq:cam-soft-well}
\end{equation}
The coefficients \(c_e,c_m\) are the kinetic and magnetic coefficients of the constant-angle effective Hamiltonian, namely \(L_0H_{\rm eff}=c_e g_m^2(-\Delta_c)+c_m g_m^{-2}V_4(c)\). Both equal \(1/2\) in the Wilson normalization of \eqref{eq:model-scale}. The notation \(W_g\) suppresses its dependence on the stopping scale and the finite lattice dimensions. On the rank-one cone its modulus of continuity is \(C_gd(1+|\log d|)\), where \(d\) is the angular distance. For \(c\) in the chamber, let \(D\) be its nearest point in \(\mathcal R_1\). The frozen energy satisfies
\begin{equation}
 W_g(c)\ge W_g(D)-C_g n^3\operatorname{dist}(c,\mathcal R_1)^2 .
 \label{eq:cam-soft-rank1}
\end{equation}
The normal-mode integral factor \(G_\hper (c_\bullet)\) for a time-dependent path satisfies
\begin{equation}
 -\log\frac{G_\hper (c_\bullet)}{G_1(0)}
   =\frac1n\sum_j W_g(c_j)
       +O\!\left(C_g n^5\sum_j|\Delta_\hper c_j|^2\right)+O_M(g^M)
 \label{eq:cam-soft-logdynamic}
\end{equation}
uniformly for each prescribed \(M>0\). If the static Haar factor is included in \eqref{eq:cam-soft-logdynamic}, the constant-angle Haar correction is counted only once in the subsequent estimates.
\end{lemma}

\begin{proof}
\emph{Linear transverse and longitudinal changes of variables.}
For each nonzero real spatial Fourier block of a flat angle \(D\), the tangent map is
\[
 (w,\psi)\longmapsto \mathcal J_Dw+d_D\psi ,
                  \qquad w\in\ker d_0^* .
\]
Here \(d_D\) is the spatial gauge derivative in the left-translated metric, \(P_D\) is the orthogonal projection onto the complement of \(\operatorname{ran}d_D\), and \(B_D=P_D\mathcal J_D|_{\ker d_0^*}\) is invertible. Splitting the map into its transverse and longitudinal blocks gives Jacobian
\(\lvert\det B_D\rvert\sqrt{\det(d_D^*d_D)}\).
The linear change of variables \(w\mapsto B_Dw\) cancels the first factor. Write \(\mathcal J_Dw-B_Dw=d_D\chi(w)\) and simultaneously shift the temporal variable \(z\) by \(\Delta_t\chi(w)\), with \(\Delta_t\) the forward time difference on the static background. This shift has Jacobian one, and the curvature range of the quadratic action is the range of the covariant curl. When \(Q=I\), Wilson Gauss integration gives \(\det(d_D^*d_D)^{-1/2}\), which cancels the second factor exactly. Each complex-conjugate Fourier pair is counted as one real block and all determinants use that real block. The logarithmic change of the residual Haar factor for the three spatial constant angles is at most \(Cq|D|^2/n^2\), because the decomposition leaves three factors \(j(D_i/n)\) per slice.

If \(Q\ne I\), write \(d_p\) for the curl restricted to the three-dimensional normal complement and \(\bar P(p)\) for the projection onto \(\operatorname{ran}d_p\). Whitening \(d_p\) by its positive Gram matrix on the original positive branch gives
\begin{equation}
 \frac{\det(d_p^*Q(p)d_p)}{\det(d_p^*d_p)}
 =\det_{\wedge^2}\big(I+\bar P(p)(Q(p)-I)\bar P(p)\big),
 \quad
 \bar P(p)=\frac{d_1(p)d_1(p)^*}{\sum_\mu|e^{{\rm i}p_\mu}-1|^2}.
 \label{eq:cam-soft-projectdet}
\end{equation}
The kernel of \(d_1\) is the full gauge direction, and the stated normal subspace is a complement to it. Thus the full \(d_1\) gives the same range projection. Let \(\Psi(p)\) be the logarithm of the right-hand side of \eqref{eq:cam-soft-projectdet}. By \(Q(0)=I\), the Hessian lower bound, and \(\partial^j\bar P=O_j(|p|^{-j})\), we have \(|\partial^j\Psi(p)|\le C_g|p|^{1-j}\) for \(0\le j\le6\). Set \(\Psi(0)=0\) by continuous extension. On a dyadic shell the integral is bounded by \(C_g r^5\), and the integral of sixth derivatives by \(C_g/r\). Split at \(r=|w|^{-1}\) and integrate by parts six times on the larger shells to obtain
\begin{equation}
               |\widehat\Psi(w)|\le C_g(1+|w|)^{-5}.              \label{eq:cam-soft-projfourier}
\end{equation}
Reality of the action coefficients makes \(\Psi\) even. The spatial Fourier sum selects only the indices \(nw\). After averaging over continuous time frequency, the change with angle of the normal energy per step multiplied by \(n\) is
\[
 \left|C n^4\sum_{w\ne0}\widehat\Psi(0,nw)
                   \bigl(\cos(w\cdot\phi)-1\bigr)\right|
 \le C_g n^{-1}
 \begin{cases}d^2\log(2/d),&d\le1,\\1,&d>1,\end{cases}
 \quad d=\operatorname{dist}(\phi,2\pi\mathbb Z^3).
\]
This follows from the shell bound on \(\sum |w|^{-5}\min(2,d^2|w|^2)\). The additional winding contribution for the actual discrete period \(q\) to the logarithm of the total integral is at most
\(C_g n^3 q\sum_{l\ne0,w}(1+q|l|+n|w|)^{-5}\le C_g/q\).
To exclude the spatial mode \(k=0\) from this comparison, take the \(\pm\phi\) average of \(\Psi(p_0,\phi/n)-\Psi(p_0,0)\) in the same time-frequency sum. On \(|p_0|\le2d/n\) use \(C_g|p|\) and the frequency count \(1+O(qd/n)\). On the remaining frequencies use the second derivative bound \(C_g/|p_0|\). Restoring the factor \(n/q\) gives \(C_g[d/q+(d^2/n)(1+\log(n/d))]\). Here and below, \(|\phi|\) and the distance \(d\) to the center are uniformly comparable in the enlarged chamber.

\emph{Positive winding expansion of the Wilson contribution.}
Define
\[
 \eta_n(k,\phi)=4\sum_i \sin^2((2\pi k_i+\phi_i)/(2n)),
 \quad \gamma(x)=2\operatorname{arsinh}(\sqrt x/2).
\]
The temporal electric kernel at \(Q=I\) for a transverse eigenvalue \(x\) is
\((2\pi)^{-1/2}\exp[-(a-b)^2/2-x(a^2+b^2)/4]\).
Matching coefficients with Mehler's kernel shows that the ground-state eigenvalue of one real transverse mode is \(e^{-\gamma(x)/2}\). The two transverse polarizations and the two charged roots give the physical zero-point potential
\[
  2n \sum_k[\gamma(\eta_n(k,\phi))-\gamma(\eta_n(k,0))].
\]
The continuous-time random-walk kernel on \(\mathbb Z\) is
\[
 p_s(j)=\frac{1}{2\pi}\int_{-\pi}^{\pi}
       e^{-2s(1-\cos x)-{\rm i}jx}\,dx .
\]
Here \(s\) is the random-walk time, while \(t\) remains the small-curvature threshold. Integrating \(\gamma'(x)=\int_0^\infty e^{-sx}p_s(0)\,ds\) from zero to \(x\) gives the winding expansion of the finite Fourier sum
\begin{equation}
 \begin{aligned}
 \sum_k[\gamma(\eta_n(k,\phi))-\gamma(\eta_n(k,0))]
 &=\sum_{w\ne0} a_n(w)[1-\cos(w\cdot\phi)],\\
 a_n(w)&=n^3\int_0^\infty p_s(0)\prod_i p_s(nw_i)\,\frac{ds}{s}.
 \end{aligned}
 \label{eq:cam-soft-positive-winding}
\end{equation}
All coefficients are positive. Set \(M=n|w|_\infty\). For \(M^2\le s\le2M^2\), the substitution \(x=s^{-1/2}\xi\) leaves the Gaussian peak as the leading contribution to the Fourier integral, and \(|j|\le \sqrt s\) implies \(p_s(j)\ge c/\sqrt s\). On a sufficiently large fixed window in \(\xi\), the rescaled integral uniformly approximates the positive Gaussian Fourier transform. Positivity covers the remaining finitely many \(j\) on the bounded time interval. For the upper bound, shift the contour by a purely imaginary amount of magnitude \(\lambda\asymp\min(1,|j|/s)\) to obtain
\(p_s(j)\le C(1+s)^{-1/2}e^{-c j^2/(s+|j|)}\). For \(s\le1\), the Bessel series with positive coefficients gives \(p_s(M)\le Cs^M/M!\). On \(1\le s\le M\), the largest absolute coordinate gives \(O(e^{-cM})\), and the rest of the integral is controlled by \(\int_M^\infty s^{-3}e^{-cM^2/s}\,ds\). Hence
\(a_n(w)\asymp 1/(n|w|^4)\).
For small \(d\), restrict the sum in \eqref{eq:cam-soft-positive-winding} to \(|w|\le(2d)^{-1}\). The inequality \(1-\cos(w\cdot\phi)\ge c(w\cdot\phi)^2\) and the permutation symmetry of the lattice ball give a lower bound \(cd/n\). For larger angles take \(w=\pm e_i\). Using \(1-\cos(w\cdot\phi)\le\min(2,d^2|w|^2/2)\) gives the upper bound \(Cd/n\). The difference bound \(\min(2,d'|w|)\) further gives \(C d'(1+|\log d'|)/n\), where \(d'\) is the distance between the two angles.

For \(k\ne0\), \(\gamma(\eta_n(k,\phi))\ge c/n\), so the logarithmic correction from excited-mode occupations at the actual heat time is at most \(Cn^3e^{-cq/n}\) relative to the zero-point contribution. At \(k=0\), \(2n\gamma(\eta_n(0,\phi))=2|\phi|+O(|\phi|^3/n^2)\). In the original \(t_a=-{\rm i}\sigma_a/2\) metric, \(|\phi|=|D|\), giving precisely the leading term \(4\sqrt{c_ec_m}|D|\). Conversion to the unit-quaternion angle multiplies both sides by the same factor. Subtract this spatial zero-mode term from the full sum to give the frozen potential in \eqref{eq:cam-soft-well}, with the four charged constant modes retained. The finite-\(q\) absolute correction \(C_g n/q^2\) and the exponential tail are both \(o(\rho)\). The normal-mode constant in \eqref{eq:cam-soft-coarea} and the calculation in \eqref{eq:cam-soft-projectdet} use the same reference across the chambers. For a nonflat static background, the inverse bound from \eqref{eq:cam-soft-normal} and two angle derivatives give a remainder \(C_g n^{-2}\omega_k^{-4}\) per time and frequency. Summing over \(k\) and multiplying by \(n\) gives \(C_gn^3\). The first normal derivative vanishes. Indeed, at a rank-one matrix, common color rotations about its axis admit no invariant linear functional in the directions normal to the rank-one cone. At \(D=0\) use the full color rotation group. This proves \eqref{eq:cam-soft-rank1}.

\emph{The time-dependent trace and reflection.}
Fix \(k\ne0\) and write the resolvent of \(A_k\) as \(R_z=(z+A_k)^{-1}\). Conjugation by the exponential weight in the actual time distance, as in \eqref{eq:cam-soft-fastinv}, gives weighted resolvent bound \(C_g/(z+C_g^{-1}\omega_k^2)\). Consider \(F_{k,u}=\Tr(P_u\log A_k)\), assigning integer times to spatial variables and half-integer times to temporal links. On the lift centered at observation time \(u\), move the seam holonomy outside the observation window. Choose a radius
\(r_k=C_{M}C_g\omega_k^{-1}\log(1/g)\ll q\). Retain the actual path within \(r_k/2\) and interpolate smoothly outside \(r_k\) to the constant \(c_j\) or \((c_j+c_{j+1})/2\). Since
\(r_k\sup|\Delta c|=o(1)\), the interpolation remains in the local domain of \eqref{eq:cam-soft-normal}. The resolvent identity places the difference between the original and truncated matrices between two exponentially weighted resolvents, giving \(C_g\omega_k^{-2}e^{-C_g^{-1}\omega_k r_k}\). Choose \(C_M\) large enough that the error remains \(O_M(g^M)\) after summation over \(q n^3\) observation points.

The second directional derivative is
\[
 D^2(\log A)[v,w]=\int_0^\infty\bigl(
 R_z A''[v,w]R_z
 -R_z A'[v]R_z A'[w]R_z
 -R_z A'[w]R_z A'[v]R_z\bigr)dz .
\]
Add exponential weights at the two outside ends emanating from \(P_u\) and control \(\partial A\) and \(\partial^2A\) by the local path length moments. With the directional norm \(\|v\|_{k,u}=\sup_s e^{-c C_g^{-1}\omega_k d(u,s)}|v_s|\), the Schur row and column bounds and the \(z\) integral give
\[
 |D^2F_{k,u}[v,w]|\le C_g n^{-2}\omega_k^{-4}
                                   \|v\|_{k,u}\|w\|_{k,u}.
\]
Reflection about an integer time makes the linear coefficient \(L_r(c)\) of the frozen trace satisfy \(L_r=L_{-r}\). Pairing \((j,j+r)\) with the reverse pair rewrites the linear term exactly as
\[
 \tfrac12\sum_{j,r}[L_r(c_j)-L_r(c_{j+r})]\,(c_{j+r}-c_j).
\]
In the same second-derivative estimate take one constant direction and one impulse at \(j+r\) to get
\[
 \operatorname{Lip}(L_r)\le C_gn^{-2}\omega_k^{-4}e^{-cC_g^{-1}\omega_k|r|}.
\]
At half-integer times put \(b_{j+1/2}=(c_j+c_{j+1})/2\). Reflection gives \(L_r=L_{1-r}\). Expand \(c_{j+r}-b_{j+1/2}\) in the linear term as the mean of the increments relative to \(c_j,c_{j+1}\). Reindex each summand by its endpoint and pair it with its reverse. This leaves the difference of two \(L_r\) evaluated at half-time centers separated by at most \(|r|+1\). The same second-derivative bound controls the difference of the static half-time value from the mean of the adjacent integer-time values by \(C_gn^{-2}\omega_k^{-4}|\Delta c_j|^2\).

Across the periodic seam, use the same color transport along the short arc on both sides. The remaining common conjugations preserve these scalars. Using
\[
 \sum_j|c_{j+r}-c_j|^2\le r^2\sum_j|\Delta_\hper  c_j|^2,
\]
the sum of the integer- and half-integer-time linear terms is at most \(C_gn^{-2}\omega_k^{-7}\sum_j|\Delta_\hper  c_j|^2\).
For the quadratic remainder, use Cauchy--Schwarz along short arcs and exponential distance decay and then sum over the base time \(j\) to get at most \(C_gn^{-2}\omega_k^{-6}\sum|\Delta_\hper  c_j|^2\). Apply \(\sum_{k\ne0}n^{-2}\omega_k^{-7}\le Cn^5\) and the normal determinant exponent \(-1/2\) to obtain \eqref{eq:cam-soft-logdynamic}.
\end{proof}

\subsection{Integration outside the center neighborhoods and the eight centers}

First consider crossings between holonomy chambers. At time \(j\), let \(T_i(j)\) denote the average fundamental trace of the \(n^2\) parallel period loops along direction \(i\). The change of an individual loop between adjacent slices is bounded by the sum over the \(n\) actual spacetime plaquettes. Cauchy--Schwarz on each loop gives
\begin{equation}
 n|\Delta T_i(j)|^2\le C\sum_x|F_{0i}(x,j)|^2 .
 \label{eq:cam-soft-crossing}
\end{equation}
If the time oscillation in all three directions is smaller than a fixed \(\varepsilon\), the three trace signs on the initial slice select the center. Each spatial loop is connected to the based loop by \(O(n^2)\) plaquettes, with resulting discrepancy \(O(n^2t)\), so the enlarged chamber of this center contains all time slices. Otherwise sum \eqref{eq:cam-soft-crossing} along an arc of length at most \(q\) in one of the directions to obtain total squared curvature at least \(c\varepsilon^2n/q\).

To use the curvature lower bound for the analytic action, take localization functions whose squares sum to one with supports in four-dimensional cubes of side \(R=g^{-\zeta}\), with \(0<\zeta<1/12\). Parallel transport between curvatures is isometric. The second moments of the original exponentially weighted row sums bound the IMS error by \(C_gR^{-2}\|F\|_2^2\). Discarding paths of length \(>R/10\) costs less than any prescribed positive power of \(g\). Lift each cube, enlarged by the path ranges, to a contractible region. A fixed axial tree gauge puts each link within \(O(R^2t)\) of the identity, and hence changes the parallel transports between curvatures from their identity-background values by at most \(O(R^3t)\). Comparing each localized quadratic curvature form with its flat symbol, \eqref{eq:cam-soft-path} gives a global form lower bound. The higher-order curvature terms and the quadratic terms from earlier scales are small perturbations. Thus, relative to the flat action value,
\[
  S(U)-S(1)\ge C_g^{-1}\beta_m \sum_{\text{plaquettes}}|F|^2-o(1).
\]
The cross-chamber contribution is at most \(\exp(-c C_g^{-1}g^{-4/3})\) times the same flat constant. For a denominator in the same measure, restrict all links to a product of Haar balls about the identity of radius \(g^K\), taking fixed \(K\) large enough that the action increment on this set is \(o(1)\). The logarithmic volume loss of the product is at most \(C_K n^3q\log(1/g)=g^{-1-o(1)}\). It is smaller than the curvature suppression exponent by a positive power of \(g\), so the relative cross-chamber mass is smaller than every prescribed positive power of \(g\). Within a single chamber we use the more precise common denominator below.

\paragraph{Two-sided estimate for the constant-angle action.}
On the actual lift \(c_{j+q}=\operatorname{Ad}_\hper  c_j\) write \(A=\max|c_j|\), \(u_A=A/n\), \(E_i(j)=\Delta_\hper c_i(j)/n\), and \(B_{ik}(j)=[c_i(j),c_k(j)]/n^2\). The leading Baker--Campbell--Hausdorff terms of the electric and magnetic plaquettes are \(E_i\) and \(B_{ik}\). Each higher electric term vanishes when adjacent angles agree, and the twice-integrated Duhamel formula for the group commutator allows one to factor out \([c_i,c_k]\) from each remaining mixed word in the magnetic plaquette. Thus the first relative corrections in both cases are \(O(u_A)\) and the remaining relative error is \(O(u_A^2)\). For a color transport word of length \(r\), the product of orthogonal factors gives a second-order remainder \(Cr^2u_A^2\) without any exponential factor in \(r\). Reflection acts by \(c\mapsto-c\). The first relative electric and magnetic square corrections, the leading mixed electric--magnetic term and the leading three-electric term are odd and vanish exactly. The remaining mixed electric--magnetic term is bounded by \(C_gu_A\sqrt{Q_\hper Q_B}\). Each remaining three-slot term has at least one magnetic slot from which one extracts \(Cu_A^2\), and the four-slot terms give at least two amplitude factors. Spatial reflection diagonalizes the magnetic directions at zero spatial frequency, and time reflection makes their symbols even. Zero-frequency normalization and the second row moments bound the symbol difference between the magnetic square and the direct quartic term by \(C_g\omega(p)^2\). Using
\(\lvert\Delta_\hper [c_i,c_k]\rvert\le CA|\Delta_\hper c|\), this error is at most \(C_g u_A^2 Q_\hper \). Combining the quadratic terms inherited from earlier scales with \eqref{eq:cam-soft-logdynamic} gives
\begin{equation}
 \bigl|S_{\rm slow}-Q_\hper -Q_B-\tfrac1n\sum_jW_g(c_j)\bigr|
 \le C_g[(u_A^2+g^{4/3})(Q_\hper +Q_B)+u_A\sqrt{Q_\hper Q_B}]+R_{\mathrm{Haar}},
 \label{eq:cam-soft-slowform}
\end{equation}
where \(Q_\hper \ge C_g^{-1}g_m^{-2}n\sum|\Delta_\hper c|^2\),
and \(Q_B=c_m g_m^{-2}n^{-1}\sum V_4(c_j)\). The separate constant-angle Haar term is bounded by \(C q A^2/n^2\). At the chosen \(n\), \(C_g n^5\sum|\Delta_\hper c|^2\) is bounded by \(C_g g^{5/3}Q_\hper \), so the bound with \(g^{4/3}\) is conservative. Set \(\omega(p)=2|\sin(p/2)|\). For the eigenphase \(\alpha\in[-\pi,\pi]\) of \(\operatorname{Ad}_\hper \), take temporal frequencies \(p=(2\pi j+\alpha)/q\). After factoring out \(g_m^{-2}n/2\), the symbol of the electric form is \(\omega(p)^2 a_g(p)\), with \(a_g(0)=1\), \(a_g\ge C_g^{-1}\), and
\(\lvert a_g(p)-1\rvert\le C_g\omega(p)^2\). These properties follow from \eqref{eq:cam-soft-path} on the electric curvature subspace and time reflection.

\begin{lemma}[Outer-holonomy tail with a common center normalization]\label{lem:cam-soft-tail}
Let \(Z_{\mathrm{cent}}\) be the unnormalized mass of a fixed soft ball at one center after factoring out the same normal-mode integral as above. The original small-curvature mass at \(\max_j|c_j|\ge\rho\) divided by \(Z_{\mathrm{cent}}\) is \(O(e^{-c g^{-1/9+o(1)}})\).
\end{lemma}

\begin{proof}
The upper bound on the full normal domain and the lower bound on the narrower center domain use the same normal-mode integral in \eqref{eq:cam-soft-coarea}. The denominator on a fixed soft ball at the center is bounded below by a fixed positive multiple of the product of this factor, the free soft Gaussian integral over nonzero temporal frequencies, and the volume of the soft constant-angle window. The independent Gaussian verification of this lower bound is given below. Divide all comparisons here by these common factors.

On the shell \(A/2\le\max|c_j|\le A\) with \(A\ge\rho\) take \(\epsilon=g^{1/100}\). Young's inequality in \eqref{eq:cam-soft-slowform} reduces the electric form to a factor \(1-C_g[\epsilon^{-1}(A/n)^2+g^{4/3}]\) and the quartic potential to a factor \(1-O(C_g\epsilon)\). The logarithmic cost in the electric Gaussian normalization is \(C_gq[\epsilon^{-1}(A/n)^2+g^{4/3}]=o((q/n)A)\). By \eqref{eq:cam-soft-minors-bound}, the second-order loss in \eqref{eq:cam-soft-rank1} is absorbed by \(\epsilon g^{-2}V_4+C_g\epsilon^{-1}g^2n^6\). The second term is \(o(\rho)\) per physical time unit. The separate constant-angle Haar term obeys \(C_g qA^2/n^2=o((q/n)A)\).

To normalize the pure electric Gaussian reference, add \(g^J/(g_m^2 n)\) times \(\sum_j|c_j|^2\) to the action, with \(J\) a sufficiently large fixed exponent. Removing this mass on the shell contributes only a factor \(e^{o(1)}\). Truncate the exponentially decaying temporal matrices at distance \(r=C_g\log^2(1/g)\). Relative to the smallest eigenvalue, the omitted part is smaller than any prescribed positive power of \(g\). The truncated matrices for the actual seam and for the identity seam differ only within \(O(r)\) of the seam. The finite-rank estimate gives relative log determinant \(O_J(r\log(1/g))=g^{-o(1)}\). The same argument applies to the positive matrices obtained by compressing to four coordinates and adding an oscillator mass. The auxiliary mass gives a uniform positive lower bound, allowing a seam holonomy that mixes the four coordinates with the others. The ratio of the massive integral of the electric zero mode to the soft constant-angle window volume is \(g^{-O_J(1)}\), whose logarithm is bounded by \(g^{-o(1)}\).

If two slices on the lift differ by more than \(\epsilon A\), Cauchy--Schwarz along the intervening time arc in the actual electric form gives a variance bound \(C_gg^2|j-j'|/n\) for a directional increment. The union bound gives at most \(Cq^2 e^{-C_g^{-1}g^{-4/3}\epsilon^2 A^2}\). This beats \(\exp(C_g(q/n)A)\) because \(A\ge g^{5/9}\), and the relative bound is \(e^{-c(q/n)A}\). All remaining paths stay inside a ball of radius \(O(\epsilon A)\). For a ball away from the rank-one cone, \eqref{eq:cam-soft-minors-bound} gives \(V_4\ge c\epsilon^2 A^4\). Since \(g^{-2}\epsilon^2A^3\to\infty\), its action controls the negative part of the frozen hard-mode potential. Near the rank-one cone, fix an anchor \(D\) and rotate the row and column axes to make \(D_{11}=|D|\). The sum of the corresponding squared \(2\times2\) minors gives
\[
 V_4(c_j)\ge(1-O(\epsilon))|D|^2 |Z-Z_*(Y)|^2,\qquad
 (Z_*)_{i\ell}=c_{i1}c_{1\ell}/c_{11}\quad (i,\ell\ge2).
\]
Write \(Z\) for the four entries of the lower-right \(2\times2\) block and \(Y\) for the other five coordinates. Condition on the entire \(Y_\bullet\) path and write \(\Gamma\) for the conditional covariance of \(Z_\bullet\). Let \(\varkappa\) denote the quadratic precision increment in these four coordinates induced by the magnetic potential, with weight \(e^{-\varkappa\|Z-Z_*(Y)\|^2/2}\), where the squared path norm is summed over slices. Completing the square shows that the time-dependent displacement from the conditional mean contributes a factor at most one. Hence
\(\mathbb E[e^{-\varkappa\|Z-Z_*(Y)\|^2/2}\mid Y]\le\det(I+\varkappa\Gamma)^{-1/2}\).
The finite-rank seam estimate above permits evaluation of this determinant with the identity seam. The mass parameter in each charged direction is
\(\omega_*=2\sqrt{c_ec_m}|D|/n\,[1+o(1)]\).
The periodic product identity is
\[
 \prod_{p\ne0}\big(1+\omega_*^2/\omega(p)^2\big)
   =\left[\frac{2\sinh(q\,\operatorname{arsinh}(\omega_*/2))}
                {q\omega_*}\right]^2 .
\]
By \(|a_g(p)-1|\le C_g\omega(p)^2\), restoring the actual temporal symbol changes the logarithm by at most \(C_g q\omega_*^2=o((q/n)A)\). Each charged direction contributes \(-q\omega_*/2+o((q/n)A)\) to the logarithm. Their sum exactly accounts for the charged term in \eqref{eq:cam-soft-well}. The cost of the modulus of continuity of the rank-one frozen potential is
\(C_g\epsilon A(1+|\log(\epsilon A)|)\), and the loss in \eqref{eq:cam-soft-rank1} is absorbed in the magnetic error above. This leaves the strictly negative exponent \(-c(q/n)A\). On each shell the fixed-dimensional rank-one anchor set is covered by \(O(\epsilon^{-9})\) balls, using a single anchor for the entire path. Summing over dyadic shells and the finite collection of chambers gives the upper bound \(\exp(-c(q/n)\rho)\), with \((q/n)\rho\asymp g^{-1/9}\). Closed center-valued cochains permute the eight centers and give them equal weights. Cross-chamber paths were controlled in \eqref{eq:cam-soft-crossing}.
\end{proof}

\subsection{Growing center balls and the soft-mode measure with a seam holonomy}

We keep \(\max_j|c_j|\le \rho\) and use the exact change of variables \(c=g_m^{2/3}X\) from \eqref{eq:model-scale}. The effective step size at the stopping scale is \(d_s=g_m^{2/3}/n\) and the full effective time is \(u_m=q d_s=Tg_m^{2/3}/L_0\). At the lattice times chosen in the main text, \(u_m\) lies in a fixed compact subset of \((0,\infty)\). After integrating out all nonzero spatial modes and dividing out the common flat normalization, the reference electric action for a central slow path is
\((2d_s)^{-1}\langle X,J_\hper X\rangle\), with \(J_\hper (p)=\omega(p)^2a_g(p)\) and \(a_g\) as defined after \eqref{eq:cam-soft-slowform}. The reference magnetic action is \(c_m d_s\sum_j V_4(X_j)\). For the orders of errors we use \(d_s\asymp g^{3/4}\).

\paragraph{Uniform soft-path tails.}
First bound the tail of the nonlocal reference for \(1\ll M\le g^{-1/9}\). On a shell \(A/2\le\max |X_j|\le A\), add to the electric form
a mass \(\mu_A^2=q^{-2}(1+A)^{-2}\). Removing this auxiliary mass on the shell contributes only a factor \(e^{O(1)}\). For a seam phase \(\alpha\), the quasiperiodic lattice Laplacian with any mass \(\mu^2>0\) satisfies
\begin{equation}
 \prod_{j=0}^{q-1}\!
   \left[4\sin^2\!\frac{2\pi j+\alpha}{2q}+\mu^2\right]
   =2\{\cosh(q\eta)-\cos\alpha\},
 \qquad \eta=2\operatorname{arsinh}(\mu/2).                     \label{eq:cam-soft-twist}
\end{equation}
This follows by factoring \(4\sin^2(p/2)+\mu^2=e^\eta(1-e^{-\eta+{\rm i}p})(1-e^{-\eta-{\rm i}p})\) and taking the product over quasiperiodic frequencies. In the term with \(a_g(p)\), first extract \(\prod_p a_g(p)\). By \(a_g(p)/a_g(0)-1=O(C_g\omega(p)^2)\), the difference in the residual logarithms is at most \(C_g q \mu_A^2=o(1)\). The seam dependence of \(\prod_p a_g(p)\) is carried by the Fourier coefficients of \(\log a_g\) with indices that are nonzero multiples of \(q\), and hence is smaller than any prescribed positive power of \(g\). Since \(\cosh(q\eta)-\cos\alpha\) takes its minimum at \(\alpha=0\), the total Gaussian mass relative to the standard nonzero-temporal-frequency integral times the volume of the fixed constant-angle window is at most \(C(1+A)^9\).

The supremum of increments of this Gaussian process on a fixed effective time interval has a uniform tail in \(g\). The near-zero mode in each phase branch is controlled by
\[
 |e^{{\rm i}\alpha s/u_m}-e^{{\rm i}\alpha s'/u_m}|^2/
                  (\alpha^2+\mu_A^2 q^2)\le C|s-s'|^2
\]
The variance from the other low frequencies is \(O(|s-s'|)\), by
\(\sum_{j\ne0} j^{-2}\min(1,|j||s-s'|)^2\le C|s-s'|\).
Here we use the uniform lower bound on \(a_g\) for \(|p|\le g^{\eta_0}\) with \(\eta_0>0\) sufficiently small. The high-frequency pointwise variance is at most \(C_gq^{-1}g^{-2\eta_0}\), and we take the union bound over \(q\) points. For the low frequencies use an increment threshold \(c A 2^{-k/4}\) at dyadic level \(k\). Summing over the at most \(C2^k\) events at each level gives \(\exp(-cA^2)\). Thus the part of the shell with displacement from an anchor exceeding a fixed small fraction of \(A\) has mass bounded by \(C(1+A)^9e^{-cA^2}\).

The remaining paths stay in a ball of radius \(O(\epsilon_0 A)\). Away from the rank-one cone, \eqref{eq:cam-soft-minors-bound} gives \(d_s \sum V_4\ge c A^4\). Near the cone use the four minor coordinates of Lemma \ref{lem:cam-soft-tail}, where a tail of order \(e^{-cA}\) suffices here. Condition on the five full tangential paths and on the time points at the seam. Since
\(J_\hper \le C L_\hper +C_g L_\hper ^2\), where \(L_\hper \) is the quasiperiodic nearest-neighbor Laplace form, the compression to the four coordinates away from the seam is bounded above by
\(C L_D+C_g L_D^2\) plus a fixed-rank endpoint term. In fact, the difference between squaring before and after the compression is
\(P L_\hper (I-P)L_\hper  P\) and involves only the fixed number of nearest-neighbor endpoints. The \(j\)th eigenvalue of each \(L_D\) is at most \(Cj^2/q^2\). By finite-rank interlacing, at least \(cA-O(1)\) eigenvalues of the conditional precision matrix are at most \(CA^2/q^2\), because \(C_g A^2/q^2=o(1)\) throughout the range up to \(A=g^{-1/9}\). The four-dimensional oscillator adds precision \(c d_s^2 A^2\). Completing the square gives a factor uniformly less than one for each of these eigenvalues, and displacing the center contributes only a nonpositive exponent. This yields the shell bound \(C(1+A)^9e^{-cA}\). The actual central slow action after integrating out the fast modes is bounded by \eqref{eq:cam-soft-slowform}. The change in the electric coefficient times \(q\) is \(o(1)\), and the static term is at most
\(C_g(q/n)n^3\rho^2=O(g^{7/36-o(1)})\). Both the constant-angle Haar correction and the nonlinear normal integration cost are \(o(1)\). The shell estimate therefore holds throughout the center ball. Summing the shells gives tail \(Ce^{-cM}\).

For continuous Brownian loops and the nearest-neighbor reference the same argument works on all shells. Below the lattice frequency limit, the Dirichlet eigenvalues in effective time are at most \(Cj^2\), yielding \(e^{-cA}\) again. For the discrete reference at \(A\ge q\), all \(q-1\) eigenvalues in each conditional coordinate are at most \(Cq^2\), so the shell mass is at most
\[
 C(1+A)^9[1+c(A/q)^2]^{-2(q-1)} .
\]
Summing over dyadic \(A\ge q\) gives \(O(q^9e^{-c q})\). The continuous reference has no frequency truncation. The denominator on the actual soft ball has a fixed positive lower bound in the same normalization. Start within a fixed small ball, which also contains the \(\hper\) image. The free heat kernel at fixed positive effective time has a positive lower bound there. A Brownian bridge stays within a slightly larger fixed ball with positive probability, and the potential weight has a positive lower bound in that ball. For the nearest-neighbor reference take the sampled event of the same bridge. The local whitening and density square-root comparison below give the bound on the same fixed ball for the nonlocal slow reference restricted to the center. This lower bound is established before using exterior suppression.

\paragraph{Density square-root error at the center.}
Take \(M=C_D(1+|\log g|)\) and compare on paths with \(\max|X|\le M\). The static central correction has vanishing first derivative at \(c=0\). The same second-order bound as in \eqref{eq:cam-soft-rank1} gives at most
\[
     C_g(q/n)n^3 g^{4/3}M^2=O(g^{5/12-o(1)}M^2).
\]
In \eqref{eq:cam-soft-slowform} use Young's inequality with parameter \(g^{3/8}\). In the logarithmic window \(Q_B\le CM^4\), the difference of the slow action from the common real reference satisfies
\(\lvert R\rvert\le r_g+\delta_e Q_\hper \), where
\[
 r_g=O(g^{3/8-o(1)}(1+M)^4),\quad
 \delta_e\le C_g(g^{-3/8}g^{4/3}M^2/n^2+g^{4/3}),\quad
 q\delta_e=O(g^{3/8-o(1)}(1+M)^2).
\]
Add the auxiliary soft mass \(q^{-2}(1+M)^{-2}\) to the free electric reference. For the same actual seam holonomy, the unnormalized mass is at most \(C(1+M)^9\) as above, and removing the mass on this window contributes a bounded factor. Under the normalized Gaussian law, \(Q_\hper \) has the form
\(\frac12\sum_{k=1}^{9q}\lambda_kZ_k^2\), with \(0\le\lambda_k\le1\). The standard Gaussian moment generating function gives
\(\mathbb E[(1+Q_\hper )^2e^{C\delta_e Q_\hper }]\le Cq^2\).
The pointwise inequality
\(\lvert e^{-R/2}-1\rvert^2\le R^2e^{|R|}/4\) thus gives the following squared error for the unnormalized density square roots:
\[
 C(1+M)^9(r_g^2+q^2\delta_e^2).
\]
The Hellinger cost of the normal-mode integration is \(O(g^{5/12-o(1)})\) and is added separately. For each actual seam holonomy normalize using the same lower bound on the fixed soft ball, and include the already bounded mass outside the central neighborhood. Taking \(C_D\) large enough makes the tail beyond the logarithmic ball smaller than every prescribed power of \(g\). The density square-root error in this step is therefore \(O(g^{3/8-o(1)})\).

\subsection{Local whitening in time, short bridges, and the full observables}

\begin{lemma}[Comparison of the full Coulomb chart with the two-endpoint heat kernel]\label{lem:cam-soft-compare}
At the two original cuts retain every link used by the prescribed observables, and let \(C_{N,m}\) be the square root of the joint density of these coarse-grained observables. In an equivariant common Hilbert space it differs from
\[
 \frac1{\sqrt8}I_{\mathrm{cent}}\otimes
 \frac{e^{-u_mH_4/4}}{\sqrt{\operatorname{Tr}_{\mathrm{sing}}e^{-u_mH_4/2}}}
\]
in Hilbert--Schmidt norm by \(O(g^{3/8-\delta-o(1)})\). At each endpoint the original observable uses \(O(g^{-2\delta})\) coarse time slices, with \(0<\delta\le 10^{-3}\). Here \(I_{\mathrm{cent}}\) is the identity on \(\mathbb C^8\), which indexes the eight center copies, and \(\operatorname{Tr}_{\mathrm{sing}}\) is the trace on the color-singlet subspace. The comparison includes the nonzero spatial modes and the color-orbit coordinates of the full coarse-grained observables. It has the same error order with the prescribed time \(u\) of the main text in place of \(u_m\).
\end{lemma}

\begin{proof}
\emph{Local whitening in time.}
Let \(B_\hper \) be the quasiperiodic convolution operator for seam holonomy \(\hper\), with symbol \(B(p)=(a_g(0)/a_g(p))^{1/2}\). The change of variables \(X=B_\hper Y\) gives the exact electric symbol \(a_g(0)\omega(p)^2\). At the chosen renormalization stage, the time matrix has exponential row bounds and \(a_g\) has a positive lower bound. A complex Taylor estimate on a narrower strip of width \(C_g^{-1}\) shows that both \(B\) and \(B^{-1}\) are analytic there and bounded uniformly by a constant of order \(C_g\). When the convolution is truncated at distance \(g^{-\delta'}\), the absolute Fourier tail is \(e^{-g^{-\delta'+o(1)}}\), with \(0<\delta'<\delta\). At the actual \(q\) frequencies, the dependence of \(\log\det B_\hper \) on \(\hper\) is carried by the Fourier coefficients of \(\log B\) at indices that are nonzero multiples of \(q\). It is therefore smaller than any prescribed positive power of \(g\). Normalize both sides with the identity seam as reference, then integrate over \(\hper\) with Haar measure. After the convolution is truncated, each endpoint observable uses only the variables in its own short time window.

Under the massive references for \(X=B_\hper Y\) and \(Y\), the bound \(B(p)-1=O(C_g\omega(p)^2)\) removes the temporal near-zero frequency singularity in the covariance. At each site,
\(\mathbb E|(B_\hper -I)Y_j|^2\le C_g d_s\), with the same estimate for the inverse change of variables. Write
\(P(Y)=c_m d_s\sum_j V_4(Y_j)\). When both paths stay within the ball of radius \(2M\), the square roots of the magnetic weights satisfy
\[
 |e^{-P(B_\hper Y)/2}-e^{-P(Y)/2}|^2
 \le C M^6 d_s^2 q \sum_j|(B_\hper -I)Y_j|^2 .
\]
The common unnormalized Gaussian mass is bounded by \(C(1+M)^9\). Including it gives square-root error
\(O(d_s^{1/2}C_g(1+M)^{15/2})\). When just one path is in the amplitude window, add the mass on that side and bound the maximum of the Gaussian difference of the paths. When both amplitudes are large, bound each side by its own shell tail. The resulting error is \(O(g^{3/8-o(1)})\). The covariance and positive-weight comparison for truncating the convolution also proceeds frequency by frequency in the massive reference. The relative symbol errors are smaller than every prescribed power of \(g\), and at most \(Cq\) eigenvalues occur. Undoing the convolution retains the same Lipschitz bound on the magnetic weight.

\emph{Short Brownian bridges.}
Compare the nearest-neighbor discrete quartic weight with the sampled continuous model on the same grid \(j d_s\). Both use the same free one-step heat kernel. Given adjacent samples \(x,y\) of norm at most \(2M\), put \(L_r=x+(r/d_s)(y-x)\) and let \(B_r\) be a Brownian bridge with zero endpoints. Set \(I=c_m\int_0^{d_s}V_4(L_r+B_r)\,dr\), \(I_0=c_m\int_0^{d_s}V_4(L_r)\,dr\) and \(I_{\rm tr}=c_m d_s(V_4(x)+V_4(y))/2\). By \(\|\nabla^2V_4(z)\|\le C|z|^2\), the trapezoidal error is at most \(C d_s M^2|y-x|^2\). The odd random terms in the Taylor expansion have zero expectation, and fixed-order Gaussian bridge moments give
\[
 |\mathbb E(I-I_0)|\le C d_s^2(1+M)^2,\qquad
 \mathbb E|I-I_0|^2\le C d_s^3(1+M)^6 .
\]
Since \(I,I_0,I_{\rm tr}\ge0\), the first two derivatives of \(e^{-x}\) are bounded by one on the relevant interval. The preceding three estimates therefore bound \(\lvert \mathbb E e^{-I}-e^{-I_{\rm tr}}\rvert\). Gaussian tails for the actual massive sampling allow restriction of all adjacent increments to \(C_D\sqrt{d_s\log(1/g)}\). After removing the mass in the \(M\) window, the discarded mass is smaller than every prescribed power of \(g\). Both single-step weights are bounded away from zero, so their logarithms may be compared directly. Summing over the steps gives \(C d_s(1+|\log g|)^3\) and controls the density square roots on the common truncation window. Using the common positive normalization preserves this order. Conditional on the sample points, the bridge segments are independent, so the comparison holds for the full sample law and for any sub-observable. At the final step the actual common seam holonomy \(\hper\) simply rotates a fixed endpoint, preserving all the moment bounds and the Haar integration.

\emph{Freezing the normal-mode integrals at the endpoints.}
After comparing the nonlinear normal integral on the full small-curvature domain with its Gaussian approximation, fix the central slow path. Compare the resulting normal marginals on the two endpoint neighborhoods with those at the zero field. Retain all \(w,z\) variables at each nonzero spatial frequency \(k\). The row difference bound \(C_g\rho/n\) and the Hessian lower bound \(C_g^{-1}\omega_k^2\) give the relative precision and covariance error
\(\varepsilon_k\le C_g n\rho/|k|^2\).
Compression to the marginals preserves this relative inequality. For two \(d\)-dimensional Gaussians with relative spectrum in \([1-\varepsilon,1+\varepsilon]\), the affinity is the product of \(\sqrt{2\sqrt\lambda/(1+\lambda)}\) over the eigenvalues. The squared error of the density square roots is thus at most \(C d\varepsilon^2\). Retaining \(O(g^{-2\delta})\) time slices for each \(k\) and using \(\sum_{k\ne0}|k|^{-4}<\infty\) gives squared error \(O(g^{17/18-2\delta-o(1)})\). By \eqref{eq:cam-soft-fastinv}, once the time seam is placed outside the two observed neighborhoods, both the marginal error caused by freezing the seam and the error from cross-endpoint correlations are \(g^{-O(1)}e^{-C_g^{-1}q/n}\). This comparison can first whiten by the common covariance for each \(k\) and then compress. The resulting normal variables are two local Gaussian noises independent of the slow endpoints. Equations \eqref{eq:cam-soft-coarea}--\eqref{eq:cam-soft-time-haar} supply the same spatial and temporal gauge-orbit Haar variables and preserve the relative colors within each observable.

\emph{A common density square root for model endpoints and full observables.}
The positive measure comparison above uses all variables in both endpoint neighborhoods. Given a short path segment of the nine-dimensional model, the truncated whitening convolution, local normal Gaussian noise, gauge-orbit Haar parameters, and inverse coordinates of \eqref{eq:cam-soft-chart-bound} generate variables of the original observable type. The original fundamental winding traces determine the center exactly. Four heat-propagation segments at the model endpoints make the two observables conditionally independent given the central model endpoints. The diffusion time in these four segments is \(O((g^{-2\delta}+g^{-\delta'})d_s)\). For \(u_m\) in a fixed compact subset of \((0,\infty)\), log-convexity of the model heat trace bounds the Hellinger error from these segments by \(O(g^{3/8-\delta})\). The nine-dimensional transverse oscillator lower bound gives finite heat traces at those times and at shorter positive times.

Let \(K_{u_m}(y,y')\ge0\) denote the integral kernel of the normalized model heat operator relative to ground-state probability \(\pi\). More explicitly, divide the integral kernel of \(e^{-u_mH_4/4}\) on the color-singlet subspace by \(\sqrt{\operatorname{Tr}_{\mathrm{sing}} e^{-u_mH_4/2}}\) and by the ground-state values at both endpoints. This gives \(K_{u_m}\). Let \(J_y(dr)\) be the conditional law of the observable given the model endpoint \(y\). Use the common reference measure \(\upsilon(dy,dr)=\pi(dy)J_y(dr)\) at each end. The square root of the joint observable density is
\[
 B_{\mathrm{rep}}(r,r')=
 \sqrt{\mathbb E_{\upsilon\otimes\upsilon}[K_{u_m}(y,y')^2\mid r,r']}.
\]
Define \(M_1=\mathbb E[K_{u_m}\mid r,r']\) and \(M_2=\mathbb E[K_{u_m}^2\mid r,r']\). The squared error conditioned on the two observables is \(2M_2-2\sqrt{M_2}M_1\le2(M_2-M_1^2)\). Hence on the same product space,
\begin{equation}
 \|K_{u_m}-B_{\mathrm{rep}}\|_2^2
 \le 2\mathbb E\operatorname{Var}(K_{u_m}\mid r,r').
 \label{eq:cam-soft-condvar}
\end{equation}
Within each endpoint observable, local inverse filtering recovers the central color orbit of \(Y\). Since the original \(X\) is obtained from \(Y\) by the same short filter, the absolute row sum of the composition of the two truncated filters minus the identity is at most \(e^{-g^{-\delta'+o(1)}}\). Even after multiplying by the polynomial number of sites, this is smaller than every prescribed positive power of \(g\). The \(L^2(\pi)\) endpoint tails are controlled by the heat-kernel diagonal at fixed positive time. On a logarithmically growing ball one can therefore use local inverse filtering to construct the predictor in \eqref{eq:cam-soft-condvar}. Both the heat kernel and its first derivatives have uniform bound \(e^{C(1+M)^4}\) on the ball of radius \(2M\). The upper bound uses the free kernel in the Feynman--Kac formula and a first-derivative Duhamel estimate with integrable time singularity \(s^{-1/2}\). For the lower bound, use the positive probability that the free bridge between two endpoints in the ball stays in a slightly larger ball, together with the magnetic potential bound \(C(1+M)^4\). Integrating one endpoint over a fixed ball of positive ground-state probability gives the corresponding bounds on the ground state and its reciprocal. The locally truncated \(K_{u_m}\) therefore has Lipschitz constant at most \(e^{C(1+M)^4}\). The squared norm outside the ball is bounded by the model heat tail \(Ce^{-cM}\) established above. Thus the recovery error is smaller than every prescribed positive power of \(g\). Since the conditional variance is taken with respect to the full actual observables, the estimate controls the full joint density square root.

The action and whitening each contribute an error \(O(g^{3/8-o(1)})\), the normal noise contributes \(O(g^{17/36-\delta-o(1)})\), and conditional decoupling at the endpoints contributes \(O(g^{3/8-\delta})\). Closed center-valued cochains permute the eight centers equivariantly, and cross-center mass is smaller than every prescribed positive power of \(g\). Haar averaging the seam holonomy projects to the color-singlet subspace. Normalizing the half-period heat operator by the full-period heat trace and including the eight equally weighted centers gives the operator in the lemma. By \eqref{eq:paper-coupling} and \eqref{eq:paper-longtime}, \(|u_m-u|=O((\log b)/b+g^{2/3})\). Since the heat trace is finite at every positive time, the Hilbert--Schmidt derivative of the normalized model heat operator is uniformly bounded on each compact subset of \((0,\infty)\). For instance, \(\|H_4 e^{-u_mH_4/4}\|_{\rm HS}\) is controlled by a shorter positive-time heat trace. Replacing \(u_m\) by \(u\) adds an error \(O((\log b)/b+g^{2/3})\) of smaller order than the error in the lemma.
\end{proof}
\section{Further Proofs for the Finite-Volume Vacuum Representation}
\label{app:cam-vacuum}

\subsection{Finite Blocking Histories and Their Conditional Laws}
\label{app:cam-fixed-law}

Fix the square torus, the terminal inverse coupling, and the finest blocking level \(f\) used to generate any of the observables. Let \(U_i\) be the configuration at level \(i\), and let \(\mathcal B_{N,i}(U_{i+1},dU_i)\) be the original probabilistic blocking kernel from level \(i+1\) to level \(i\). For fixed \(f\), write \(r_{N,f}\) for the level-\(f\) density after removing a configuration-independent scalar. The published fixed-level results give convergence along the full cutoff sequence for every path coefficient and every activity coordinate with its covering weight \cite[Proposition 5.7 and Lemma 13.24]{armstrong}. On this fixed finite lattice, group the activities by support and admissibility. The tails of the logarithmic regular factor are controlled by the anchored absolute sums. For the forced factor, use \(\sum_{k\ge0}A^k/k!=e^A\), where \(A\) is the absolute sum of the weights over all supports. In subtracting two exact-cover products, keep one difference factor at a time. The remaining absolute sum is bounded by \(e^{A+A'}\). Thus \(r_{N,f}\) is Cauchy in Haar \(L^1\). The small-curvature domains and admissibility events on the fixed lattice are the same for both densities. Requiring all free off-tree links to lie in a sufficiently small ball about the identity prevents every large-field test from being triggered and gives a uniform upper bound on the action after removal of the scalar. This ball has fixed positive Haar measure, so \(\inf_N\int r_{N,f}>0\). The normalization estimate
\[
 \left\|\frac{r}{\int r}-\frac{r'}{\int r'}\right\|_1
     \leq \frac{2\|r-r'\|_1}{\min(\int r,\int r')},
\]
transfers the Cauchy property to the actual probability laws of the observables. This step depends only on the densities after removal of the field-independent scalar.

At each fixed level, the blocking conditional law first selects the local branch with its original probability and then adds independent group noise to each coarse link. Let \(\mathsf b>0\) be the fixed blocking penalty coefficient and \(w(W)=1-\tfrac12\operatorname{Re}\Tr W\) the link action of the noise variable \(W\). Its density relative to normalized Haar measure,
\(q_\beta=e^{-\mathsf b\beta w}/\int e^{-\mathsf b\beta w}\), satisfies
\[
 \partial_\beta q_\beta=-\mathsf b(w-\mathbb E_{q_\beta}w)q_\beta,
 \qquad \|\partial_\beta q_\beta\|_{L^1}\leq 2\mathsf b .
\]
Replacing one coarse link at a time and averaging over the original branches gives the conditional \(L^1\) bound, uniform in the given fine field,
\[
 \|\mathcal B_{N,i}-\mathcal B_{N',i}\|_{L^1}
 \le 2\mathsf b |E_i|\,|\beta_{N,P,i+1}(v)-\beta_{N',P,i+1}(v)|,
\]
where \(E_i\) is the set of positively oriented links at level \(i\). Let \(\mu_{N,f}\) be the normalized probability measure at level \(f\). The full fixed finite segment
\(d\mu_{N,f}(U_f)\prod_{i<f}\mathcal B_{N,i}(U_{i+1},dU_i)\)
converges in total variation. Hence the expectation of any bounded Borel function of these variables converges. If the recorded variable takes values in a Haar-null image, retain the links and auxiliary variables that generate it and take the measurable pushforward of their jointly convergent laws.

\subsection{Short Arms of the Vacuum Paths and Observable Kernels}
\label{app:cam-bridge}

At level \(N\), use the normalized ground state \(\Omega_N>0\) of the full Gauss-law space for the ground-state transform and write \(S_N(t)\) on the ground-state probability space \(L^2(\pi_N)\). Its kernel \(k_t(x,y)\) is symmetric with respect to \(\pi_N\), preserves constants, and satisfies the semigroup identity. On the short time circle \(T_0=L_0\), the joint density of the opposite slices relative to \(\pi_N\otimes\pi_N\) is \(k_t(x,y)^2/Z_N^0(T_0)\), where \(t=T_0/2\). Under the reference law, conditional on the two slices, draw the four short arms independently using the ground-state transform. Each starts at its respective slice and has duration \(h\), rounded to a lattice time.

Given the slices \(x,y\), let \(A_1,A_2\) be the two evaluations of \(k_{t-2h}\) at the pairs of remote arm endpoints on the two sides. They are conditionally independent and satisfy \(\mathbb E(A_i\mid x,y)=C=k_t(x,y)\). With respect to the common reference probability, the actual extended path and the approximate path with conditional independence given the fine slices have densities \(A_1A_2/Z_N^0(T_0)\) and \(C^2/Z_N^0(T_0)\), respectively. Write \(a_i=\mathbb E(\sqrt{A_i}\mid x,y)\leq\sqrt C\). Conditional independence gives the squared Hellinger distance multiplied by \(Z_N^0(T_0)\) as \(2\int(C^2-Ca_1a_2)\,d\pi_N(x)d\pi_N(y)\). Moreover,
\[
 C^2-Ca_1a_2\leq C^{3/2}\sum_i(\sqrt C-a_i)
       =\tfrac12 C\sum_i\mathbb E\big[(\sqrt{A_i}-\sqrt C)^2\mid x,y\big].
\]
Use \(C/(\sqrt{A_i}+\sqrt C)^2\leq1\) to bound the last expression by the corresponding sum of the squares \((A_i-C)^2\) with coefficient \(\tfrac12\). Under the independent reference law, the two remote endpoints in each evaluation again have marginal \(\pi_N\). Thus
\[
 \int\mathbb E(A_i-C)^2\,d(\pi_N\otimes\pi_N)
       = Z_N^0(T_0-4h)-Z_N^0(T_0).
\]
This gives the forward error used in the main text. The same calculation applies to a long time circle with an integer number of lattice time steps.

To retain an observable band of fixed positive width centered within one half-period, first fix the entire original path and the blocking variables in that band. Let \(w\) be the half-width of the band and choose \(t-2w>s_*>0\) with \(Z_N^0(s_*)\leq C_*\). The three intervals without observables have lengths \(s_1=s_2=t/2-w\) and \(s_3=t\). After adding the short arms at the two slices, the kernel between the remote endpoints of interval \(j\) is \(A_j=k_{s_j-d_j^{\mathrm{arm}}}\), with \(0\leq d_j^{\mathrm{arm}}\leq 2h\) the total duration of the arms at the ends of that interval. Given its original endpoints, \(\mathbb E A_j=C_j=k_{s_j}\). Replace the \(A_j\) by \(C_j\) in succession. In each difference term, integrate out all paths other than the given interval, including the observable band, using the actual semigroup. This gives the kernel \(k_{T_0-s_j}\) on its two endpoints. These integrals include the values of the observables, so the estimate applies to their entire joint law. By Cauchy--Schwarz, the \(L^1\) norm of the difference term before normalization is at most
\[
 \big[Z_N^0(2s_j-2d_j^{\mathrm{arm}})-Z_N^0(2s_j)\big]^{1/2}
                 Z_N^0(2(T_0-s_j))^{1/2}.
\]
Each \(2s_j\) exceeds \(s_*\) by a fixed positive amount. Monotonicity and log-convexity of the heat trace give \(O(h^{1/2})\) for this term, and the common normalization satisfies \(Z_N^0(T_0)\geq 1\). The forward Hellinger error with the observable retained is therefore \(o(1)\).

For the reverse comparison take the auxiliary observable level \(f\) finer than the fixed level of the observable. In the actual blocking history, conditional on the entire level-\(f\) field, a coarser observable uses only new independent blocking variables. It does not change the conditional law of the finer field. Retain the observable as shared information in the scale-by-scale endpoint localization. After the observable bins have been integrated out, the copies have the original full-field marginals. The local exceptional events can therefore be averaged against those marginals without a condition-number loss depending on the bins.

Let \(p=p(x,u_1,y,u_2,\gamma)\) be the joint density including the observable \(\gamma\) in the middle time band, and let \(p_{xy\gamma}\) and \(p_{u_1u_2\gamma}\) be the corresponding marginals. The two approximate laws are
\[
    p_{xy\gamma}\,r_x(du_1)r_y(du_2),
   \qquad
    p_{u_1u_2\gamma}\,s_{u_1}(dx)s_{u_2}(dy),
\]
where \(r\) generates the one-end coarse-grained observable using ground-state paths, and \(s\) is the fine-slice conditional law given the coarse-grained observable at that end under the finite-period measure. On replacing a candidate kernel by the original one-end conditional kernel, contraction of total variation under marginalization bounds the additional error at each end by one copy of the original comparison error. Compare the square-root kernels of the two approximate laws in the same endpoint space. The endpoint isometries extend them by the density square roots \(\sqrt{r_x}\), \(\sqrt{s_{u_1}}\), and their analogues at the other end. To extract the observable kernel from the joint law, compare the actual density \(p\) with the reference density \(\widetilde p\) using
\[
 \left\|\int|\mathcal R(p)_\gamma-\mathcal R(\widetilde p)_\gamma|\,d\gamma\right\|_{L^2}
       \leq3\|\sqrt p-\sqrt{\widetilde p}\|_{L^2},\qquad
 \mathcal R(p)_\gamma=\frac{p(\cdot,\gamma)}{\sqrt{\int p(\cdot,\gamma')\,d\gamma'}} .
\]
For the proof, fix a pair of endpoints, let \(A,B\) denote the two corresponding square roots as functions of \(\gamma\), and suppose \(\|A\|\geq\|B\|\). Use \(\||A|^2-|B|^2\|_1\leq(\|A\|+\|B\|)\|A-B\|\), then change the second denominator to \(\|B\|\). The total error is at most \(3\|A-B\|\). The same bound holds directly when a marginal vanishes.

Let \(\mathcal T_N(\gamma)(x,y)\) denote the density in \(\gamma\) of the heat kernel with the observable inserted, with reference measure \(\pi_N\) for the slice variables. Cutting the path integral at the two slices gives
\[
 p_{xy\gamma}=\frac{\mathcal T_N(\gamma)(x,y)k_t(y,x)}{Z_N^0(T_0)}.
\]
It follows that \(\mathcal R(p)\) is the observable kernel with the same factor \(1/\sqrt{Z_N^0(T_0)}\) as the normalized half-period heat operator. For fixed \(f\), the total variation limit of the previous subsection and the square-root kernel bound give the joint operator limit on the fixed comparison space. Then \(f\to\infty\) removes the endpoint approximation error. The constructions preserve the center cochain actions, so the comparison is obtained under the positive unprojected law and then restricted to zero electric flux.

\subsection{Contractible Loop Traces Separate the Finite-Level Center Quotient}
\label{app:cam-loop}

On a three-dimensional slice or an open time band, choose at each vertex \(x\) two basic plaquette holonomies \(P_x,Q_x\) in distinct spatial orientations. Identify \(\mathrm{SU}(2)\) with the unit quaternions. Use Gram--Schmidt on \(\Im P_x,\Im Q_x\) to obtain the covariant orthonormal frame \(e_1(x),e_2(x),e_3(x)\), taking the third axis to be the cross product of the first two. Under product Haar measure, almost surely the two imaginary parts are not collinear and each basic plaquette trace is nonzero. Indeed, each exceptional set is defined by a nonidentically zero real-analytic equation. Each of the two chosen plaquettes has a link unique to it, so its holonomy can be varied independently of the other. Only finitely many conditions are required on the entire lattice.

At the edge \(e=(x,y)\), read the matrix
\[
        \bigl(e_i(x)\cdot\operatorname{Ad}_{U_e}e_j(y)\bigr)_{ij}.
\]
For two plaquette holonomies \(A,B\), the required inner product is determined by
\[
 \Im A\cdot\operatorname{Ad}_{U_e}\Im B
   =(\Re A)(\Re B)-\Re(AU_eBU_e^{-1}).
\]
For a cross-product term, replace \(A\) by \(P_xQ_x-Q_xP_x\) and expand by linearity. The normalization denominators are also given by the plaquette traces at the same endpoint and the traces of their products. All the resulting closed paths lie in the local contractible region of the two adjacent vertices.

Suppose two link configurations have the same local traces. Use vertex gauge transformations to align their frames. Since their adjoint matrices on each link agree, the links differ only by center signs \(z_e\in\{\pm1\}\). The basic plaquette traces are nonzero and agree, giving \(\prod_{e\in\partial p}z_e=1\). Thus \(z\) is a closed \(1\)-cochain. Exact cochains are vertex center gauges, and the remaining classes are the three center actions of the spatial torus. Consequently, these finitely many contractible traces separate the quotient by gauge and center transformations on a full-measure set. Take the image of the compact finite-lattice configuration space under these continuous observables. The standard Borel structure on the quotient and injectivity on a full-measure set give a Borel inverse. A bounded Borel observable can then be approximated in \(L^2\) by bounded continuous functions for the finite reference measure on this compact trace space.

The last blocking step transfers this approximation to the actual observables. At fixed observable level and fixed \(b,v\), the density of every final Haar noise variable is uniformly bounded, also conditional on the parent field and the branch. For any fixed finite collection of observed links, multiply the bound a finite number of times. The resulting bound is preserved under mixing over the conditioning. The same argument applies at the observable level on stationary vacuum paths. For two observables \(F,G\), conditional Jensen with respect to the actual initial slice gives
\[
 \|(\mathcal T_N(F)-\mathcal T_N(G))\Omega_N\|^2
     \leq \mathbb E_{\mathrm{vac}}|F-G|^2 .
\]
One may therefore replace Borel representatives of the constant-mode angles by bounded continuous combinations of contractible Wilson loop traces while retaining the strict margin in the first-cluster observable matrix.

\subsection{Reflection Norms of Separated Local Insertions}
\label{app:cam-local-mark}

Let \(V_d\) be the space of adjoint-invariant polynomials in the Taylor coefficients of the connection in radial gauge, of engineering dimension at most \(d\). The coefficient of the homogeneous part of degree \(j+1\) of the connection has weight \(j+2\). For a real label \(\mathbf c\in V_d\), the cutoff point insertion of the published construction at the central lattice site \(y\) is
\[
 \mathcal O_N(\mathbf c;y;U)
   =a_N^{-d}P_d^{\mathrm{loc}}\bigl[t_N^{-\mathsf N_d}\mathsf Z_{N,P,d}^{-1}\mathbf c\bigr](y;U).
\]
The local section \(P_d^{\mathrm{loc}}\) realizes polynomial labels as gauge-invariant functions on a fixed lattice stencil. The operator \(\mathsf N_d\) multiplies a monomial by its number of factors, and \(\mathsf Z_{N,P,d}\) carries cutoff-level labels to the reference level at the fixed renormalization scale. The normalizations are those of the published construction \cite[Definitions 2.7--2.8]{armstrong}.

On the short square torus, choose a reflected pair of points separated by a positive physical distance. At a fixed chessboard blocking level \(k_{\mathrm{cb}}\), take the two points in distinct cells and a positive distance from their cell boundaries. Write \(\theta_{\mathrm E}\) for Euclidean time reflection composed with complex conjugation. Below, we abbreviate the insertion at the fixed position and label by \(\mathcal O_N\). The original odd blocking ratios and even initial period ensure that the lattice size in each direction is compatible with the reflection cells. As the fixed starting level \(f\) becomes finer, the physical dependence radius of the original blocking tends to zero. We may thus choose a bounded coarse polynomial at that level and keep all its finer-scale dependencies within the same cell. The replication estimate in the published bounded-approximation proof \cite[Lemmas 14.14 and 14.29, Eq.~(14.54)]{armstrong} uses the same coefficient bound on a single admissible finite period. Here it suffices to check the finer-scale dependence regions and reflection compatibility of the two cells on that fixed period. Use the terminal parameter for which the original label comparison is established, and keep it fixed throughout the cutoff limit. First fix the error \(\varepsilon>0\), then fix the intermediate levels needed for \(f,k_{\mathrm{cb}}\), and finally remove the cutoff along the original sequence. Chessboard Cauchy--Schwarz and the original cutting and gluing identity give
\[
 \frac{\|\mathcal T_N(\mathcal O_N)-\mathcal T_N(F_{\varepsilon,N})\|_{\HS}^2}{Z_N^0(T_0)}
    =
 \mathbb E_{T_0}\bigl[\theta_{\mathrm E}(\mathcal O_N-\mathcal P_{N,f}F_{\varepsilon,N})
                         (\mathcal O_N-\mathcal P_{N,f}F_{\varepsilon,N})\bigr]
       \leq\varepsilon^2 .
\]
Here \(\mathcal P_{N,f}\) is conditional averaging by the positive blocking kernel, and the auxiliary variables in the two separated dependence regions are conditionally independent. This estimate concerns reflected insertions at separated points and requires no coincident-point variance.

To establish convergence without preselecting a coefficient subsequence for the bounded coarse approximation, take two candidate limits \(F,F'\) in the fixed finite-dimensional coefficient space and set \(G=F-F'\). The actual mixed inner product \(\mathbb E_{T_0}[\theta_{\mathrm E}(\mathcal P_{N,f}G)\mathcal O_N]\) converges along the full cutoff sequence. In the published stopping argument for fixed-level labels \cite[Lemmas 14.31--14.32]{armstrong}, choose a reference level no coarser than \(f\). Given that level, average all later coarsening variables according to their own finite positive conditional law. The additional test still has supremum bound \(\|G\|_\infty\). The original proof controls the coefficient absolute sum for the stopping error at the reference level before taking the Haar integral. Multiplication by this test costs at most \(\|G\|_\infty\). The fixed-level limit then follows from the \(L^1\) comparison for finite conditional kernels in Appendix~\ref{app:cam-fixed-law} and convergence of the fixed-level labeled coordinates. The mixed moment is thus obtained using the original point insertion throughout.

On each of the two subsequences, replace the actual labeled insertion by its bounded approximation, and then subtract. Let \(C_*=\sup_N Z_N^0(T_0)\). The resulting estimate is
\[
 \lim_N\|\mathcal T_N(G)\|_{\HS}^2
       \leq 2\varepsilon\sqrt{C_*}\,\lim_N\|\mathcal T_N(G)\|_{\HS} .
\]
The common operator comparison for bounded fixed-level observables gives both norm limits. Thus the candidate observables are at reflection-norm distance at most \(2\varepsilon\sqrt{C_*}\). Fix one candidate and let \(\varepsilon\to0\) to define the buffered labeled operator on the same vacuum space. Composition with a positive-time contraction preserves the approximation. For finite words, replace the factors one at a time. The low-spectrum matrices are then transferred by continuous spectral projections and the same endpoint isometries. This construction gives the local-field cyclic subspace and its spectral lower bound used in the main text.
\phantomsection
\addcontentsline{toc}{section}{References}
\begingroup\small
\begin{thebibliography}{99}
\raggedright
\setlength{\itemsep}{.25\baselineskip}

\bibitem{wilson}
K.~G. Wilson.
Confinement of quarks.
\emph{Phys. Rev. D} \textbf{10} (1974), 2445--2459.
\href{https://doi.org/10.1103/PhysRevD.10.2445}{DOI}.

\bibitem{luscher-transfer}
M.~L\"uscher.
Construction of a selfadjoint, strictly positive transfer matrix for Euclidean lattice gauge theories.
\emph{Commun. Math. Phys.} \textbf{54} (1977), 283--292.
\href{https://doi.org/10.1007/BF01614090}{DOI}.

\bibitem{osterwalder-seiler}
K.~Osterwalder and E.~Seiler.
Gauge field theories on a lattice.
\emph{Ann. Phys.} \textbf{110} (1978), 440--471.
\href{https://doi.org/10.1016/0003-4916(78)90039-8}{DOI}.

\bibitem{luscher-torus}
M.~L\"uscher.
Some analytic results concerning the mass spectrum of Yang--Mills gauge theories on a torus.
\emph{Nucl. Phys. B} \textbf{219} (1983), 233--261.
\href{https://doi.org/10.1016/0550-3213(83)90436-4}{DOI}.

\bibitem{koller-vanbaal}
J.~Koller and P.~van Baal.
A rigorous nonperturbative result for the glueball mass and electric flux energy in a finite volume.
\emph{Nucl. Phys. B} \textbf{273} (1986), 387--412.
\href{https://www.sciencedirect.com/science/article/pii/055032138690252X}{Journal}.
\newblock
SU(2) spectroscopy in intermediate volumes.
\emph{Phys. Rev. Lett.} \textbf{58} (1987), 2511--2514.
\href{https://doi.org/10.1103/PhysRevLett.58.2511}{DOI}.

\bibitem{vanbaal-koller-torus}
P.~van Baal and J.~Koller.
QCD on a torus, and electric flux energies from tunneling.
\emph{Ann. Phys.} \textbf{174} (1987), 299--371.
\href{https://doi.org/10.1016/0003-4916(87)90032-7}{DOI}.

\bibitem{vanbaal-review}
P.~van Baal.
QCD in a finite volume (2000).
\href{https://arxiv.org/abs/hep-ph/0008206}{arXiv:hep-ph/0008206}.

\bibitem{formal}
OpenAI.
Lean construction of four-dimensional pure Yang--Mills theory, 28 October 2026.
\href{https://github.com/yang-mills-group/OpenAI-YM-20261028/tree/3456e5e}{Formal source, revision 3456e5e}.

\bibitem{armstrong}
The Yang--Mills informalization project.
\emph{The construction of four-dimensional quantum Yang--Mills theory}.
Third draft, 13 November 2026.
Project maintained by S.~Armstrong in collaboration with the humanise-yang-mills project.
\href{https://mgubi.org/yang-mills-informalization/manuscript/main.html}{Manuscript}.

\bibitem{jaffe-witten}
A.~Jaffe and E.~Witten.
Quantum Yang--Mills Theory.
\href{https://www.claymath.org/wp-content/uploads/2022/06/yangmills.pdf}{Clay Mathematics Institute, official problem description}.

\bibitem{thermodynamics}
G.~Boyd, J.~Engels, F.~Karsch, E.~Laermann, C.~Legeland, M.~L\"utgemeier, and B.~Petersson.
Thermodynamics of \(\mathrm{SU}(3)\) lattice gauge theory.
\emph{Nucl. Phys. B} \textbf{469} (1996), 419--444.
\href{https://arxiv.org/abs/hep-lat/9602007}{arXiv:hep-lat/9602007}.

\bibitem{a-theorem}
Z.~Komargodski and A.~Schwimmer.
On renormalization group flows in four dimensions.
\emph{JHEP} \textbf{12} (2011), 099.
\href{https://arxiv.org/abs/1107.3987}{arXiv:1107.3987}.

\end{thebibliography}
\endgroup
\end{document}
